% !TEX program = xelatex
% 2026 年中国研究生数学建模竞赛 F 题
% 算力约束下提升大语言模型能力的资源配置建模
% 编译：xelatex -> bibtex(可选) -> xelatex -> xelatex
\documentclass[bwprint]{gmcmthesis}

\hypersetup{hidelinks}
% 模板的 \thesubsection 里含 \thinskip（\hskip 0.05em），随标题写入 PDF 书签串时
% hyperref 会对每个 \kern 报一次 “Token not allowed in a PDF string”（全文 226 条）。
% 书签串里把 \thinskip 退化为空格即可，正文排版不受影响。
\pdfstringdefDisableCommands{\let\thinskip\space}

% ---- 页面设置：左右边距 2.5 cm（上下与页眉间距沿用模板默认）----
% 模板默认左右各 22.5 mm；此处按中文论文常用的 2.5 cm 设置，正文宽 160 mm。
%\geometry{top=30.0mm,bottom=25.0mm,left=25.0mm,right=25.0mm,headsep=8mm}

% ---- 论文额外需要的宏包（模板已含 amsmath/amssymb/graphicx/booktabs/
%      longtable/tabularx/caption/geometry/needspace 之外的补充）----
\usepackage{mathtools}
\usepackage{mathrsfs}
\usepackage{needspace}
\usepackage{float} % 使用 [H] 时需要
\usepackage[noend]{algpseudocode}
\usepackage{algorithm}
\floatname{algorithm}{算法}
\renewcommand{\algorithmicrequire}{\textbf{输入：}}
\renewcommand{\algorithmicensure}{\textbf{输出：}}

% ---- 定理类环境 ----
% 模板已定义 definition/theorem/lemma/corollary/assumption/proof，但为全文连续编号；
% 本文按节编号并共享同一计数器（与图 1 技术路线图中标注的编号一致）。
% 做法：先用独立名字新建所需环境，再重绑定模板已占用的环境名。
\makeatletter
\let\proof\relax \let\endproof\relax
\makeatother
\usepackage{amsthm}
\renewcommand{\proofname}{证明}
\theoremstyle{definition}
\newtheorem{thmdef}{定义}[section]
\newtheorem{thmassumption}[thmdef]{假设}
\theoremstyle{plain}
\newtheorem{thmlemma}[thmdef]{引理}
\newtheorem{thmtheorem}[thmdef]{定理}
\newtheorem{thmproposition}[thmdef]{命题}
\newtheorem{thmcorollary}[thmdef]{推论}
\theoremstyle{remark}
\newtheorem*{thmremark}{说明}
\renewenvironment{definition}{\begin{thmdef}}{\end{thmdef}}
\renewenvironment{assumption}{\begin{thmassumption}}{\end{thmassumption}}
\renewenvironment{lemma}{\begin{thmlemma}}{\end{thmlemma}}
\renewenvironment{theorem}{\begin{thmtheorem}}{\end{thmtheorem}}
\newenvironment{proposition}{\begin{thmproposition}}{\end{thmproposition}}
\renewenvironment{corollary}{\begin{thmcorollary}}{\end{thmcorollary}}
\newenvironment{remark}{\begin{thmremark}}{\end{thmremark}}

% ---- 正文使用的记号 ----
\newcommand{\pvec}{\mathbf p}
\newcommand{\one}{\mathbf 1}
\newcommand{\Prob}{\mathbb P}
\newcommand{\E}{\mathbb E}
\newcommand{\Var}{\operatorname{Var}}
\newcommand{\Cov}{\operatorname{Cov}}
\newcommand{\Loss}{\mathcal L}
\newcommand{\DeltaJ}{\Delta\widehat J}
\newcommand{\dd}{\mathop{}\!\mathrm d}
\DeclareMathOperator*{\argmin}{arg\,min}
\DeclareMathOperator*{\argmax}{arg\,max}
\newcolumntype{P}[1]{>{\raggedright\arraybackslash}p{#1}}
\setlength{\emergencystretch}{3em}

% ===================== 封面信息（提交前填写）=====================
\title{算力约束下提升大语言模型能力的资源配置建模}
\baominghao{\hspace{6em} 26104590182}
\schoolname{\hspace{5.8em} xxxx}
\membera{\hspace{6em} xxx}
\memberb{\hspace{6em} xxx}
\memberc{\hspace{6em} xxx}
% ==========================段排序=================================

\newcounter{obsitem}
\newcommand{\obsitem}[1]{%
    \refstepcounter{obsitem}%
    \par\noindent\textbf{(\arabic{obsitem})#1}
}
% ==========================subequation=================================
%\usepackage{amsmath}


\begin{document}

\maketitle

\begin{abstract}
本文研究算力约束下大语言模型的资源配置问题，围绕数据质量、领域配比、参数规模、训练数据量与算力预算之间的关系，依次讨论\textbf{质量与配比的可识别性}、\textbf{质量修正下的资源替代}、\textbf{预算约束下的资源配置}以及\textbf{能力前沿建模}。四个问题通过质量标尺、有效数据量和训练算力等中间变量衔接，形成“\textbf{识别边界--有限资源替代--预算配置--条件能力前沿}”的分析路线。

针对数据质量与领域配比问题，本文首先构造统一的数据质量代理，并将领域配比视为概率单纯形上的变量。对于部分领域缺少质量观测的情形，不对未知质量直接作点估计，而是给出\textbf{完整质量的部分识别区间}. 同时证明，在域质量固定时，\textbf{配比效应与质量交互效应不能由现有配比试验分别识别}. 据此，本文将质量评价与配比响应分开建模，并在平方根坐标下建立领域配比响应模型. 对于有限候选集，进一步给出\textbf{一致预测误差向配比选择后悔的传递界}，从而将预测误差与最终决策误差联系起来.

针对数据质量与训练资源之间的关系，本文引入\textbf{质量修正的有效数据量}，并在来源内建立广义标度模型. 在给定模型结构和质量标尺的条件下，利用 Loss 对有效数据量的单调性，推导质量、数据量与领域配比之间保持 Loss 不变的\textbf{有限等价条件},  进一步给出有限质量变化对应的\textbf{等效参数规模及其存在条件}. 为检验共同质量响应假设，本文还给出\textbf{中心化均值矩阵的秩一必要条件}. 上述结果主要用于描述同一模型设定下不同资源状态之间的可比较关系，并不意味着不同数据来源之间具有相同的质量机制。

针对算力预算约束下的资源配置问题，本文在固定质量下将原问题化为一维尺度优化，并证明相应\textbf{最优资源对的唯一性}。进一步将质量选择表示为关于参数规模的一族\textbf{仿射成本函数的下包络}，由此得到在既定假设下的\textbf{预算比较静态结果}：预算增加时，全局最优参数规模严格增加，最优质量不下降. 在此基础上，本文给出\textbf{有限预算达到最高质量水平的充要条件}，并分析质量饱和后的\textbf{高预算渐近结构}. 对于参数估计误差，进一步通过局部稳定性分析和有限情景后悔准则考察最优配置的稳健性.

针对模型能力演进问题，本文首先统一评测口径，并区分模型规模变化、时间残差与训练 Loss 的不同作用. 在此基础上，采用\textbf{条件分位数}描述能力前沿，并通过\textbf{有限顶点约束}保证声明域内多个分位前沿保持有序，同时将联合拟合写成\textbf{凸优化问题}. 进一步利用目标函数在共享斜率方向上的曲率，将数值求解的目标间隙转化为相应的\textbf{参数误差界}. 第三问得到的训练算力路径仅作为第四问的条件输入, 对于超出历史支持范围的长期结果，本文将其解释为\textbf{给定模型和算力情景下的条件外推}，而非已验证的预测.

综上，本文的主要理论成果包括：\textbf{部分信息下的数据质量识别界与配比决策后悔界}，\textbf{质量修正下的有限资源等价条件}，\textbf{基于仿射成本下包络的预算单调性与质量饱和判据}，以及\textbf{具有有序约束和全局可解性的条件分位数能力前沿}. 上述结果均在相应的数据口径、模型假设和声明域内成立，并通过数值计算对其适用性和敏感性进行检验，为算力约束下大语言模型资源配置提供了一套相互衔接的分析方法.

\keywords{大语言模型、数据质量、广义标度律、有效数据量、仿射下包络、比较静态、决策后悔、条件分位数前沿.}
\end{abstract}

\pagestyle{plain}
\setcounter{tocdepth}{2}
\maketoc
\clearpage




\section{问题重述与建模思路}\label{sec:overview}

本题研究算力受限时参数规模 $N$、训练数据量 $D$、数据质量 $Q$ 和领域配比 $\pvec$ 对语言模型能力的共同作用. 本文围绕\textbf{质量识别与配比决策、有效数据量折算、全局资源配置、有序能力前沿}四个递进环节展开. 先确定数据能够识别的质量与配比信息，再建立其与损失的替代关系，进而证明最优配置随预算变化的规律，最后构造具有形状约束和误差控制的能力前沿. 前三问的响应量是验证交叉熵损失 $\mathcal L$，第四问的响应量是基准测试得分 $S$, 二者通过可比资源情景及同口径检验建立联系. 

\subsection{问题重述}
\noindent\textbf{问题一:质量评价与配比选择.}
将 22 项质量指标同向化，建立文档及域级评分，分析指标冲突并在扩展数据上检验. 在 17 域概率单纯形上刻画配比与 Loss 的关系. 数学上需要解决质量域映射不完整时的识别问题，以及响应面预测误差如何影响最终配比选择. \\
\noindent\textbf{问题二：质量、数据与参数的替代关系.}
建立包含 $N,D,Q,\pvec$ 的广义标度律，衔接经典 $N$--$D$ 模型，给出边际效用、弹性及领域替代或互补的解释，并计算 $Q$ 提高 $0.1$ 的等效参数量. 核心是把质量和配比统一为有效数据状态，区分局部替代率与有限变化的精确等价关系. \\
\noindent\textbf{问题三：预算最优配置及结构性转移.}
在训练、质量处理与注意力成本的总和不超过 $C$ 时选择资源配置，比较三类质量成本、至少三档预算及 C7 给定的上下文长度，并说明配比是否能够联立优化. 需要确定全局最优解的预算变化方向、质量达到模型上界的条件，以及高预算下参数量和数据量的增长规律. \\
\noindent\textbf{问题四：历史增长分解与能力前沿.}
统一综合能力、开放模型筛选、预训练与后训练类型及时间口径，量化规模与非规模因素的观察性份额，检验 Loss--Benchmark 映射，并给出算力增长放缓时 12, 24四个月的条件前沿. 数学上需要同时处理分位函数的有序性、全局求解和区间校准的适用条件. 

\subsection{建模框架与跨问衔接}

质量与领域配比首先通过训练分布作用于验证损失. 对领域 $a$，以 $P_a^{\boldsymbol\xi_a}$ 表示语料状态 $\boldsymbol\xi_a$ 下的分布.令 $Z=(H_{\mathrm{ctx}},Y)$ 表示由上下文 $H_{\mathrm{ctx}}$ 与下一 Token $Y$ 构成的训练样本, 并定义单样本交叉熵损失$\ell_{\boldsymbol\theta}(Z):=-\ln p_{\boldsymbol\theta}(Y\mid H_{\mathrm{ctx}})$. 在模型规模与训练规程固定时, 有
\begin{subequations}\label{eq:v40-data-loss-mechanism}
\begin{align}
R_{\mathrm{tr}}(\boldsymbol\theta;\pvec,\boldsymbol\xi)
&:=\sum_{a=1}^{17}p_a\,\E_{Z\sim P_a^{\boldsymbol\xi_a}}\ell_{\boldsymbol\theta}(Z),\\
\boldsymbol\theta^*(\pvec,\boldsymbol\xi)
&\in\argmin_{\boldsymbol\theta}R_{\mathrm{tr}}(\boldsymbol\theta;\pvec,\boldsymbol\xi),\\
\mathcal L_{\mathrm{val}}(\pvec,\boldsymbol\xi)
&:=\E_{Z\sim P_{\mathrm{val}}}\ell_{\boldsymbol\theta^*(\pvec,\boldsymbol\xi)}(Z).
\end{align}
\end{subequations}
式~\eqref{eq:v40-data-loss-mechanism} 给出引入数据因素的机制依据, 效应方向和大小仍由具体响应模型及来源内数据确定. 

\Needspace{8\baselineskip}
第一问提供训练锚点 $\pvec^{(0)}$ 上的条件质量均值 $Q_0$，以及固定 A5 标尺上的响应差 $\psi(\pvec):=\DeltaJ(\pvec)=\widehat J(\pvec)-\widehat J(\pvec^{(0)})$. 第二问按来源 $s$ 估计质量响应 $g_s$；记 $n:=N/10^9$、$d:=D/10^9$、$c:=C/10^{18}$，前三问依次构成
\begin{subequations}\label{eq:unified-chain-v9}
    \begin{align}
\mathcal G_s(Q)&:=e^{-g_s(\ln Q)/\beta_s},\quad
d_{\mathrm{eff},s}:=d\mathcal G_s(Q)e^{-\lambda_{\mathrm{mix}}\psi(\pvec)},\\
\mathcal L_s&:=E_s+A_sn^{-\alpha_s}+B_sd_{\mathrm{eff},s}^{-\beta_s},\\
(n^*,d^*,Q^*)&\in\argmin_{\substack{n,d>0,\ d[kn+h(Q)]\le c\\Q_0\le Q\le1}}
\mathcal L_{B6}(n,d,Q;\pvec^{(0)}).
\end{align}
\end{subequations}
其中 $k:=6+\eta_{\mathrm{att}}L_{\mathrm{ctx}}$，$h(Q):=[c_Q(Q)-c_Q(Q_0)]/10^9$ 为质量处理增量成本，$c_Q$ 与质量响应 $g_s$ 分别描述代价和收益. 锚点数据不能识别配比传递系数 $\lambda_{\mathrm{mix}}$，故第三问主优化固定 $\pvec=\pvec^{(0)}$，其他配比作为条件情景. $Q_0$ 与 B6 质量等级共用标尺也是额外假设. 第四问接收最优配置的训练算力 $C_{\mathrm{train}}=6ND$，再用附件 C 估计能力前沿, 包含处理与注意力开销的总预算不直接作为训练算力使用. 

\subsection{核心理论与方法贡献}
四问的理论链由识别边界、有效数据替代、预算下包络和有序能力前沿组成. 各结论的精确条件与证明放在对应问题的正. 本节说明它们如何支撑题目答案，数值证据见表~\ref{tab:v42-answer-index} 所指位置. \\
\textbf{质量与配比:}
命题~\ref{prop:background-quality-bound} 把未映射领域的完整质量限制在\textbf{可达区间}内. 命题~\ref{prop:q1-quality-collinearity-v15} 说明固定域质量的配比试验不能分别识别质量与配比系数. 第一问因此分别构造质量代理与\textbf{平方根配比响应面}. 对有限候选，命题~\ref{prop:q1-decision-regret-v37} 将一致预测误差传给选择后悔；\textbf{交叉验证 RMSE} 只评价平均预测表现，不能充当该命题的一致误差前提. \\
\textbf{质量与资源的有限替代:}
第二问按来源建立有效数据量. 定理~\ref{thm:q2-log-marginal-v14} 由 Loss 对有效数据量的\textbf{严格单调性}得到质量、Token 和配比的等 Loss 条件. 命题~\ref{prop:q2-finite-substitution-v15} 再给出有限质量提升的等效参数规模及存在条件. \textbf{中心化秩一条件}~\ref{prop:q2-rank-one-v37} 可检验共同质量响应的必要结构. B6 的分组\textbf{误差支持曲率}响应的来源内选型，不证明跨来源质量机制相同. \\
\textbf{预算下包络与有限饱和:}
第三问的关键结构是将质量方案写成关于参数规模 $n$ 的仿射成本族，
$\mathbf{\mathscr H(n)=\min_{Q\in[Q_0,1]}e^{g(\ln Q)/\beta}[kn+h(Q)]}$.
固定质量的唯一资源对由定理~\ref{thm:q3-unique-scale} 保证. 下包络使不同预算的全局解可以直接比较，得到参数量严格增加、质量不下降的定理~\ref{thm:q3-global-monotonicity-v37}，并导出有限预算达到质量上界的充要条件~\ref{thm:q3-finite-saturation-iff-v37}\textbf{ 局部触边}、首次全局转移与充分预算界是不同量，均需与第三问的程序结果逐项核对. \\
\textbf{有序能力前沿:}
第四问用统一六任务得分拟合\textbf{联合分位}. 定理~\ref{thm:q4-ordered-quantile-qp} 在声明域内以有限顶点约束保证分位有序，并将拟合化为可求\textbf{全局解的凸规划}. 推论~\ref{cor:q4-slope-gap-v37} 把数值目标间隙传为\textbf{共享斜率误差界}. 形状与求解保证不等于历史预测有效性，后者由开发期、复用晚期和整组织留出分别检验, 未来点位于声明域外时仅作条件情景解释. 

\subsection{技术路线与理论成果}\label{sec:overview-roadmap}

图~\ref{fig:four-question-roadmap} 将四问组织为“质量识别与配比决策、有限资源替代、预算内全局配置、条件能力前沿”的递进路线. 左列概括核心定理与方法，右列给出定稿程序保存的测试结果，跨问箭头标明实际传递量，各项解析保证与数值证据按其适用范围对应. 

\begin{figure}[H]
\centering
\includegraphics[width=\linewidth]{fig_roadmap.pdf}
\caption{四问的技术路线、核心理论成果与程序测试结果 }
\label{fig:four-question-roadmap}
\end{figure}

\subsection{题目要求与解答概述}\label{sec:answer-index}

表~\ref{tab:v42-answer-index} 把四问的直接要求与本文答案、出处逐条对应，便于核对结论的适用范围；每条答案的证据等级（模型内解析、来源内经验、跨来源条件迁移）在该条备注中保留. 此处只给出结论与出处，推导见相应小节，模型的优缺点与改进方向见第~\ref{sec:conclusion} 节. 

\begin{table}[H]\centering\scriptsize
\setlength{\tabcolsep}{3pt}
\caption{题目要求与本文答案对照及出处表}\label{tab:v42-answer-index}
\begin{tabular}{P{3.8cm}P{8.9cm}P{2.2cm}}\toprule
题目要求 & 本文答案 & 出处\\\midrule
一·质量评分与聚合 & 22 项指标统一为“越高越好”，按模型／DSIR／表面三组加权合成并 clip；域级取文档等权均值，全量去重后 261086 篇、7 个质量域 & 3.1.1，表~\ref{tab:q}、\ref{tab:qtok}\\
一·冲突定义与消解 & $H_i$ 为组间分歧诊断（$0.70/0.30$ 阈值下 A1 为 $2.99\%$），$P_i,A_i$ 为两类污染型矛盾的显式扣分；消融显示惩罚仅改变 $Q_0$，定位为局部修正 & 3.1.2，表~\ref{tab:q1-penalty-ablation}\\
一·配比与 Loss 的关系 & 平方根坐标全二次岭：重复 CV 宏平均 RMSE $0.04453$；A6/A7 上 RMSE $0.037812$、Spearman $0.960201$；选定第 239 行且观测后悔为零 & 3.2，表~\ref{tab:model}、\ref{tab:scale}\\
一·外推稳健性（A12--A15） & 负相关（10B/70B 为 $-0.67/-0.79$）来自跨域水平通道，配比形状通道在五档规模上保持 $+0.68$ 至 $+0.98$；故只支持形状效应迁移，不支持绝对 Loss 预测 & 3.3.1，表~\ref{tab:q1-extrapolation-diagnostic-v41}\\
二·广义标度律 & $\mathcal L_s=E_s+A_sn^{-\alpha_s}+B_sd_{\mathrm{eff},s}^{-\beta_s}$，$d_{\mathrm{eff}}=d\mathcal G_s(Q)e^{-\lambda_{\mathrm{mix}}\psi}$；$Q=1$ 且锚点配比时精确退化为经典形式 & 4.1.2，定义~\ref{def:q2-quality-response}\\
二·边际效用与弹性 & $\gamma(Q_0)=1.88464$、$\gamma(1)=2.40532$；B1 分支的经典指数与文献四位有效数字一致 & 4.3.3，表~\ref{tab:q2-v26-validation}\\
二·质量与规模的替代条件 & 命题~\ref{prop:q2-finite-substitution-v15} 给出有限等价与等效规模；$N=1$B、$D=150$B、$Q_0$ 下质量提高 $0.1$ 等效于旧质量下 $1.23452$B；等效解存在的充要条件为 $H>0$ & 4.2，命题~\ref{prop:q2-finite-substitution-v15}\\
二·领域替代或互补 & 锚点处 17 个单域探针与 136 个双域交互：74 个负交互（互补）、62 个正交互（替代），两模型同向 $112/136$组同向 & 3.3.3，表~\ref{tab:q1-v39-single}、\ref{tab:q1-v39-pairs}\\
三·三档预算的配置 & 曲率响应下 $Q^*=1$，$(N,D)=(1.825,0.568)/(7.102,172.718)/(20.764,6163.27)$（B）；常指数响应给出 $42$ 个上界解与 $3$ 个内点解 & 5.6.1，表~\ref{tab:q3-v26-config}、\ref{tab:q3-v41-transfer}\\
三·$L_{\mathrm{ctx}}^{\mathrm{crit}}$ 与敏感性 & 解析值 $30000$ Token；C7 五个可行值中 $32768$ 为最邻近的高侧取值；五档上下文的配置见表~\ref{tab:q3-v39-context} & 5.1.1，式~\eqref{eq:q3-v39-critical-context}\\
三·结构性转移 & \textbf{条件答案}：名义（B6 指数与曲率响应）下三档间无质量维度转移，结构变化是质量开支份额的系统性退出（$20.75\%\to6.31\%\to2.25\%$）；把质量响应换成常指数形式后首次转移约为 $1.36\times10^{19}$，把规模指数换成真实轨迹估计后 $10^{19}$ 档最优质量退回下界 $Q_0$，两种替代设定下转移都重新出现在区间内 & 5.6.1、5.6.4、5.6.5，表~\ref{tab:q3-v41-transfer}、\ref{tab:q3-exponent-source}\\
三·成本函数的影响 & 三类成本的 $10^{19}$ 档配置：指数 $1.8255/0.5683$、幂 $1.9338/0.5238$、对数 $1.2823/0.9366$（B，$Q^*=1$）；损失口径不同不可横向比较 & 5.6.1，表~\ref{tab:q3-v39-cost}\\
三·数据--参数配比核对 & 三档 $D^*/N^*$ 分别为 $0.311$、$24.3$、$296.8$ token/参数；仅 $10^{22}$ 档与经典算力最优量级 $20$ 一致，$10^{19}$ 档属模型外推而非可执行配方 & 5.6.1，表~\ref{tab:q3-v41-dn-ratio}\\
四·规模与非规模贡献 & 90 日窗口观察性分解：规模 $-2.250630$、时间残差 $+7.832644$、合计 $+5.582014$；有符号净份额 $-40.3\%$ 与 $+140.3\%$（绝对值口径 $22.3\%/77.7\%$） & 6.2.1，式~\eqref{eq:q4-shapley-v9}、图~\ref{fig:q4-decomposition-v11}\\
四·Loss--Benchmark 桥接 & C6 高可比层 7 点留一：均值基线 RMSE $0.424448$ 优于 Loss 线性 $0.552834$ 与 Loss 岭 $0.439516$，据此在高可比层把 Loss 通道权重置零；中可比层两个 Loss 模型优于本层均值基线，但标尺不可与高层合并 & 6.2.2，表~\ref{tab:q4-bridge-v8}\\
四·12/24 月前沿与不确定性                          & 主模型为开发期选出的近 180 日类型经验分位，不含时间趋势，故 12 与 24 月点估计相同：预训练 $33.1968$ 分、后训练 $29.4380$ 分；曲线抽样范围由 $200$ 次有效联合抽样给出，个体带未校准（历史实测覆盖 $50.0\%$--$75.0\%$，低于名义 $90\%$）                                     & 6.5.2，表~\ref{tab:q4-v39-slowdown}、\ref{tab:q4-v39-coverage}                \\
四·算力情景的资源路径                               & 训练算力增长 $\zeta=0/0.5/1$ 对应的资源倍率（以 $C_0=10^{22}$、8192 Token 的 $Q^*=1$ 分支为锚点）：$\zeta=1$ 时 12 月 $N/D$ 倍率 $1.363/2.948$、24 月 $1.878/8.600$；这些是条件资源情景，其能力分数未经长期验证                               & 6.5.1，表~\ref{tab:q4-v39-slow-resources}、\ref{tab:q4-v26-resources}         \\
四·经验分位基准 & 三个入选口径给出不同答案：严格口径 $33.1968/29.4380$（$n=8/4$）、许可去重口径 $18.4544/38.1002$（$n=98/773$）、最宽口径 $19.4413/33.4021$（$n=182/1873$）；后训练上的符号在两种口径间反转，故不能据此断言与经验分位外推相当 & 6.5.2，表~\ref{tab:q4-baseline-caliber}\\
\bottomrule
\end{tabular}
\end{table}

\clearpage

\section{数据选取规则、模型假设与符号约定}\label{sec:data-notation}

在进入具体模型之前, 须先固定哪些量可以相互比较. 附件 A、B、C 的生成机制、观测单位和可信度不同, 半合成关系不能直接解释为真实训练规律, 不同含义的 Loss 与算力也不能直接合并比较. 本章规定证据边界与跨问题统一记号. 通用数学理论集中于附录~\ref{sec:math-background}, 第 3--6 章只引用所需条件与结论. 

\subsection{证据等级与建模约定}
A、B、C 三类附件分别承担质量与配比、标度律、历史能力与算力口径. 本文区分直接观测、半合成、外推与参考映射. 半合成参数和外推结果只支持条件结论, 不因后续解析推导而提升证据等级. 不同附件的 Loss、Benchmark、时间和算力口径仅在正文明确建立桥接时连接. 

\subsection{模型假设与适用范围}
本文只作完成四问所必需的结构假设, 并把不能由附件识别的部分保留为条件参数. \\
\textbf{1. Loss 的可比性分层:} 同一来源、同一验证口径内的交叉熵 Loss 可直接用于拟合和比较. 跨来源 Loss 不预设同一截距或同一质量响应, 而是通过来源下标 $s$ 分层估计. C6 的跨口径桥接只作独立诊断, 不反向改变前三问的 Loss 参数.\\ 
\textbf{2. 质量代理的有限解释:} 文档质量代理 $Q_i$ 经固定标尺归一到 $[0, 1]$, 但未映射领域的质量不作点估计, 只通过部分识别区间进入不确定性分析. 第二问的 $Q$ 因而是模型输入尺度, 不被解释为可直接观测的“真实数据质量”. \\
\textbf{3. 标度律的局部稳定性:} 幂律形式只在各附件的观测支持域及明确标记的外推情景内使用. 第三问主优化固定 $\pvec=\pvec^{(0)}$, 因为锚点数据不足以识别跨来源配比迁移系数 $\lambda_{\mathrm{mix}}$. 配比变化仅作敏感性讨论, 不与质量效应强行联合估计. \\
\textbf{4. 算力成本的可加代理:} 第三问采用题面给定的训练、质量处理与注意力代理成本, 并假设三者作用于同一训练 Token 数 $D$. $L_{\mathrm{ctx}}$ 为外生量, 按 C7 支持范围做敏感性分析. 由该代理模型得到的最优解只在该成本定义下成立. \\
\textbf{5. 能力演进的观察性解释：} 第四问以六项 Benchmark 完整记录的算术均值作为综合能力, 规模—时间分解是拟合值的观察性记账, 不把时间残差解释为已识别的纯算法因果效应. 12/24 月结果均是给定算力增长、资源映射和能力前沿预测模型的条件情景. \\
\textbf{6. 跨来源质量标尺与计量口径的三项叠加假设:} 第二、三问仅在下列三点同时成立时才可联合解释, 任一点不成立都会缩小结论范围：
\setcounter{obsitem}{0}
\obsitem{}第一问六个已映射配比域的条件质量 $Q_0$ 与 B6 半合成质量等级共用同一标尺, 而附件不含重叠锚点, 无法检验.\obsitem{}第一问的配比响应差 $\psi(\pvec)$ 可跨来源迁移, 即 $\lambda_{\mathrm{mix}}$ 虽不可识别但机制不变.\obsitem{}第三问按 Token 计量质量处理代价, 而 $Q_0$ 取自文档等权均值, 其词数加权对应值为 $0.6617$. 第 (3) 点在表~\ref{tab:q3-v41-q0} 作敏感性核查. 

\subsection{统一符号与单位}
\textbf{记号规范:} 本文按表~\ref{tab:notation} 采用“标量斜体、向量粗体、矩阵粗体大写”的对象区分；定义等号写作 $:=$,  概率与期望分别记作 $\Prob$ 与 $\E$, 微分写作直立体 $\dd$； 自然对数记作 $\ln$, 十进制对数明确写作 $\log_{10}$；单侧极限写作 $x\to a^+$ 或 $x\to a^-$； Loss 统一用花体 $\mathcal L$, 上下文长度保留 $L_{\mathrm{ctx}}$；分位水平统一用 $\tau$, 标准化时间变量用 $t$. 概率单纯形、KKT 与分位回归的术语和记号分别参照信息几何、Boyd--Vandenberghe 优化教材与 Koenker 分位回归体系\cite{amari, boyd, koenker}. \par

\begin{table}[H]\centering\footnotesize
\caption{四问共用的数学对象、单位与下标约定}\label{tab:notation}
\begin{tabular}{P{2.9cm}P{8.4cm}P{2.0cm}}\toprule
记号 & 定义与单位 & 所属层次\\\midrule
$i, m, r, a, j, b$ & 文档、原始指标、配方行、配比域、验证域和质量域的索引；局部双域扰动另在该式定义 & 第一问\\
$M_i, U_i, R_i, Q_i$ & 文档的模型组、DSIR 组、表面组及综合质量代理, 均无量纲 & 第一问\\
$\bar Q_b, m_0, \mathcal M, \pi$ & 质量域 $b$ 的文档等权均值；锚点已映射份额；六个有映射的配比域集合及其质量域映射 & 第一、二问\\
$Q_0, Q, Q_{\mathrm{all}}^{(0)}$ & 六映射域在锚点配比下的条件基线；标度律质量输入或第三问的决策量；仅区间识别的锚点完整质量 & 第一至三问\\
$\pvec, \pvec^{(0)}, \boldsymbol\phi$ & $K=17$ 个领域的配比、锚点及平方根坐标；锚点分量为 $p_a^{(0)}$, $\phi_a(\pvec):=\sqrt{p_a}$ & 第一至三问\\
$\delta_{\mathrm{probe}}$ & 第一问锚点配比的有限替换幅度, 取 $0.05$ & 第一问\\
$\bar{\mathcal L}_j, Y_{rj}, \widehat J, \DeltaJ$ & A5 验证域 Loss 的训练折均值、归一化 Loss、平方根响应面及锚点函数差 & 第一、二问\\
$N, D, n, d$ & 参数个数、Token 数；$n:=N/10^9$、$d:=D/10^9$ & 第二至四问\\
$s$ 及来源参数 & $s$ 为数据来源；$(E_s, A_s, B_s)$ 为 Loss 系数, $\alpha_s, \beta_s$ 为规模指数 & 第二、三问\\
$\lambda_{\mathrm{ridge}}, \lambda_{\mathrm{mix}}, \chi_{ND}$ & Q1 岭惩罚、跨来源配比迁移系数、Q2 加性检验的交互系数 & 第一、二问\\
$g_s, \mathcal G_s, a_s, \gamma_s$ & 对数质量响应、有效数据倍率、质量斜率 $a_s=-g_s'$ 及弹性 $\gamma_s(Q)=a_s(Q)/\beta_s$ & 第二、三问\\
$\delta, \delta_{\min}, \kappa_Q$ & 曲率质量函数的端点斜率参数及二次项系数；$\psi:=\DeltaJ$ & 第二、三问\\
$C, c, k, \eta_{\mathrm{att}}, c_Q, h$ & 总 FLOPs, $c:=C/10^{18}$；$k:=6+\eta_{\mathrm{att}}L_{\mathrm{ctx}}$；单位 Token 质量成本 $c_Q$, $h(Q):=[c_Q(Q)-c_Q(Q_0)]/10^9$ & 第三问\\
$w, \mathscr H, n_{\rm env}, c_{\rm tr}$ & 有效数据倍率的倒数、质量成本下包络、固定尺度饱和阈值和首次全局饱和预算 & 第三问\\
$C_{\mathrm{train}}, S, G, T, t, \tau$ & $C_{\mathrm{train}}:=6ND$；六项能力均分；类型 $G$（预训练 P、后训练或微调、基础模型合并 M）；C1 提交日期 $T$、标准化时间和分位水平 $\tau\in(0, 1)$ & 第四问\\
$\lambda_{\mathrm{QR}}, r_C, \kappa_C, \xi_C$ & 联合分位斜率惩罚系数；C4 训练 FLOPs 前沿年增长斜率、算力--参数关系系数及年份系数 & 第四问\\\bottomrule
\end{tabular}
\end{table}

\textbf{损失与风险的术语: } 前三问的 $\mathcal L$ 表示验证交叉熵损失；第四问的 $\rho_\tau(S-a)$ 是单条观测的 Pinball 损失, $\mathrm{PB}_\tau$ 是样本平均 Pinball 损失. 命题~\ref{thm:bg-pinball-quantile} 中的 $R_x(a)$ 是条件期望损失, 也可称风险函数；第三问的 $\operatorname{CVaR}_\tau$ 则是作用于情景后悔值的条件风险价值. 三者的观测对象、单位与优化目标各不相同. 
\section{问题一：数据质量评估与领域配比建模}\label{sec:q1}

第一问包含两类观测单位不同、且无法逐样本匹配的数据. 文档级质量信号用于刻画数据本身的质量特征,  配比实验则用于描述领域比例变化对验证 Loss 的影响. 因此本问先对两类数据分别建模,  再通过可核验的领域映射与共同锚点建立接口. 这样做的目的不是强行构造一个“质量--配比联合回归”, 而是先明确哪些量可识别、哪些量必须留到第二问借助额外质量扰动估计. 

\subsection{完整质量的部分识别与不可识别性}

设第 $r$ 组训练配比为 $\pvec_r=(p_{r1}, \ldots, p_{r17})$, 满足
\begin{equation}\label{eq:simplex}
    p_{ri}\ge 0, \quad \sum_{i=1}^{17}p_{ri}=1.
\end{equation}
质量信号文件给出文档级的 22 项指标, 配方文件给出 17 域训练比例与 13 个可观测验证域的交叉熵损失. 两类文件来自不同语料体系, 不能将各文件 (按领域) 直接逐项拼接. 本文分别建立文档质量评分与配比响应面, 再讨论二者之间能够成立的关联. 

\subsubsection{数据来源、证据层级与比较范围} 质量侧共有 A1 的 51230 条记录、A2 的 17523 条 arxiv 记录和 A3 的 203752 条 github 记录. 其中 A1 的 1419 条 arxiv 与 10000 条 github 已被扩展文件包含, 因此联合统计按 ID 去重, 最终把 A1 作为含原文的校准层、A2/A3 作为相应域的全量层. 配比侧 A4/A5 含 512 组 1M 训练配方 , A6/A7、A8/A9、A10/A11 分别提供 256、256、64 组实测检验. A12--A15 的 10B、70B 仅为已有尺度模型的估算压力测试. 配比与 Loss 只在同一文件对内按 \texttt{index} 对齐, 行和 0.996--1.003 的舍入误差在建模前归一化.

\begin{definition}[已映射份额的条件均值]\label{def:background-mapped-quality}
设 $\pvec=(p_1, \ldots, p_K)$ 为一组总和为 $1$ 的非负配比, $\mathcal M$ 为有领域质量代理分数 $z_a\in[0, 1]$ 的领域集合, 并设已映射份额 $m_{\mathcal M}:=\sum_{a\in\mathcal M}p_a>0$. 定义
\[
Q_{\mathcal M}:=\frac{\sum_{a\in\mathcal M}p_az_a}{m_{\mathcal M}}.
\]
这是已映射份额内部的加权均值. 若其余领域的质量代理分数未知, 它不自动等于完整配方的平均质量. 
\end{definition}

\begin{proposition}[完整质量的部分识别区间]\label{prop:background-quality-bound}
在定义~\ref{def:background-mapped-quality} 的条件下, 若未映射领域只知 $z_a\in[0, 1]$, 则完整配方平均质量 $Q_{\mathrm{all}}:=\sum_a p_az_a$ 必满足
\begin{equation}\label{eq:background-quality-bound}
m_{\mathcal M}Q_{\mathcal M}
\le Q_{\mathrm{all}}\le
m_{\mathcal M}Q_{\mathcal M}+1-m_{\mathcal M}.
\end{equation}
令所有未映射代理为 $0$ 或 $1$ 即可达到左右两个端点；该区间是信息缺失下的识别范围, 不是置信区间. 
\end{proposition}
\begin{proof}
把 $Q_{\mathrm{all}}$ 分解为已映射项 $m_{\mathcal M}Q_{\mathcal M}$ 与未映射项 $\sum_{a\notin\mathcal M}p_az_a$. 后者在 $0$ 与 $\sum_{a\notin\mathcal M}p_a=1-m_{\mathcal M}$ 之间, 故得式~\eqref{eq:background-quality-bound}. 
\end{proof}


\subsubsection{一般约束下的质量识别集}
为掌握有依据的域质量区间或校准关系, 令这些约束构成非空紧集 $\mathcal U\subseteq[0, 1]^K$, 定义
\begin{equation}\label{eq:q1-general-identification-v37}
Q_-(\pvec):=\min_{\boldsymbol z\in\mathcal U}\pvec^\top\boldsymbol z, 
\quad
Q_+(\pvec):=\max_{\boldsymbol z\in\mathcal U}\pvec^\top\boldsymbol z.
\end{equation}
连续性保证两个端点都能达到. 若 $\mathcal U$ 还是凸集, 则线性函数的值域 $[Q_-, Q_+]$. 当约束为线性时, 两端点由线性规划计算. 命题~\ref{prop:background-quality-bound} 是仅有未知域约束时的特例, 属于部分识别而非补齐缺失真值\cite{manski_partial}. 没有新的校准信息时, 不能仅依靠改变聚合公式缩窄该范围, 其进入第二问之前还需检查质量坐标是否一致.

\subsubsection{配比与质量效应的识别边界}
A16 对 17 个配方域仅给出 3 个直接映射、3 个近似映射. 其余 11 个域没有质量侧同名数据, 但部分可以在额外语义假设下使用代理域. 本文以附件明示的六域映射为主口径, 将扩展映射作为条件敏感性情景, 区分数据可得性与映射规则的选择. 

\begin{proposition}[固定域质量下的配比--质量系数不可识别]\label{prop:q1-quality-collinearity-v15}
若所有配方 $r$ 的第 $i$ 域质量恒为 $z_i$, 则模型中的两列 $p_{ri}$ 与 $p_{ri}z_i$ 线性相关；观测损失至多识别合并系数 $\beta_i+z_i\theta_i$, 不能分别识别配比系数 $\beta_i$ 和质量交互系数 $\theta_i$. 
\end{proposition}
\begin{proof}
对每一行 $r$, 有 $\beta_i p_{ri}+\theta_i p_{ri}z_i=(\beta_i+z_i\theta_i)p_{ri}$, 故设计矩阵中对应两列满足 $\mathbf x_2=z_i\mathbf x_1$. 由命题~\ref{prop:bg-linear-identifiability}, 设计矩阵存在非零零空间方向. 具体地, 对任意实数 $t$, 把 $(\beta_i, \theta_i)$ 改为 $(\beta_i-tz_i, \theta_i+t)$ 后预测不变, 因此两系数没有唯一数据解. 
\end{proof}
据此, 本问分开构造质量代理与配比响应. 质量效应由第二问的 B6 质量扰动另行估计, 并保留其半合成来源标记. 

\subsubsection{质量信号与综合评分}

\textbf{质量信号、方向与标尺: }
令 $x_{im}$ 为文档 $i$ 在原始指标 $m$ 上的值. 八个列表字段先按官方数据卡的类别顺序压缩\cite{metacard, metarater}. 四组 ModernBERT 六级 logits $z_{im\ell}$ 转为连续等级
\begin{equation}\label{eq:logit}
    r_{im}=\frac{1}{5}\sum_{\ell=0}^{5}\ell
    \frac{\exp(z_{im\ell}-z_{im}^{\max})}{\sum_{k=0}^{5}\exp(z_{imk}-z_{im}^{\max})}.
\end{equation}
二分类 \texttt{fluency\_en} 与 \texttt{ad\_en} 均取第 1 类概率 $\sigma(z_1-z_0)$, 分别代表流畅和无广告. \texttt{fineweb\_edu} 取唯一数值；\texttt{qurater} 四项依次为写作风格、所需专业度、事实知识、教育价值, 来自 QuRating 式的主观质量维度\cite{qurating}, 本文以四项均值作为该原始指标的一维摘要. 所需专业度的高低具有任务依赖性, 此处将其视为通用预训练的一个弱代理, 后文检验权重敏感性. 

DSIR 三项原值明显随文本长度改变, 故以 $\max\{w_i, 20\}$ 个词作近似长度校正. 词数与句数先作 $\ln(1+x)$ 变换. 以 A1 的有限值估计每项的 $1\%$、$99\%$ 分位数 $a_m, b_m$, 定义
\begin{equation}\label{eq:normalize}
    v_{im}=\min\!\left\{1,  \max\!\left\{0,  \frac{r_{im}-a_m}{b_m-a_m}\right\}\right\}.
\end{equation}
教育、流畅、无广告、四组 ModernBERT、QuRating、三组 DSIR,  以及词数、句数、词熵、独特词比例、终止标点比例取 $u_{im}=v_{im}$. 无字母词比例、最高频二元及三元片段比例、大写字母比例取 $u_{im}=1-v_{im}$. 数字字符比例与平均词长没有普遍单调方向,  取中间区间适宜度 $u_{im}=1-|2v_{im}-1|$. 由此所有 $u_{im}$ 都在 $[0, 1]$ 内, 并按``越高越好''解释. 方向与适宜区间是模型假设, 并非附件提供的真值标签. A1 的 18 条记录及 A3 的 1 条记录各有一个 ModernBERT 列表为非有限值, 仅该项用 A1 对应指标的中位数补齐, 其余指标均有效. 固定 A1 标尺后不在 A2/A3 重新拟合. 

将八项模型评分、三项 DSIR 与十一项文本表面信号分别记为 $M_i, U_i, R_i$ 的组内等权平均. 这里先把“连续质量证据”和“明确冲突惩罚”分开. 连续部分给出基础得分, 惩罚项只处理少数满足预先规则的矛盾样本. 等组权版本为 $Q_i^{\rm eq}=(M_i+U_i+R_i)/3$. 因 DSIR 更接近参考语料重要度而非通用质量标签, 主评分降低其权重：
\begin{subequations}\label{eq:quality}
   \begin{equation}
           Q_i=\operatorname{clip}_{[0, 1]}\left(0.50M_i+0.15U_i+0.35R_i-0.25P_i[M_i-R_i-0.40]_+-0.05A_i\right), 
   \end{equation}
   \begin{equation}
        P_i=\mathbf 1_{\{M_i\ge0.70, \ R_i\le0.30\}}, \quad A_i=\mathbf 1_{\{u_{i, \mathrm{edu}}\ge0.70, \ u_{i, \mathrm{noad}}\le0.30\}}.
   \end{equation}
\end{subequations}
式~\eqref{eq:quality} 可以理解为
\[
\text{综合质量代理}=\operatorname{clip}\{\text{三组连续证据的加权和}-\text{规则化冲突惩罚}\}.
\]
因此权重决定三类证据的相对贡献, $P_i, A_i$ 只决定特定矛盾样本是否额外扣分. 二者在统计含义上应分开做敏感性分析. 令 $\mathcal I_b$ 为质量域 $b$ 的文档集合、$n_b:=|\mathcal I_b|$, 域级质量取 $\bar Q_b=n_b^{-1}\sum_{i\in\mathcal I_b}Q_i$. 标准误差按同域文档的样本标准差除以 $\sqrt{n_b}$ 计算. 质量评分是明确口径下的代理量, 不等同人工标注的真实质量. 

\begin{table}[H]\centering\small
\caption{七个质量域的综合评分及扩展样本对照表}\label{tab:q}
\begin{tabular}{lrrrrr}\toprule
质量域 & A1 条数 & A1 平均 $Q$ & A1 标准误 & 扩展条数 & 扩展平均 $Q$\\\midrule
arxiv & 1419 & 0.8206 & 0.0010 & 17523 & 0.8209\\
book & 171 & 0.7011 & 0.0051 & -- & --\\
c4 & 10000 & 0.6118 & 0.0010 & -- & --\\
commoncrawl & 9640 & 0.6612 & 0.0008 & -- & --\\
github & 10000 & 0.5049 & 0.0009 & 203752 & 0.5041\\
stackexchange & 10000 & 0.6330 & 0.0007 & -- & --\\
wikipedia & 10000 & 0.5718 & 0.0009 & -- & --\\\bottomrule
\end{tabular}
\end{table}

表~\ref{tab:q} 显示,  两域抽样与扩展样本的平均评分高度接近,  其中 arxiv 仅相差 $0.0003$,  github 仅相差 $-0.0008$. 这一结果表明,  所构建的评分体系在两个同源嵌套样本上具有较好的均值稳定性与一致性. 由于扩展集包含相应 A1 记录,  本文将其作为同源样本扩展下的一致性检验.进一步地,  将三组权重调整为各 $1/3$,  或完全去除 DSIR 后,  book、c4 等域的绝对评分会出现一定变化. 这一敏感性分析说明,  模型能够有效反映不同评分口径对各领域质量水平的影响,  同时主要领域的评分结果具有清晰、可解释的结构响应,  从而增强了综合质量指标的可分析性与可复核性.


域级质量 $\bar Q_b$ 以文档为等权单位. 由于第三问按 token 计量质量处理的代价, 本文另报告以原文词数 $w_i$（\texttt{rps\_doc\_word\_count} 的未变换原值, 作为 token 体量的近似）为加权均值
\[\bar Q_b^{\mathrm{tok}}=\frac{\sum_{i\in\mathcal I_b}w_iQ_i}{\sum_{i\in\mathcal I_b}w_i}, \]
%两种口径都保留. 
表~\ref{tab:qtok} 给出按文档 ID 去重后七个域的两种结果：去重表保留 A2/A3 中的记录, 故 A1 中已包含于 A2、A3 的 1419 条 arxiv 与 10000 条 github 记录不重复计入, 七域合计 261086 条. 文档等权列与表~\ref{tab:q} 中相应域的均值一致, 词数加权列是新增口径. 

\begin{table}[H]\centering\small
\caption{按文档 ID 去重后七个质量域的文档等权均值与原文词数加权均值(保留 A2/A3 )}\label{tab:qtok}
\begin{tabular}{lrrrr}\toprule
质量域 & 去重条数 & 文档等权 $\bar Q_b$ & 词数加权 $\bar Q_b^{\mathrm{tok}}$ & 差值\\\midrule
arxiv & 17523 & 0.8209 & 0.8250 & $+0.0041$\\
book & 171 & 0.7011 & 0.7145 & $+0.0134$\\
c4 & 10000 & 0.6118 & 0.6635 & $+0.0518$\\
commoncrawl & 9640 & 0.6612 & 0.6926 & $+0.0313$\\
github & 203752 & 0.5041 & 0.5683 & $+0.0641$\\
stackexchange & 10000 & 0.6330 & 0.6600 & $+0.0270$\\
wikipedia & 10000 & 0.5718 & 0.6579 & $+0.0861$\\\bottomrule
\end{tabular}
\end{table}

词数加权后,  七个质量域的平均得分均高于文档等权结果,  但不同领域的变化幅度有所差异. 其中,  wikipedia 由 $0.5718$ 升至 $0.6579$（$+0.0861$）,  github 由 $0.5041$ 升至 $0.5683$（$+0.0641$）,  C4 由 $0.6118$ 升至 $0.6635$（$+0.0518$）,  而 arxiv 由 $0.8209$ 升至 $0.8250$（$+0.0041$）；在扩展表中单独计算时,  A3 的 github 升幅同样为 $+0.0641$. 这一结果表明,  文档等权与 token 加权分别刻画不同层面的平均质量. 前者以文档为分析单位,  描述典型文档的平均质量. 后者以 token 为分析单位,  描述典型 token 所处语料的平均质量. 两种口径相互补充,  使质量评价能够同时覆盖文档层面与语料规模层面的特征.因此第一问主体分析采用文档等权口径,  以保持文档级质量评价的一致性. 第三问涉及按 token 计量的质量处理成本时,  进一步引入 token 加权结果作为补充分析,  从而使质量评价口径与后续成本计量保持对应.



\textbf{评分稳健性与域内分辨力:}

组权属于模型假设, 故在保持两类显式惩罚规则不变的条件下, 比较主权重 $(0.50, 0.15, 0.35)$、三组等权、去除 DSIR、模型主导、表面信号主导和三个单组极端口径, 共八种设定. 主口径与等权、去除 DSIR、模型主导、表面主导四种设定的七域排序 Spearman 相关依次为 $0.964$、$1.000$、$0.964$ 和 $0.964$；与纯模型、纯 DSIR、纯表面三种极端口径的相关依次为 $0.857$、$0.786$ 和 $0.929$. arxiv 在除纯 DSIR 外的七种设定下居首, 纯 DSIR 下由 book 居首. github 在除纯表面外的七种设定下居末, 纯表面下由 wikipedia 居末. 因此极端单组口径会改变首末域, 不宜把主口径排序说成与组权无关. 表 \ref{tab:7-filed-score} 展示了这七个域的分数和名词变化

\begin{table}[H]
    \centering
    \caption{固定惩罚项后改变模型、DSIR 与文本表面三组权重时的质量排名}
    \renewcommand{\arraystretch}{1.3}
    \begin{tabular}{l *{7}{>{\centering\arraybackslash}p{1.2cm}}}
        \toprule
         & arxiv & book & common\-crawl & c4 & stack\-exchange & github & wiki\-pedia \\
        \midrule
        main\_50\_15\_35         & 1 & 2 & 3 & 5 & 4 & 7 & 6 \\
        equal\_third             & 1 & 2 & 3 & 4 & 5 & 7 & 6 \\
        no\_dsir\_60\_00\_40     & 1 & 2 & 3 & 5 & 4 & 7 & 6 \\
        model\_led\_70\_10\_20   & 1 & 2 & 4 & 5 & 3 & 7 & 6 \\
        surface\_led\_30\_10\_60 & 1 & 2 & 3 & 4 & 5 & 7 & 6 \\
        model\_only              & 1 & 3 & 4 & 6 & 2 & 7 & 5 \\
        dsir\_only               & 3 & 1 & 2 & 4 & 6 & 7 & 5 \\
        surface\_only            & 1 & 2 & 3 & 4 & 5 & 6 & 7 \\
        \bottomrule
    \end{tabular}
    \label{tab:7-filed-score}
\end{table}

在 A1 的 $51230$ 篇文档上按质量域进行单因素平方和分解, $Q_i$ 的总平方和为 $580.893$, 其中域内 $363.705$（$62.6\%$）, 域间 $217.188$（$37.4\%$）. 该恒等分解表明 $Q_i$ 不只是域标签的替身, 域内仍有变异. 它不能证明域内高分文档能改善模型训练, 后者仍需独立训练或人工盲评. 图~\ref{fig:q-groups} 展示三组信号在各域内的分布，其中虚线对应组间冲突诊断阈值. 

\begin{figure}[htbp]\centering
\includegraphics[width=0.94\linewidth]{q1_group_score_distributions.pdf}
\caption{A1 七个质量域在模型评分 $M_i$、DSIR $U_i$ 和文本表面评分 $R_i$ 上的分布 }\label{fig:q-groups}
\end{figure}
%图~\ref{fig:q-groups} 展示 A1 七个质量域在模型评分 $M_i$、DSIR 评分 $U_i$ 和文本表面评分 $R_i$ 上的分布,  虚线表示组间冲突的诊断阈值.
\textbf{表面指标适宜度函数的敏感性分析:}\label{sec:q1-function-sens-v8}

以原始平均词长和数字字符比例为对象, 固定其余二十项指标转换、组权 $(0.50, 0.15, 0.35)$ 与惩罚结构, 比较含主口径在内的十一组函数情景. 平均词长 $x$ 的候选平台在 $3\le x\le10$ 取满分, 低于 3 时取 $x/3$, 高于 10 时以宽度 $s=5$ 或 10 线性衰减到零. 数字字符比例在 A1 同域样本的 80 分位以下不扣分, 80 至 95 分位线性衰减, 95 分位以上取零. 阈值均仅从 A1 拟合, 扩展集 A2/A3 不参与定阈值. 平均词长平台借鉴 Gopher 对网页文本采用的宽接受区间\cite{gopher}. 平台外斜率以及数字比例的域分位是本文试算假设, 并非质量真值. 

表~\ref{tab:q1-function-sens-v8} 列出全部情景. $Q_0$ 仍按六个已映射配比域、锚点配比和文档等权均值计算, 其覆盖的锚点份额为 $31.94\%$. 全部情景均按 ID 去重为 261086 篇, 十一组情景中七域排序的 Spearman 相关均为 1, 但逐文档分数和 $Q_0$ 的绝对标尺会改变, 其中联合平台方案与主口径的逐文档相关最低. 因而“排序稳定”不等于“质量评分被外部证实”. 
\begin{table}[htbp]\centering\small
\caption{两项文本表面函数的十一组独立试算}\label{tab:q1-function-sens-v8}
\begin{tabular}{lrr}\toprule
情景 & $Q_0$ & 与主口径逐文档 Spearman\\\midrule
现行主口径 & 0.607474 & 1.000000\\
两项三角峰移至归一化值 0.6 & 0.605174 & 0.999399\\
两项三角峰移至归一化值 0.7 & 0.603435 & 0.998314\\
仅词长平台, $s=5$ & 0.630529 & 0.995200\\
仅数字比例域内 80/95 分位 & 0.628034 & 0.990344\\
联合：$s=5$、域内 80/95 分位 & 0.651087 & 0.988656\\
联合：$s=10$、域内 80/95 分位 & 0.651250 & 0.989524\\
联合：$s=5$、全局 80/95 分位 & 0.649757 & 0.990929\\
联合：$s=5$、域内 75/90 分位 & 0.649301 & 0.986049\\
联合：$s=5$、域内 85/97.5 分位 & 0.652371 & 0.990585\\
删除两项表面信号 & 0.629372 & 0.994303\\
\bottomrule\end{tabular}\end{table}
%表~\ref{tab:q1-function-sens-v8} 列出两项文本表面函数的十一组独立试算. 所有情景均按 ID 去重为 261086 篇文档. $Q_0$ 仍按照六个已映射配比域、锚点配比和文档等权均值计算,  对应的锚点份额覆盖为 $31.94\%$.十一组情景下七个质量域的排序 Spearman 相关系数均为 $1$,  表明领域相对排序对函数形式具有较高稳定性. 不同函数设定会对逐文档分数和 $Q_0$ 的绝对数值产生一定影响,  其中联合平台方案与主口径的逐文档相关最低. 因此,  本文同时报告排序一致性与绝对标尺变化,  从相对排序和数值尺度两个层面对质量评分的稳健性进行检验.

联合替代使 $Q_0$ 在表~\ref{tab:q1-function-sens-v8} 的十一组试算中取得最高值, 是评分函数变换导致的尺度位移, 不代表真实训练质量提升, 本次没有新增人工盲评或等算力下游训练. 主评分规则和平方根二次配比响应面保持冻结, 故第二至四问的正式计算继续使用表~\ref{tab:q1-penalty-ablation} 的主口径 $Q_0$, 不把本节试算分数当作已验证的新输入. 

\subsubsection{质量冲突、领域映射与核验}

\textbf{组间冲突与显式惩罚: }
对文档 $i$ 的模型、DSIR 和表面三组质量代理 $(M_i, U_i, R_i)$, 定义组间分歧指示函数
\[
H_i=\mathbf 1_{\{\max(M_i, U_i, R_i)\ge0.70, \ \min(M_i, U_i, R_i)\le0.30\}}.
\]
$H_i$ 只表示三个代理相互不一致. 式~\eqref{eq:quality} 的 $P_i, A_i$ 才触发评分中的显式惩罚. 因此分歧诊断与评分惩罚对应两个不同集合. 

按上述冲突标记, A1 的 c4 域冲突率为 $9.62\%$, github 为 $1.79\%$, arxiv 为 $0.49\%$；在扩展集中, github 为 $1.94\%$, arxiv 为 $0.26\%$. 多数 $H_i=1$ 情形的最高组是 DSIR、最低组是模型评分, 例如 A3 github 有 3810 条属于此型. 它可能代表参考域相似性与通用语言评分的差异, 不能直接判为污染. 式~\eqref{eq:quality} 仅对 $P_i$、$A_i$ 两种更明确的污染型矛盾施加惩罚, 复杂推理文本与低可读性并存也只作分歧记录. 

指示函数 $H_i$ 的高低阈值取 $0.70$ 与 $0.30$ 是本稿的建模选择, 故表~\ref{tab:q1-penalty-ablation} 一并报告标记结果对阈值的敏感性：阈值放宽 $0.05$ 使标记比例约增至 $2.1$ 倍, 收紧 $0.05$ 则降至约 $0.42$ 倍, 可见``冲突''的绝对水平高度依赖阈值设定, 跨表比较只应在同一阈值下进行, 上文各域的冲突率均对应 $0.70/0.30$. 显式惩罚条件由式~\eqref{eq:quality} 单独定义, 不取决于 $H_i$ 的诊断阈值：三个阈值下均有 33 条记录满足惩罚条件, 其中 24 条位于 A1、9 条位于 A3（计数包含重复记录）. 在主阈值下, $H_i$ 标记集合与惩罚集合的交集为 12 条：A3 的 9 条 github 记录全部落入交集, A1 仅 3 条. 

两类显式惩罚与诊断标记并不嵌套,  $P_i$ 的条件（$M_i\ge0.70$ 且 $R_i\le0.30$）自动满足 $H_i=1$, 故 $P_i$ 记录全在交集内, 而 $A_i$ 的条件（教育价值高、无广告）不涉及三组信号的组间极差, 其记录中有相当一部分三组得分都落在 $[0.30, 0.70]$ 内, 因而 $A_i=1$ 而 $H_i=0$. 因此``冲突诊断''与``冲突惩罚''在本文中是两套不重合的规则, 这一点在下文的消融实验中保留. 其余未受惩罚的组间分歧只由固定组权合成, 不附加文档级扣分. 

\textbf{显式惩罚的消融: \quad}由于式~\eqref{eq:quality} 

的两类惩罚只作用于极少数记录, 须报告其对最终评分与下游输入的实际影响, 否则"冲突消解"只是公式上的机制. 在全量记录上按 ID 去重后的 261086 篇文档上分别计算四种情形, 四种情形下的七域均值与域间排序完全一致, 故正文只报告锚点量, 结果见表~\ref{tab:q1-penalty-ablation}. 

\begin{table}[H]\centering\small
\caption{显式冲突惩罚项的消融结果}\label{tab:q1-penalty-ablation}
\begin{tabular}{lrrrr}\toprule
量 & 两项均保留 & 仅保留 $A_i$ & 仅保留 $P_i$ & 两项均去掉\\\midrule
受惩罚记录数 & $33$ & $21$ & $12$ & $0$\\
$Q_0$（锚点） & $0.6074742$ & $0.6074768$ & $0.6074789$ & $0.6074814$\\\bottomrule
\end{tabular}
\end{table}
表~\ref{tab:q1-penalty-ablation} 给出显式冲突惩罚的消融结果,  各列标题表示对应设定中保留的惩罚项. 受惩罚文档数按包含 A1/A2/A3 重复记录的 272505 条原始记录统计, $Q_0$ 则基于 261086 条去重记录,  按六个已映射配比域及锚点配比计算. 各质量域的具体结果见表~\ref{tab:q}.
消融给出一个明确但不利于"惩罚是主要机制"的结论：去掉两类惩罚后 $Q_0$ 由 $0.6074742$ 变为 $0.6074814$, 
变化量 $7.2\times10^{-6}$（相对 $1.2\times10^{-5}$）；七域均值中变化最大的是 c4 与 wikipedia, 
也仅为 $8.5\times10^{-5}$ 与 $5.0\times10^{-6}$, 域间排序在四种情形下完全不变. 原因是可核验的：
在去重后的 261086 篇文档上, 满足 $P_i$ 的记录 $11$ 条、满足 $A_i$ 的记录 $21$ 条, 两类无交集, 
合计 $32$ 条, 占 $1.23\times10^{-4}$；若按含 A1/A2/A3 重复记录的 272505 条原始记录计, 
对应 $12$、$21$、$33$ 条, 占 $1.21\times10^{-4}$. 无论采用哪种计数口径, 受影响记录都不超过万分之一点三. 
因此本文把式~\eqref{eq:quality} 的定位明确为\textbf{诊断指标 $H_i$ 之上的局部修正}, 
而不是质量评分的主要来源：$H_i$ 用于刻画组间分歧（$0.70/0.30$ 阈值下 A1 为 $2.99\%$、扩展集 A2+A3 为 $1.81\%$、三条文件合计为 $2.03\%$）, 
$P_i$、$A_i$ 只对最明确的两类污染型矛盾施加扣分. 若希望冲突消解在数值上真正参与评分, 
需要把离散惩罚改为按分歧幅度连续的稳健聚合；
主口径不采用该改写, 以免在无人工真值的情况下改变已验证的评分标尺, 下文将其列为条件试算.


\textbf{扩展域映射与连续冲突权重的条件试算: \quad}\label{par:q1-review-v44}

为量化上述建模选择, 保留 A16 六域为正式口径, 另将十个未映射域作纯语义代理压力测试：\texttt{dm\_mathematics}、\texttt{pubmed\_central}、\texttt{pubmed\_abstracts} 指向 arxiv, \texttt{philpapers}、\texttt{freelaw} 指向 book, \texttt{hackernews}、\texttt{europarl}、\texttt{uspto\_backgrounds}、\texttt{enron\_emails}、\texttt{ubuntu\_irc} 指向 commoncrawl；\texttt{nih\_exporter} 保持未映射. 这些跨体裁对应没有新的质量真值, 后五项尤其只是宽泛网页代理, 不能提升实测覆盖率. 该情景的条件映射份额为 $0.997006$、$Q_0=0.692706$；若所指定的代理质量值全部准确, 剩余未映射域仅约束在 $[0, 1]$, 完整质量的算术区间为 $[0.690632, 0.693626]$, 不能把这个较窄区间解读为由附件本身识别出的真实质量. 

另以三组信号的极差 $s_i:=\max(M_i, U_i, R_i)-\min(M_i, U_i, R_i)$ 构造连续综合评分情景
\[
Q_i^{(\alpha)}=(1-\alpha s_i)(0.50M_i+0.15U_i+0.35R_i), 
\quad \alpha\in\{0.1, 0.25, 0.5\}.
\]
它对全部 $261086$ 篇去重文档施加幅度相关的折减, 并替代式~\eqref{eq:quality} 的两项离散惩罚. 这是无人工真值条件下的结构敏感性, 不是重新校准的主评分. 三档 $\alpha$ 的七域排序相对主口径 Spearman 均为 $1$, 但 $Q_0$ 依次降至 $0.588713, 0.560559, 0.513637$.

图~\ref{fig:q-conflict} 将诊断标记与显式惩罚绘于不同信号平面. A1 的组间 Spearman 秩相关中, $M_i$--$U_i$ 为 $0.382$, $M_i$--$R_i$ 为 $0.288$, $U_i$--$R_i$ 为 $0.625$. 这种相关结构与大量 DSIR 高、模型组低的分歧一致. 图中星号稀少, 也说明质量评分的惩罚是针对明确条件的局部修正, 而非对所有 $H_i=1$ 样本统一扣分. 

\begin{figure}[htbp]\centering
\includegraphics[width=0.94\linewidth]{q1_conflict_geometry.pdf}
\caption{A1 文档三组质量信号的关系与冲突标记图 }\label{fig:q-conflict}
\end{figure}
图~\ref{fig:q-conflict} 展示 A1 文档三组质量信号之间的几何关系. 左图以模型组和 DSIR 组为两个坐标轴,  颜色表示文本表面组, 右图区分组间分歧标记与受到显式惩罚的文档.
A1 原文可用于局部核验. 按域比较 $Q$ 最高、最低及组间信号分歧较大的文档后, 低分 github 示例多为 XML/HTML 模板片段, 高分 arxiv 和 c4 示例多为连贯的学术段落. 另从 A1 原文计算三项事后诊断指标：重复长行比例、广告短语标记和至少三个 URL 标记. 这三项均不参与 $Q$ 的计算, 也不参与质量评分算法. 

将 A1 的 51230 条记录在各域内按 $Q$ 分为四组, 再合并同分位记录, 四组文档数依次为 12808、12808、12806、12808. 重复长行比例的四组均值依次为 $2.318\%$、$2.007\%$、$2.339\%$、$2.301\%$,  广告短语和至少三个 URL 的标记率也不呈一致的单调变化. 因此, 这些事后诊断不能作为 $Q$ 有效性的支持性证据. 另一方面, 域间次序在 11 组函数情景中不变, 域内平方和仍占总平方和的 $62.6\%$；故 $Q$ 可作为本附件内域间相对排序的条件代理, 不能据此确认为文档级质量真值或跨来源训练收益. 文本质量仍需独立人工盲评或下游训练实验核验；A2/A3 不含原文, 无法复算这三项诊断指标. OpenScholar 的专项知识问答结果说明检索证据与训练数据组织会影响具体任务表现\cite{asai_openscholar}, 但不单独识别本文 $Q$ 的因果作用. 

%\subsection{Fisher--Rao 配比响应与选择后悔界}

上一节得到的是文档质量及其有限领域映射, 但无法回答“改变十七域配比后验证 Loss 如何变化”. 为避免把缺失质量域纳入进回归, 本文转而直接使用 A4/A5 的配比--Loss 实验建立独立响应面, 信息几何只负责选择适合单纯形约束的坐标, 不替代统计验证. 

\subsubsection{平方根响应面与模型选择}

A4/A5 给出 $512$ 组十七域训练配方及其在十三个验证域上的 Loss. 令 $\bar{\mathcal L}_j$ 为 A5 第 $j$ 个验证域 Loss 的训练折均值, 并定义 $Y_{rj}:=\mathcal L_{rj}/\bar{\mathcal L}_j$. 本文以
\begin{equation}\label{eq:objective}
J(\pvec_r)=\frac1{13}\sum_{j=1}^{13}Y_{rj}, \qquad
\pvec_r\ge0, \quad\sum_{i=1}^{17}p_{ri}=1
\end{equation}
作为配比选择的等权目标, 同时对十三个 $Y_{rj}$ 分别拟合. 这里的归一化先消除十三个验证域自身 Loss 量级差异, 再取等权平均, 因此 $J(\pvec)$ 比较的是“相对各域典型水平的平均损失”, 而不是让数值较大的验证域自动占据更大权重. 折内的 $\bar{\mathcal L}_j$ 只从该折训练行计算, 避免验证行参与目标归一化. 


\subsubsection{信息几何坐标与平方根响应面}

由定义~\ref{def:background-simplex} 与定理~\ref{thm:bg-sqrt-isometry}, 配比平方根坐标
$\boldsymbol\phi(\pvec):=(\sqrt{p_1}, \ldots, \sqrt{p_{17}})^\top$
落在单位球面正象限, 并与 Fisher--Rao 局部几何相容. 因此本文用平方根坐标构造响应面. 这一几何依据只说明坐标选择合理, 预测优劣仍由交叉验证决定. 

置 $\phi_i=\sqrt{p_i}$. 本文的主响应面是平方根坐标的全二次岭回归, 其标准化特征经训练折拟合后, 可在原始配比上写成
\begin{equation}\label{eq:sqrtquad}
\widehat Y_j(\pvec)=b_{j0}
 +\sum_{i=1}^{17}a_{ji}\sqrt{p_i}
 +\sum_{i=1}^{17}d_{ji}p_i
 +\sum_{i<k}c_{jik}\sqrt{p_i p_k}.
\end{equation}

为区分模型结构与领域效应, 令 $\boldsymbol\phi(\pvec)=(\sqrt{p_1}, \ldots, \sqrt{p_{17}})^\top$、$\widehat J(\pvec)=13^{-1}\sum_{j=1}^{13}\widehat Y_j(\pvec)$, 并定义 $\DeltaJ(\pvec):=\widehat J(\pvec)-\widehat J(\pvec^{(0)})$. 十三域平均响应可写成
\begin{equation}\label{eq:q1-matrix-v9}
\widehat J(\pvec)=b_0+\boldsymbol a^\top\boldsymbol\phi(\pvec)
+\boldsymbol\phi(\pvec)^\top\mathbf B\boldsymbol\phi(\pvec), \qquad \mathbf B=\mathbf B^\top.
\end{equation}
式~\eqref{eq:q1-matrix-v9} 把响应面分成三个层次：$b_0$ 是共同基线, $\boldsymbol a^\top\boldsymbol\phi$ 描述平方根坐标中的一阶方向, $\boldsymbol\phi^\top\mathbf B\boldsymbol\phi$ 描述曲率与域间联合变化. 由于这些项是在球面而不是整个 $\mathbb R^{17}$ 上观测, 系数表示并不唯一. 单纯形约束给出 $\phi_a\ge0$、$\sum_a\phi_a^2=1$, 即平方根坐标落在单位球面的正象限部分. 因而对任意实数 $t$, $(b_0, \mathbf B)$ 与 $(b_0-t, \mathbf B+t\mathbf I)$ 在全部可行配比上给出同一函数. 这正是定义~\ref{def:bg-identifiability} 所述的参数不可识别：存在非零参数方向而预测函数在全部可行设计点上不变. 岭惩罚确定一个数值表示, 但单个对角系数并非可识别的领域效应. 

\textbf{模型性质：球面约束下的一阶梯度}
在球面内部点, 令 $f(\boldsymbol\phi)=b_0+\boldsymbol a^\top\boldsymbol\phi+\boldsymbol\phi^\top\mathbf B\boldsymbol\phi$. 其欧氏梯度为 $\nabla f=\boldsymbol a+2\mathbf B\boldsymbol\phi$. 单位球面在 $\boldsymbol\phi$ 处的一阶可行方向满足 $\boldsymbol\phi^\top\boldsymbol z=0$. 由命题~\ref{prop:bg-orthogonal-projection}, 约束梯度为
\begin{equation}\label{eq:q1-riem-grad-v13}
\operatorname{grad} f(\boldsymbol\phi)
=(\mathbf I-\boldsymbol\phi\boldsymbol\phi^\top)
(\boldsymbol a+2\mathbf B\boldsymbol\phi).
\end{equation}
若该向量非零, 则切向量中的最陡下降方向是
$-\operatorname{grad} f/\|\operatorname{grad}f\|_2$. 式~\eqref{eq:q1-riem-grad-v13} 因而只保留不改变 $\|\boldsymbol\phi\|_2$ 的切向分量；径向分量不对应可行的配比一阶变化. 本文仍只在已观测候选集内选配方, 以避免把几何上的可优化性解释为统计上的可外推性. 

训练折内对连续多项式特征逐列标准化, 截距不惩罚, 岭估计的目标为
\begin{equation}\label{eq:q1-ridge-v9}
(\widehat b_j, \widehat{\boldsymbol\theta}_j)=\argmin_{b, \boldsymbol\theta}
\bigl\|\boldsymbol Y_j-b\one-\widetilde{\mathbf X}\boldsymbol\theta\bigr\|_2^2
+\lambda_{\mathrm{ridge}}\|\boldsymbol\theta\|_2^2, 
\qquad \lambda_{\mathrm{ridge}}=0.1.
\end{equation}
式~\eqref{eq:q1-matrix-v9} 给出可检查的预测函数, 不赋予冗余系数独立的领域贡献解释. 在 $p_a>0$ 的单纯形内部, 对满足 $\one^\top\boldsymbol v=0$ 的可行切向量, 锚点函数差的方向导数为
\begin{equation}\label{eq:q1-tangent-v9}
D\DeltaJ(\pvec)[\boldsymbol v]
=\sum_{a=1}^{17}\frac{[\boldsymbol a+2\mathbf B\boldsymbol\phi]_a}{2\sqrt{p_a}}v_a.
\end{equation}
%该式由 $d\sqrt{p_a}/dp_a=1/(2\sqrt{p_a})$ 和二次型的对称性推出. 它只沿单纯形切向量有配比意义；当某个 $p_a=0$ 时分母可发散, 本文改用有限替换路径而不外推边界导数. 

模型含 $170$ 个多项式特征, 岭惩罚为 $\lambda_{\mathrm{ridge}}=0.1$. 式~\eqref{eq:sqrtquad} 的结构可直接检查；由于 $\sum_i p_i=1$ 带来系数冗余, 系数值并非独立的领域贡献. 函数在闭单纯形上连续, 但含根号的项在零比例边界一般不可微. 下文均以 $\widehat J$ 表示这一平方根坐标下的综合响应面, 不另设重复符号. 

候选集共 $16$ 个配置：原有 $13$ 种模型, 加上式~\eqref{eq:sqrtquad} 的同一平方根全二次结构在 $\lambda_{\mathrm{ridge}}=0.1, 1, 10$ 下的三个岭回归配置. 最终选型在全部 $512$ 组 A4/A5 数据上进行：五折交叉验证重复三次, 共形成 $15$ 个验证折, 以逐折十三域宏平均 RMSE 的平均值排序, 并在全部数据上重拟合胜出模型. 另在同一数据上做一次外层五折、内层三折的嵌套验证：每个外层训练集只用其内层三折选择候选, 再在该外层留出折上评价选型流程；外层结果不参与最终选型. 表~\ref{tab:model} 给出主要候选的重复交叉验证成绩及历史 1M 参照结果. 
\begin{table}[H]\centering\small
\caption{A4/A5 交叉验证与 A6/A7 历史参照表}\label{tab:model}
\begin{tabular}{lrrr}\toprule
模型 & CV macro RMSE & 1M macro RMSE & 1M 目标 Spearman\\\midrule
平方根二次岭, $\lambda_{\mathrm{ridge}}=0.1$ & \textbf{0.04453} & \textbf{0.03781} & \textbf{0.9602}\\
既有平方根 RBF 核岭 & 0.04583 & 0.03822 & 0.9528\\
固定极端随机树基线 & 0.08096 & 0.07344 & 0.9274\\
旧正交对比坐标二次岭 & 0.08497 & 0.07654 & 0.7976\\
常数均值 & 0.14725 & -- & --\\\bottomrule
\end{tabular}
\end{table}
表~\ref{tab:model} 汇总 A4/A5 的重复交叉验证结果以及 A6/A7 的历史参照结果. RMSE 统一按 A5 域均值归一化,  1M 目标的 Spearman 相关系数采用同一固定 A5 标尺计算. A6/A7 标签未参与模型拟合,  在本文中作为历史参照结果单独报告.

平方根二次面相对既有平方根 RBF 的 CV 均值差为 $0.001304$, 即约 $2.85\%$, 小于前者的折间标准差 $0.001954$. 两者同折比较时, 前者在 $15$ 折中的 $12$ 折较低. 两种模型共享训练行和重复验证行, 不能把 $15$ 折误当独立样本, 重复交叉验证的重叠校正给出双侧 $p\approx0.179$, 且这一事后比较没有校正 $16$ 种候选的选型. 嵌套验证的五个外层折均在内层选中平方根 RBF. 故本文按完整训练集 CV 排序规则采用平方根二次面, 同时不宣称它在总体上优于平方根 RBF. 相对固定树和正交对比二次面的近半误差降幅主要属于平方根坐标族, 而不是某一种具体二次结构独有的优势. 
\begin{figure}[H]\centering
\includegraphics[width=0.95\linewidth]{q1_cv_model_comparison.pdf}
\caption{第一问重复交叉验证的模型表现与同折误差对照图}\label{fig:q1-cv-contrast-v11}
\end{figure}
图~\ref{fig:q1-cv-contrast-v11} 展示第一问重复交叉验证下的模型对照结果,  从模型族整体表现和同族模型细粒度差异两个层面验证模型选择的合理性. 左图给出四个代表候选模型的宏平均 RMSE 及折间标准差,  可以直观看出不同模型族之间较为清晰的性能差距,  体现平方根坐标建模在误差控制上的良好表现. 右图进一步比较相同验证折上平方根二次岭与平方根 RBF 的 RMSE 差值,  负值表示平方根二次岭取得更低误差；在 15 个验证折中,  平方根二次岭有 12 折表现更优,  显示出较好的跨折稳定性. 图中误差线表示重复验证折上的标准差,  用于刻画不同数据划分下模型误差的波动程度. 这种“模型族整体比较 + 同折直接对照”的验证方式既能够识别主要性能提升来源,  又能够细致呈现候选模型之间的差异,  使最终模型选择具有更充分、直观且可复核的数值依据.
\begin{figure}[H]\centering
\includegraphics[width=0.68\linewidth]{q1_mixture_1m_validation.pdf}
\caption{平方根二次岭在 A6/A7 的 256 组候选配方上的相对 Loss 预测. }\label{fig:pred}
\end{figure}
图~\ref{fig:pred} 对照所选响应面在 1M 检验候选上的预测与观测. 模型仅在给定配方附近得到验证, 故首选配方从已观测候选集合选出, 不把式~\eqref{eq:sqrtquad} 连续优化至训练支持以外的单纯形角点. 
上述定义已经给出第一问的完整数学接口：文档级信号经统一方向和组权得到质量代理, 域级聚合只在可观测域上进行；配比模型在平方根坐标中拟合, 并以已观测候选集上的验证误差决定候选方案. 后续主计算的质量域映射限制在表~\ref{tab:qmap} 可核验的六个域内；前述扩展域映射仅是条件敏感性情景. 



\subsection{模型验证、跨尺度预测与第一问结论}

建立响应面之后, 第一问真正需要回答的是两个更严格的问题：模型能否在未参与拟合的配方上保持排序能力, 以及这种能力能否跨参数尺度迁移. 前者决定响应面是否可用于候选选择, 后者决定其结果能否直接进入后续大模型情景. 本文因此将“同尺度验证”和“跨尺度压力测试”分开报告. 

对规模 $s$ 的候选集合 $\mathcal T_s$, 令 $J_s$ 按该表自身十三域均值标准化真实 Loss, 仅用于评价排序；另令 $\widetilde J_s$ 按 A5 的 $\bar{\mathcal L}_j$ 标准化同一 Loss. 模型选择均不读取两种检验目标. 给定其中任一目标 $K_s$, 候选集后悔值定义为
\begin{equation}\label{eq:regret}
R_s(K)=
\frac{K_s(\widehat{\pvec}_s)-\min_{\pvec\in\mathcal T_s}K_s(\pvec)}
{\min_{\pvec\in\mathcal T_s}K_s(\pvec)}.
\end{equation}
后悔值 $R_s(K)$ 是一个决策指标：它不要求模型准确预测每个候选的绝对 Loss, 而只衡量被冻结的选择相对该候选表真实最优值损失了多少. $R_s=0$ 表示在给定候选集与给定目标口径下选择正确；它不意味着连续单纯形上的全局最优. 不同标尺可改变真最优行与后悔值；预测首选行只由已冻结模型决定. 60M 和 1B 仅报告排序及选择质量, 不把 1M 训练得到的数值标尺当作跨规模绝对 Loss 预测. 
\Needspace{13\baselineskip}
\subsubsection{响应误差到配比选择的决策保证}

式~\eqref{eq:regret} 用已观测候选表评价选择结果. 要在选择之前控制可能的损失, 还需把预测误差转换为决策误差. 以下命题固定候选集合和目标标尺, 不跨不同尺度或归一化方式混用误差. 

\begin{proposition}[配比选择的误差与后悔界]\label{prop:q1-decision-regret-v37}
设 $\mathcal P$ 为非空有限候选集合, $J$ 与 $\widehat J$ 分别为同一口径的真实目标和预测目标. 若
\[
\max_{\pvec\in\mathcal P}|J(\pvec)-\widehat J(\pvec)|\le\varepsilon, 
\qquad
\widehat J(\widehat{\pvec})\le\min_{\pvec\in\mathcal P}\widehat J(\pvec)+\varepsilon_{\mathrm{sel}}, 
\]
其中 $\varepsilon, \varepsilon_{\mathrm{sel}}\ge0$ 分别为候选集合上的一致预测误差界和候选选择的优化误差界, 则
\begin{equation}\label{eq:q1-decision-regret-v37}
0\le J(\widehat{\pvec})-J_{\min}\le2\varepsilon+\varepsilon_{\mathrm{sel}}, 
\quad J_{\min}:=\min_{\pvec\in\mathcal P}J(\pvec).
\end{equation}
若 $J_{\min}>0$, 相对后悔值进一步不超过 $(2\varepsilon+\varepsilon_{\mathrm{sel}})/J_{\min}$. 
\end{proposition}
\begin{proof}
取 $\pvec_*\in\argmin_{\mathcal P}J$, 依次使用一致误差界和预测近优性, 得
\[
J(\widehat{\pvec})\le\widehat J(\widehat{\pvec})+\varepsilon
\le\widehat J(\pvec_*)+\varepsilon_{\mathrm{sel}}+\varepsilon
\le J(\pvec_*)+2\varepsilon+\varepsilon_{\mathrm{sel}}.
\]
非负性来自真实最优值的定义；正分母下相除得到相对界. 
\end{proof}

该界将响应面拟合和候选选择连接起来, 符合决策导向预测的基本思想\cite{elmachtoub_spo}. 如果一致误差界以概率至少 $1-\zeta$ 成立, 则后悔界在同一事件上成立；普通交叉验证 RMSE、逐点区间或单次重抽样百分位均不能直接代入为 $\varepsilon$. 

\textbf{模型集合内的相对基线改进:}
设 $\mathcal C$ 为非空的候选响应模型集合, 基线为 $\pvec^{(0)}$. 定义
\begin{equation}\label{eq:q1-safe-improvement-v37}
U(\pvec):=\sup_{\theta\in\mathcal C}
\{J_\theta(\pvec)-J_\theta(\pvec^{(0)})\}.
\end{equation}
若 $U(\pvec)<0$, 则该配比在 $\mathcal C$ 中每个模型下都严格优于基线；这是上确界定义的直接结果. 
平方根二次面、核模型及其联合参数扰动可形成明确声明的模型集合, 但有限情景不自动覆盖真实机制. 因此式~\eqref{eq:q1-safe-improvement-v37} 给出模型集合内的安全选择准则, 未改变下文已经冻结的候选行与验证结果. 跨问传递仍采用完整响应差 $\widehat J(\pvec)-\widehat J(\pvec^{(0)})$, 而非将局部探针的线性化外推到任意远配比. 


\begin{table}[H]\centering\scriptsize
\setlength{\tabcolsep}{2.5pt}
\caption{冻结平方根二次面在历史候选表上的排序表}\label{tab:scale}
\begin{tabular}{lrrrrrrrl}\toprule
& & \multicolumn{3}{c}{按本表均值 $J_s$} & \multicolumn{3}{c}{按 A5 均值 $\widetilde J_s$} &\\
\cmidrule(lr){3-5}\cmidrule(lr){6-8}
规模 & 候选数 & Spearman & 首选/真最优行 & 后悔值 & Spearman & 真最优行 & 后悔值 & 数据属性\\\midrule
A6/A7 1M & 256 & 0.960 & 239/239 & 0.000\% & 0.960 & 239 & 0.000\% & 实测\\
A8/A9 60M & 256 & 0.931 & 239/239 & 0.000\% & 0.934 & 219 & 0.445\% & 实测\\
A10/A11 1B & 64 & 0.701 & 13/46 & 0.512\% & 0.714 & 13 & 0.000\% & 实测\\
A12/A13 10B & 63 & $-0.681$ & 16/45 & 21.406\% & $-0.673$ & 45 & 21.265\% & 估算\\
A14/A15 70B & 63 & $-0.789$ & 16/32 & 31.768\% & $-0.786$ & 32 & 31.430\% & 估算\\\bottomrule
\end{tabular}
\end{table}

表~\ref{tab:scale} 汇总冻结后的平方根二次面在历史候选表上的排序结果. 前组按照对应检验表自身均值 $J\_s$ 计算,  后组按照 A5 均值 $\tilde{J}\_s$ 计算；10B 和 70B 用于检验模型在更大尺度条件下的估算表现. 各行行号均按对应候选表内部顺序自 1 起计,  便于与原始候选记录逐项对应.

\textbf{证据结论：} 表~\ref{tab:scale} 展示了清晰的跨尺度证据层次. 在 1M 候选集上,  两种标尺均稳定选中第 239 行,  Spearman 相关系数达到 $0.960$,  表明模型对该尺度下的最优配比具有较好的识别能力；在 60M 上,  排序相关仍保持在约 $0.93$,  显示出较强的排序一致性,  尽管最优行会随评价标尺变化；在 1B 上,  排序相关约为 $0.70$,  仍保留一定的候选排序信息. 当尺度进一步扩展至 10B 和 70B 时,  Spearman 相关分别约为 $-0.68$ 和 $-0.79$,  后悔值超过 $21\%$ 和 $31\%$,  从而清楚标示出当前冻结模型的跨尺度适用范围.据此,  第 239 行记为 1M 实测候选集内的最优配比 $\pvec^\star$,  用于后续相应条件情景分析；后续主计算的训练配比锚点仍采用 A4 第 170 行的 $\pvec^{(0)}$. 10B 和 70B 的估算结果则作为跨尺度压力测试,  用于识别模型外推边界,  从而使后续分析能够明确区分实测支持范围与条件外推范围.


\textbf{外推表失效机制的诊断:} 表~\ref{tab:scale} 的负相关需要区分\textbf{跨域水平}与\textbf{配比形状}两条通道, 二者对结论的含义完全不同. 对每一档规模, 把逐域预测与逐域观测分别按该表自身的 A5 训练均值归一化, 再计算两种相关：\textbf{水平通道}取每行十三域归一化损失的均值作为该配方的汇总量, \textbf{形状通道}先减去该行自身的十三域均值, 再汇总全部 $n\times13$ 个中心化元素. 结果见表~\ref{tab:q1-extrapolation-diagnostic-v41}. 

\begin{table}[H]\centering\small
\caption{响应面的跨尺度迁移诊断}
\label{tab:q1-extrapolation-diagnostic-v41}
\begin{tabular}{lrrrrr}\toprule
规模 & 候选数 & 数据属性 & \multicolumn{2}{c}{水平通道} & 形状通道 \\
\cmidrule(lr){4-5}
 & & & Spearman & 最大绝对误差 & Spearman \\\midrule
1M & 256 & 实测 & $+0.9602$ & $0.0481$ & $+0.9606$ \\
60M & 256 & 实测 & $+0.9342$ & $0.3465$ & $+0.7166$ \\
1B & 64 & 实测 & $+0.7135$ & $0.5628$ & $+0.6764$ \\
10B & 63 & 估算 & $-0.6726$ & $0.8783$ & $+0.9779$ \\
70B & 63 & 估算 & $-0.7861$ & $0.9554$ & $+0.9774$ \\\bottomrule
\end{tabular}\end{table}
表~\ref{tab:q1-extrapolation-diagnostic-v41} 给出响应面的跨尺度迁移诊断结果. 其中,  水平通道采用十三个领域归一化损失的行均值,  用于统一比较不同尺度下响应面的整体变化与迁移表现.

水平与形状双通道诊断给出一个可检验的结构性结论：\textbf{配比形状通道在全部五档规模上保持强正相关}（$+0.96$、$+0.72$、$+0.68$、$+0.98$、$+0.98$）, 而\textbf{跨域水平通道在 10B/70B 反转为负}. 逐域分解进一步确认这一点（图~\ref{fig:q1-extrapolation-diagnostic-v41}）：在 1M 档, 13 个验证域的域内 Pearson 相关为 $0.929$--$0.993$；在 70B 档仍有 12 个域为正, 仅有 hackernews 为 $-0.046$, dm\_mathematics 与 arxiv 分别保持 $0.909$ 与 $0.745$. 

因此 A12--A15 的负相关意为跨域基线排序在大于 $1B$ 的尺度上被重新排列：模型仍能正确预测“同一配方内各域的相对高低”, 但其预测的域间绝对水平与真实域间水平不再同序, 而后者在十三域宏平均汇总时占据主导. 这一机制与更早出现的标尺敏感性方向一致：60M 档的真最优行随归一化方式由第 239 行变为第 219 行, 1B 档则由第 46 行变为第 13 行（表~\ref{tab:scale}）. 由此, A12--A15 的稳健性结论可精确表述为：\textbf{外推表支持“配比形状效应可跨尺度迁移”的相对比较, 不支持绝对 Loss 预测, 也不支持据其确定大尺度最优配比}；表中 10B/70B 的最优行因此不作报告. 
\begin{figure}[H]\centering
\includegraphics[width=0.95\linewidth]{q1_extrapolation_diagnostic.pdf}
\caption{响应面跨尺度迁移的相关性诊断图}\label{fig:q1-extrapolation-diagnostic-v41}
\end{figure}
图~\ref{fig:q1-extrapolation-diagnostic-v41} 从整体结构与分域表现两个层面展示响应面的跨尺度迁移特征. 左图比较水平通道与配比形状通道在不同规模下的 Spearman 相关,  灰色区域对应估算外推表；右图给出各验证域在五档规模上的域内 Pearson 相关,  用于进一步刻画不同领域随规模变化的局部一致性.
\begin{table}[H]\centering\small
\caption{1M 候选集第 239 组的十七域归一化配比表}\label{tab:pstar}
\begin{tabular}{lr@{\qquad}lr}\toprule
训练域 & 比例 & 训练域 & 比例\\\midrule
arxiv & 0.096 & freelaw & 0.084\\
nih\_exporter & 0.000 & pubmed\_central & 0.019\\
wikipedia\_en & 0.069 & dm\_mathematics & 0.028\\
github & 0.030 & philpapers & 0.000\\
stackexchange & 0.127 & enron\_emails & 0.000\\
gutenberg\_pg\_19 & 0.002 & pile\_cc & 0.034\\
ubuntu\_irc & 0.132 & europarl & 0.025\\
hackernews & 0.000 & pubmed\_abstracts & 0.349\\
uspto\_backgrounds & 0.005 & & \\\bottomrule
\end{tabular}
\end{table}
表~\ref{tab:pstar} 给出 1M 检验候选集第 239 组在十七个领域上的配比结果. 各域比例已按行内总和统一归一化,  因而可直接比较不同领域在该候选配比中的相对占比.
\subsubsection{参考方法与对照检验}

本节用参考文献提出的四种检验方法复核前面的选择, 目的不是替换已给出的模型, 而是确认
哪些建模简化在本题数据上确实成立. 对照均沿用同一历史候选集合并注明目标口径. 

\textbf{多目标稳健性与最坏域代价（DoReMi \cite{doremi}）}
等权平均相对损失隐含了一个选择：十三个域同等重要. 若改用逐域超额损失或最坏域口径, 
候选排序会改变. 表~\ref{tab:q1-utility} 给出同一 256 组候选在四种目标下的最优行与统计量. 

\begin{table}[H]\centering\small
\setlength{\tabcolsep}{4pt}
\caption{1M 检验候选集上的多目标对照表}
\label{tab:q1-utility}
\begin{tabular}{lrrrr}
\toprule
目标 & 最优行 & 该行均值 & 该行最坏域 & 该行域间标准差\\
\midrule
等权平均相对损失 $\frac1{13}\sum_jY_j$（本文主口径） & 239 & 0.8935 & 1.0263 & 0.1238\\
最坏域 $\max_jY_j$ & 199 & 0.8938 & 1.0093 & 0.1283\\
域间标准差 $\mathrm{sd}_jY_j$ & 76 & 0.9450 & 1.0874 & 0.0680\\
平均超额损失（原始尺度） & 239 & 0.8935 & 1.0263 & 0.1238\\
平均超额损失（按各域空间归一） & 219 & 0.9088 & 1.0211 & 0.1049\\
\bottomrule
\end{tabular}
\end{table}

\textbf{多目标对照的结论}
在本文固定的十三域等权目标下, 平方根二次岭与平方根 RBF 都从同一 1M 实测候选池中选出第 239 行, 且该行也是观测最优行, 候选集内的实际后悔为零（表~\ref{tab:scale}）. 两种响应面在这一目标和候选范围内给出一致的首选. 表~\ref{tab:q1-utility} 进一步量化目标偏好带来的选择：若更重视最坏域, 可改选第 199 行, 使最坏域损失下降 $1.66\%$, 而等权平均损失仅上升 $0.038\%$；按各域可改进空间归一的超额损失则选第 219 行. 因此第 239 行是等权目标的首选, 第 199 行是最坏域优先时的备选. 与 RegMix、DoReMi 等方法及 B1、B5、C6 的补充数值对照见附录~\ref{app:q1-method-crosschecks}. 

\subsubsection{锚点效应与解释边界}

单纯形上增加某个领域的比例必须减少其余领域的合计比例. 取 A4 中最接近训练配方均值的实际配方 $\pvec^{(0)}$ 为锚点, 并令局部探针幅度 $\delta_{\mathrm{probe}}=0.05$. 这里的扰动是将原配方的 $5\%$ 权重向顶点 $\boldsymbol e_i$ 替换, 领域 $i$ 的实际份额增量为 $\delta_{\mathrm{probe}}(1-p_i^{(0)})$. 定义
\begin{subequations}
    \begin{align}\label{eq:effects}
\pvec^{(i)}&=(1-\delta_{\mathrm{probe}})\pvec^{(0)}+\delta_{\mathrm{probe}}\boldsymbol e_i, \\
\Delta_i&=\widehat J(\pvec^{(i)})
          -\widehat J(\pvec^{(0)}), \\
\Gamma_{ik}&=
\widehat J\!\left((1-2\delta_{\mathrm{probe}})\pvec^{(0)}+
               \delta_{\mathrm{probe}}\boldsymbol e_i+\delta_{\mathrm{probe}}\boldsymbol e_k\right)
-\widehat J(\pvec^{(i)})
-\widehat J(\pvec^{(k)})
+\widehat J(\pvec^{(0)}).
\end{align}
\end{subequations}
负的 $\Delta_i$ 表示该锚点与扰动强度下的模型预测 Loss 降低. 负的 $\Gamma_{ik}$ 表示联合变化低于两项单域变化之和. 双域探针值按定义满足 $\Delta_i+\Delta_k+\Gamma_{ik}$.

\begin{table}[htbp]\centering\small
\caption{平方根二次响应面在锚点附近的部分有限差分.}\label{tab:effects}
\begin{tabular}{lrr}\toprule
方向或组合 & 点估计 & 重抽样 $95\%$ 分位区间\\\midrule
arxiv 单域 & $-0.0131$ & $[-0.0262, -0.0063]$\\
hackernews 单域 & $-0.0048$ & $[-0.0118, \ 0.0063]$\\
philpapers 单域 & $-0.0040$ & $[-0.0278, \ 0.0237]$\\
nih\_exporter 与 ubuntu\_irc 交互 & $-0.0086$ & $[-0.0214, \ 0.0021]$\\\bottomrule
\end{tabular}
\end{table}

在 136 个双域探针中, 74 个 $\Gamma_{ik}<0$、62 个大于零. 数值最大的负交互本身仍有跨零区间. 与同样采用平方根坐标的既有 RBF 模型相比, 正确锚点下 17 个单域方向仅 12 个同号, 136 个双域交互有 112 个同号. %故即便主响应面是可写出函数式的参数模型, 也不能将全部局部符号作为稳定的领域效应. 第二问采用响应面在具体配方处相对锚点的直接差值；探针差分只用于解释和情景计算. 

\begin{table}[H]\centering\scriptsize
\caption{锚点处 $5\%$ 单域再分配的相对 Loss 效应}\label{tab:q1-v39-single}
\begin{tabular}{lccc}\toprule 领域 & 平方根二次面 $\Delta_i$ & 条件 $95\%$ 区间 & 平方根 RBF $\Delta_i$\\\midrule
arxiv & -0.0131 & $[-0.0262,-0.0063]$ & -0.0192 \\
europarl & -0.0069 & $[-0.0323,0.0170]$ & -0.0059 \\
hackernews & -0.0048 & $[-0.0118,0.0063]$ & 0.0007 \\
gutenberg\_pg\_19 & -0.0041 & $[-0.0101,0.0014]$ & -0.0043 \\
philpapers & -0.0040 & $[-0.0278,0.0237]$ & 0.0082 \\
pubmed\_central & -0.0039 & $[-0.0087,-0.0006]$ & -0.0070 \\
ubuntu\_irc & -0.0029 & $[-0.0174,0.0142]$ & -0.0005 \\
pubmed\_abstracts & -0.0008 & $[-0.0048,0.0027]$ & -0.0021 \\
stackexchange & -0.0007 & $[-0.0036,0.0012]$ & -0.0010 \\
uspto\_backgrounds & -0.0006 & $[-0.0040,0.0020]$ & 0.0012 \\
github & 0.0006 & $[-0.0056,0.0038]$ & -0.0032 \\
freelaw & 0.0006 & $[-0.0017,0.0024]$ & 0.0011 \\
wikipedia\_en & 0.0010 & $[-0.0015,0.0033]$ & 0.0015 \\
pile\_cc & 0.0012 & $[-0.0028,0.0039]$ & -0.0002 \\
nih\_exporter & 0.0052 & $[-0.0170,0.0314]$ & 0.0025 \\
dm\_mathematics & 0.0116 & $[0.0089,0.0181]$ & 0.0075 \\
enron\_emails & 0.0321 & $[-0.0237,0.1035]$ & 0.0017 \\
\bottomrule\end{tabular}\end{table}

\begin{table}[htbp]\centering\scriptsize
\caption{两种平方根响应面的一致交互效应}\label{tab:q1-v39-pairs}
\begin{tabular}{lrrl}\toprule 双域探针 & 二次面 $\Gamma_{ik}$ & RBF $\Gamma_{ik}$ & 局部判据\\\midrule
ubuntu\_irc + pubmed\_abstracts & -0.00118 & -0.00109 & 互补 \\
ubuntu\_irc + uspto\_backgrounds & -0.00090 & -0.00067 & 互补 \\
pubmed\_central + dm\_mathematics & -0.00088 & -0.00055 & 互补 \\
nih\_exporter + gutenberg\_pg\_19 & 0.00547 & 0.00343 & 替代 \\
arxiv + stackexchange & 0.00183 & 0.00215 & 替代 \\
arxiv + pubmed\_central & 0.00163 & 0.00196 & 替代 \\
\bottomrule\end{tabular}\end{table}

在同一锚点与 $\delta_{\mathrm{probe}}=0.05$ 下, $\Gamma_{ik}<0$ 定义为联合增配的预测 Loss 低于单域效应之和, 称局部互补；$\Gamma_{ik}>0$ 为局部替代. 表~\ref{tab:q1-v39-single} 报告全部 17 个单域扰动. 两种平方根坐标模型有 112/136 个交互方向一致, 其中 15 个的两套条件重抽样区间同向且均不跨零；表~\ref{tab:q1-v39-pairs} 列出六个代表组合. 单域中 arxiv、pubmed\_central 的降 Loss 方向与 dm\_mathematics 的升 Loss 方向在这两种模型及各自条件区间下均保持一致. %这里的筛选不含模型选择不确定性或 136 次同时比较的校正, 也不表示任意配比的因果替代关系；全量 17 与 136 项数值随复核程序归档. 

%本模型的结论局限于给定损失口径和候选集合. 质量评分没有人工质量真值；事后原文诊断指标可能将代码、公式和链接误判为低质量信号. 配比模型未观测四个训练域的单域验证 Loss, 也不能把某个配比变量的统计效应解释为该域的因果收益. 10B、70B 的负相关提示估算口径下的外推风险, 应等待真实训练实验复核. 各域误差、冲突类型、重抽样区间和完整候选配比均作为附加证据保留. 

\subsubsection{第一问结论及可复用范围}
第一问得到文档质量基线和配比响应面两项可复用结果. 质量侧覆盖去重后的 261086 篇文档；六个可映射配比域在锚点中占 $31.94\%$, 其条件文档等权均值 $Q_0=0.607474$ 给出这部分数据的质量基线. 配比侧按预定的重复交叉验证规则选中平方根二次岭, 宏平均 RMSE 为 $0.044528$；在 1M 实测候选表上, 固定 A5 标尺的 RMSE 为 $0.037812$、Spearman 为 $0.960201$. 这些数值支持给定数据与目标下的配比排序, 质量基线则只对应已映射的六个域. 

第二问主计算接收三个明确对象：训练锚点 $\pvec^{(0)}$ 上的条件质量基线 $Q_0$、固定 A5 标尺上相对该锚点的响应面差值 $\DeltaJ(\pvec)$, 以及 A4 第 170 行的训练锚点配比. 1M 候选集内实测最优的第 239 行配比 $\pvec^\star$ 用于条件情景. 
% 质量评分尚缺人工盲评或下游训练真值；60M/1B 的标尺敏感性和 10B/70B 估算表的负相关限定了跨尺度应用. 因此后续计算沿用已说明的锚点与目标口径, 跨尺度效果仍须独立实测. 
\section{问题二：质量修正的广义标度律}\label{sec:q2}

第一问给出了两个性质不同的输出：$Q_0$ 是有限映射域上的条件质量代理，$\DeltaJ(\pvec)$ 是固定 Loss 标尺下的配比响应. 二者都不能直接解释为训练收益. 故第二问的任务是建立一个中间层，把参数量、Token、质量与配比整合到同一 Loss 模型中，并明确哪些系数来自真实轨迹、哪些来自半合成质量扰动. 

\subsection{有效数据量与广义标度模型}

令 $N$ 为模型参数量，$D$ 为训练 Token 数，\(Q\in(0,1]\) 为质量指数，\(\pvec=(p_1,\ldots,p_{17})\) 为 17 个训练领域的比例向量，满足
\begin{equation}
p_i\ge 0,\quad \sum_{i=1}^{17}p_i=1.
\end{equation}
传统标度律将验证交叉熵写为 \(\mathcal L=E+AN^{-\alpha}+BD^{-\beta}\)，可描述参数规模与数据规模的幂律关系\cite{kaplan,hoffmann,bahri}. 本文研究如何将第一问得到的质量与配比信息纳入该关系. 文献关于数据配比的研究表明，不同领域对验证损失的影响并不相同，且可能出现领域交互\cite{regmix,doremi}；因此本文把 $Q$ 与 \(\pvec\) 分别作为数据质量和数据构成的输入. 

表~\ref{tab:data} 列出来源与用途: 附件 B 的主拟合数据为 B1 Pythia 训练日志，共 1176 行、8 个参数规模. B2 为 1029 行 Cerebras 半合成轨迹，B3 为 8 条各 500 点的 Pythia 插值轨迹，B4 为 57 个跨模型族收敛点，B5 为 44 个文献标度点. B6、B7、B8 为依据真实数据校准的半合成数据, 提供质量信号. B9 是模型元数据，B10 是大模型损失估算值, B11 含 18 条模型族元数据，B12 含 1386 条 Pythia 检查点索引，用于核对来源与检查点覆盖，不含验证损失，故不作为拟合观测. Pythia 公开套件覆盖 70M 至 12B 参数并发布训练检查点，B1 的观测范围与该训练体系相符. %\cite{pythia}. 

\begin{table}[htbp]
\centering\small
\caption{附件 B 中与第二问相关的数据及其证据角色}
\label{tab:data}
\begin{tabular}{P{1.2cm}P{3.3cm}P{2.0cm}P{7.0cm}}
\toprule
编号 & 内容 & 数据性质 & 本文用途\\
\midrule
B1 & Pythia 训练日志，1176 行 & 真实轨迹 & 经典 $N$--$D$ 主拟合及分组验证\\
B2 & Cerebras 轨迹，1029 行 & 半合成 & 跨族排序与损失标定检验\\
B3 & 8 条 Pythia 插值轨迹 & 插值数据 & 轨迹插值一致性核验，不作独立观测\\
B4、B5 & 跨族收敛点、文献数据 & 真实整理数据 & 跨族及文献排序检验\\
B6、B7 & 360、450 个 $N$--$D$--$Q$ 点 & 半合成 & 质量响应估计与扩展质量等级检验\\
B8 & 1704 个含外推点的 $N$--$D$--$Q$ 点 & 半合成 & 不同校准机制下的压力测试\\
B9、B10 & 大模型元数据、估算损失 & 元数据、估算 & 讨论 100B--10000B 参数范围的外推边界\\
B11、B12 & 模型族元数据、Pythia 检查点索引 & 辅助元数据 & 来源与检查点覆盖说明；不含损失观测，不参与拟合\\
\bottomrule
\end{tabular}
\end{table}

\subsubsection{第一问输出与识别边界}

第一问已在质量评估中给出域级 \(\bar Q_b\)，并在 17 域配比实验上建立响应面. 可与配比域直接或近似对应的质量域只有 6 个. 第一问扩展数据给出的 \(\bar Q_b\) 见表~\ref{tab:qmap}. 其余 11 个配比域没有对应的质量观测，因此不对缺失域填入臆测质量值. 

\begin{table}[H]
\centering\small
\caption{第一问质量评分与配比域的直接或近似映射}
\label{tab:qmap}
\begin{tabular}{lrrl}
\toprule
配比域 & \(\bar Q_b\) & 映射关系 & 质量数据\\
\midrule
arxiv & 0.8209 & 直接 & A2 扩展集\\
github & 0.5041 & 直接 & A3 扩展集\\
stackexchange & 0.6330 & 直接 & A1\\
wikipedia\_en & 0.5718 & 近似 & A1 的 wikipedia\\
gutenberg\_pg\_19 & 0.7011 & 近似 & A1 的 book\\
pile\_cc & 0.6612 & 近似 & A1 的 commoncrawl\\
\bottomrule
\end{tabular}
\end{table}

令 $\mathcal M$ 为表~\ref{tab:qmap} 的六个配比域，$\pi(a)$ 为配比域 $a$ 对应的质量域. 锚点已映射份额与文档等权条件均值定义为
\begin{equation}\label{eq:q0-mapping-v9}
m_0=\sum_{a\in\mathcal M}p_a^{(0)}=0.319361,\quad Q_0=\frac{\sum_{a\in\mathcal M}p_a^{(0)}\bar Q_{\pi(a)}}{m_0}=0.607474.
\end{equation}
这里 $m_0$ 是定义~\ref{def:background-mapped-quality} 中 $m_{\mathcal M}$ 在锚点 $\pvec^{(0)}$ 的取值. $Q_0$ 不是完整十七域配方的实测质量. 由命题~\ref{prop:background-quality-bound}，若其余 $68.06\%$ 份额的质量代理只限定在 $[0,1]$，则完整锚点质量在这些约束下的取值范围为
\begin{equation}\label{eq:q2-full-quality-bound-v15}
Q_{\mathrm{all}}^{(0)}\in[m_0Q_0,\,m_0Q_0+1-m_0]=[0.1940,0.8746].
\end{equation}
第三问采用 $Q_0$ 作质量处理起点时，额外假设这个六域条件均值可代表所处理语料的质量标尺. 式~\eqref{eq:q2-full-quality-bound-v15} 不能被误作其统计置信区间. 

结合表~\ref{tab:pstar} 的候选配比和表~\ref{tab:effects} 的响应方向，第一问的平方根二次模型在历史 1M 候选表上得到目标 Spearman $\rho=0.9602$（固定 A5 标尺）、宏平均 RMSE $0.0378$，并选中第 239 行，记为 $\pvec^\star$. 在最接近训练配比均值的实际锚点 $\pvec^{(0)}$ 上，第一问输出有限差分与完整响应面探针值. 第二问仅使用第一问已经定义的质量基线、锚点和响应面差值. 预测值与局部探针均来自同一响应面，因此它们只提供同一模型下的两种解释，不能作为跨模型稳健性证据. 

对 $\pvec^{(0)}$，六个已映射域的配比覆盖率为 $31.94\%$，其条件文档等权质量均值为 $0.6075$（按原文词数加权为 $0.6617$）；对 $\pvec^\star$，覆盖率为 $35.80\%$，条件均值为 $0.6638$（词数加权为 $0.6995$）. 这些只代表已映射份额. 对训练锚点 $\pvec^{(0)}$. 第二问因而把配比响应限定为锚点邻域的条件情景.

\subsubsection{质量响应函数、有效数据状态与来源分层}

\begin{definition}[质量响应与有效数据量]\label{def:q2-quality-response}
令 $n:=N/10^9$、$d:=D/10^9$，来源 $s$ 的参数满足 $A_s,B_s,\alpha_s,\beta_s>0$. 设 $g_s$ 在包含 $[\ln0.1,0]$ 的开区间上连续可微，且 $g_s(0)=0$. 定义
\begin{subequations}\label{eq:q2-generalized-v26}
\begin{equation}
    \mathcal G_s(Q):=\exp\{-g_s(\ln Q)/\beta_s\},\quad \psi(\pvec):=\DeltaJ(\pvec)
\end{equation}
\begin{equation}
    d_{\mathrm{eff},s}:=d\mathcal G_s(Q)e^{-\lambda_{\mathrm{mix}}\psi(\pvec)},
\end{equation}
\begin{equation}
    \mathcal L_s(n,d,Q,\pvec):=E_s+A_sn^{-\alpha_s}+B_sd^{-\beta_s} e^{g_s(\ln Q)+\beta_s\lambda_{\mathrm{mix}}\psi(\pvec)}.
\end{equation}
\end{subequations}

由 $g_s(0)=0$，则在 $Q=1$ 且 $\pvec=\pvec^{(0)}$（或 $\lambda_{\mathrm{mix}}\psi(\pvec)=0$）时，式~\eqref{eq:q2-generalized-v26} 退化为经典 $N$--$D$ 标度律；任意新配方即使 $Q=1$，配比项通常仍存在. 
质量的局部斜率与有效 Token 弹性分别为
\begin{equation}\label{eq:q2-local-quality}
a_s(Q):=-g_s'(\ln Q),\quad
\gamma_s(Q):=\frac{a_s(Q)}{\beta_s}
=\frac{\dd\ln\mathcal G_s(Q)}{\dd\ln Q}.
\end{equation}
当 $a_s\ge0$ 时，固定其余输入，提高 $Q$ 不增加模型 Loss. $\mathcal G_s$ 是给定来源标尺上的折算函数，并非由真实信息量直接观测得到. 
\end{definition}
\Needspace{21\baselineskip}
\begin{theorem}[有效数据状态的有限等价及其局部切线]\label{thm:q2-log-marginal-v14}
固定来源，设 $E\in\mathbb R$, $A,B,\alpha,\beta>0$, $\lambda_{\mathrm{mix}}\in\mathbb R$，且 $g\in C^1$. 令 $\mathcal Z:=\{\psi(\pvec):\pvec\in\Delta^{16}\}$. 对 $z\in\mathcal Z$ 定义 $d_{\rm eff}:=d\mathcal G(Q)e^{-\lambda_{\mathrm{mix}}z}$. 在同一 $n$ 下，两个可达状态 $(d_1,Q_1,z_1)$, $(d_2,Q_2,z_2)$ 的 Loss 相等，当且仅当
\begin{equation}\label{eq:q2-finite-equivalence}
\ln\frac{d_2}{d_1}
+\int_{Q_1}^{Q_2}\gamma(q)\,\frac{\dd q}{q}
-\lambda_{\mathrm{mix}}(z_2-z_1)=0.
\end{equation}
在局部把 $z$ 延拓为实数坐标，记 $\mathcal L_N:=An^{-\alpha}$, $\mathcal L_D:=Bd_{\rm eff}^{-\beta}$，则
\begin{equation}\label{eq:q2-log-differential}
\dd\mathcal L=-\alpha\mathcal L_N\dd\ln n
-\beta\mathcal L_D\{\dd\ln d+\gamma(Q)\dd\ln Q
-\lambda_{\mathrm{mix}}\dd z\}.
\end{equation}
真实配比变动还必须满足单纯形可行方向，且 $\dd z=D\psi(\pvec)[\dd\pvec]$；独立 $z$ 坐标的切线不保证任意变化都可由配比实现. 
\end{theorem}
\begin{proof}
因 $B>0,\beta>0$，映射 $v\mapsto Bv^{-\beta}$ 在 $v>0$ 上严格递减，故同一 $n$ 下 Loss 相等当且仅当有效数据量相等. 两状态取对数，并以 $\ln\mathcal G(Q_2)-\ln\mathcal G(Q_1)=\int_{Q_1}^{Q_2}\gamma(q)\dd q/q$ 代入，即得有限关系. 对 $An^{-\alpha}$ 和 $Bd_{\rm eff}^{-\beta}$ 分别使用链式法则，再用 $\dd\ln d_{\rm eff}=\dd\ln d+\gamma(Q)\dd\ln Q-\lambda_{\mathrm{mix}}\dd z$ 得全微分. 配比约束由背景中的可行方向定义给出. 
\end{proof}
固定参数量和配比时，式~\eqref{eq:q2-finite-equivalence} 给出有限质量提升的等效数据量；下文的形状约束与识别诊断决定该关系在哪些质量区间和来源内可以使用. 

\textbf{总响应与机制的区别:}
固定 $n,d,\pvec$ 时，式~\eqref{eq:q2-generalized-v26} 的质量方向导数为
\begin{equation}\label{eq:v40-quality-derivative}
\frac{\partial\mathcal L_s}{\partial Q}
=-\frac{a_s(Q)}{Q}\,B_sd^{-\beta_s}
e^{g_s(\ln Q)+\beta_s\lambda_{\mathrm{mix}}\psi(\pvec)}.
\end{equation}
故“质量提高使预测 Loss 非增”严格依赖拟合模型的 $a_s(Q)\ge0$ 形状约束. 这里的 $Q$ 汇总了若干文本信号，并未分别操纵或识别重复、噪声、有效信息密度等通道. 式~\eqref{eq:v40-quality-derivative} 是代理质量的总响应. 

经典质量幂律对应 $g_s(u)=-\delta_su$，因而 $\gamma_s=\delta_s/\beta_s$ 为常数. 质量效应若随水平改变，直接使用一个固定指数就会把局部替代率错误地用于有限质量变化. 本文比较嵌套的两个模型：
\begin{equation}\label{eq:q2-power-curved}
\mathrm P:\ g(\ln Q)=-\delta\ln Q,\quad
\mathrm C:\ g(\ln Q)=-\delta\ln Q+\kappa_Q(\ln Q)^2,\quad \ln Q\in[\ln0.1,0].
\end{equation}
以下无来源下标的曲率参数仅指 B6；B1 的真实 Pythia 轨迹仍独立拟合 $E+An^{-\alpha}+Bd^{-\beta}$. 

\begin{proposition}[曲率响应的全区间形状约束]\label{prop:q2-shape}
置 $u_{\min}:=\ln0.1$，用 $(\delta,\delta_{\min})$ 参数化曲率：
\begin{equation}\label{eq:q2-shape}
\delta_{\min}:=\delta-2\kappa_Q u_{\min},\quad
\kappa_Q=\frac{\delta-\delta_{\min}}{2u_{\min}},\quad
0\le\delta,\delta_{\min}\le5\beta.
\end{equation}
则所有 $Q\in[0.1,1]$ 均满足 $0\le a(Q)\le5\beta$，从而 $0\le\gamma(Q)\le5$，且 $\mathcal L$ 关于 $Q$ 非增. $\delta=\delta_{\min}$ 时模型 C 精确退化为模型 P. 
\end{proposition}
\begin{proof}
$a(Q)=\delta-2\kappa_Q\ln Q$ 是 $\ln Q$ 的仿射函数，其两个端点值为 $\delta_{\min}$ 与 $\delta$. 区间内每个值都是端点值的凸组合，故继承上下界. 又 $\partial_Q\mathcal L=-a(Q)\mathcal L_D/Q$，其中 $\mathcal L_D>0$ 是数据项；最后由 $u_{\min}\ne0$ 得退化结论. 
\end{proof}

配比传递采用同一响应面上的精确差值. 若第一问的二次响应为 $\widehat J(\pvec)=b+\boldsymbol a^\top\boldsymbol\phi+\boldsymbol\phi^\top\mathbf B\boldsymbol\phi$，其中 $\mathbf B$ 对称，则
\begin{equation}\label{eq:q2-exact-mixture}
\psi(\pvec)=\boldsymbol a^\top(\boldsymbol\phi-\boldsymbol\phi_0)
+(\boldsymbol\phi-\boldsymbol\phi_0)^\top\mathbf B(\boldsymbol\phi+\boldsymbol\phi_0).
\end{equation}
直接展开两个二次式即可验证该恒等式, 把 $\widehat J(\pvec)$ 本身作为迁移项，会使锚点产生无依据的额外偏移. 

\begin{proposition}[锚点数据不能识别配比传递系数]\label{prop:q2-lambda-unidentified-v15}
若某来源的全部拟合观测均使用 $\pvec^{(0)}$，则该来源的 Loss 观测不能识别 $\lambda_{\mathrm{mix}}$. 
\end{proposition}
\begin{proof}
锚点满足 $\psi(\pvec^{(0)})=0$，故式~\eqref{eq:q2-generalized-v26} 的预测与 $\lambda_{\mathrm{mix}}$ 无关，关于该系数的设计方向恒为零. 任意两个系数值给出同一观测分布；这与背景中可识别性的定义不相容. 
\end{proof}
因此第三问主结果固定 $\pvec=\pvec^{(0)}$, 其他配比只给定 $\lambda_{\mathrm{mix}}$ 作条件情景. 质量代理与 B6 等级共用标尺、配比效应可以迁移，均是额外假设. 

\subsubsection{共同质量响应的秩结构与可检验性}

式~\eqref{eq:q2-generalized-v26} 不仅规定一条质量曲线，还要求不同资源单元共享其形状. 完整的资源单元与质量等级交叉设计使这一假设可以脱离资源基线单独检查. 

\begin{proposition}[中心化质量响应的秩一必要条件]\label{prop:q2-rank-one-v37}
固定来源和配比，令 $B_p:=B\exp\{\beta\lambda_{\rm mix}\psi\}$. 对 $R$ 个完整资源单元 $(n_r,d_r)$ 和 $m$ 个共同质量等级 $Q_j$，设模型均值矩阵为 $\mathbf M=(M_{rj})$. 若满足共同质量响应模型，则
\begin{equation}\label{eq:q2-rank-one-v37}
\mathbf M\mathbf P_m=\boldsymbol b\boldsymbol f^\top\mathbf P_m,
\quad \operatorname{rank}(\mathbf M\mathbf P_m)\le1,
\quad \mathbf P_m:=\mathbf I_m-\frac1m\one\one^\top,
\end{equation}
其中 $b_r=B_pd_r^{-\beta}$, $f_j=e^{g(\ln Q_j)}$. 
\end{proposition}
\begin{proof}
令 $a_r=E+An_r^{-\alpha}$，则 $\mathbf M=\boldsymbol a\one^\top+\boldsymbol b\boldsymbol f^\top$. 右乘 $\mathbf P_m$ 后第一项为零，第二项为两个向量的外积，因此秩至多为一. 
\end{proof}

对观测矩阵 $\mathbf Y$ 的中心化 $\mathbf Y\mathbf P_m$ 作奇异值分解，以 $s_1\ge s_2\ge\cdots$ 记奇异值，表~\ref{tab:q2-rank-diagnostic-v37} 给出结构诊断. 第一方向的能量占比为 $s_1^2/\sum_js_j^2$，残差 RMSE 按全部 $Rm$ 个中心化元素计算. 
\begin{table}[htbp]\centering\small
\caption{B7 与 B6 共享基础记录表}\label{tab:q2-rank-diagnostic-v37}
\begin{tabular}{lrrrr}\toprule
数据 & 完整布局 & 第一方向能量 & $s_2/s_1$ & 秩一残差 RMSE\\\midrule
B6 & $45\times8$ & $88.18\%$ & 0.1947 & 0.04388\\
B7 扩展表 & $45\times10$ & $87.07\%$ & 0.2065 & 0.04190\\\bottomrule
\end{tabular}\end{table}
%这些是观测变化的描述量，不是原始 Loss 的预测 $R^2$，也不是模型拒绝检验. 噪声、同一资源单元内的相关性和数值舍入都可产生非零高阶奇异值；在没有独立重复测量时，噪声与结构偏离不能完全分离. B7 也不构成独立验证. 即使秩一成立，仍须另查行载荷是否服从 $d_r^{-\beta}$，以及共同质量曲线的单调性. 

共同形状若获得数据支持，其具体函数不必局限于二次式. 例如在 $u=\ln Q$ 上取非负基函数 $M_j$，记
\begin{equation}\label{eq:q2-monotone-basis-v37}
v(u)=\sum_{j=1}^Ja_jM_j(u),\quad a_j\ge0,
\quad g(u)=\int_u^0v(s)\,\dd s,
\end{equation}
就有 $g(0)=0$, $g'(u)=-v(u)\le0$. 非负积分样条是成熟的单调回归构造\cite{ramsay_monotone}. 本文正式参数仍来自式~\eqref{eq:q2-power-curved} 的嵌套比较, 式~\eqref{eq:q2-monotone-basis-v37} 只说明后续理论可容纳的形状扩展. B6 仅有八个质量等级，新增自由度必须通过原有分组外验证；若转而引入质量与资源的交互项，则需重新检查第三问的可分离结构. 

\subsection{有限替代的扩展与坐标一致性}

当参数量和配比固定时，保持 Loss 不变所需的数据量变化为
\begin{equation}\label{eq:q2-data-substitution}
\frac{d_2}{d_1}=
\exp\!\left[-\int_{Q_1}^{Q_2}\gamma(q)\frac{\dd q}{q}\right]
=\exp\!\left[\frac{g(\ln Q_2)-g(\ln Q_1)}{\beta}\right].
\end{equation}
常数情形可简化为 $(Q_1/Q_2)^\gamma$. 局部有 $\dd\ln d=-\gamma(Q)\dd\ln Q$；同时改变参数量还需由式~\eqref{eq:q2-log-differential} 按当前 $\mathcal L_N,\mathcal L_D$ 加权，不能把 $\alpha,\beta,\gamma$ 当作脱离状态的全局重要性排序. 

\begin{proposition}[有限质量提升的等效参数规模]\label{prop:q2-finite-substitution-v15}
固定 $d>0$, 配比坐标 $z$ 和初始 $n_0>0$，比较质量 $Q_a$, $Q_b$. 令
\[
H:=n_0^{-\alpha}+\frac BA d^{-\beta}e^{\beta\lambda_{\mathrm{mix}}z}
\{e^{g(\ln Q_b)}-e^{g(\ln Q_a)}\}.
\]
若 $H>0$，则在旧质量 $Q_a$ 下与 $(n_0,d,Q_b,z)$ 等 Loss 的唯一有限参数规模为 $n_{\rm eq}=H^{-1/\alpha}$. 若 $H\le0$，不存在这样的有限正参数规模；$H=0$ 仅能由 $n\to\infty$ 达到极限. 
\end{proposition}
\begin{proof}
把 $\mathcal L(n_{\rm eq},d,Q_a,z)=\mathcal L(n_0,d,Q_b,z)$ 两边公共项消去，得到 $n_{\rm eq}^{-\alpha}=H$. 左侧严格递减地把 $(0,\infty)$ 映到 $(0,\infty)$，故存在性、唯一性与边界结论成立. 
\end{proof}

\begin{proposition}[质量坐标变换与有限响应不变性]\label{prop:q2-quality-coordinate-v37}
设 $z=\varphi(Q)>0$ 严格递增, 连续可微，$\varphi'(Q)>0$ 且 $\varphi(1)=1$. 同时变换响应函数及后续质量成本：
\[
\widetilde g(\ln z)=g(\ln\varphi^{-1}(z)),\quad
\widetilde h(z)=h(\varphi^{-1}(z)).
\]
则对应状态的 Loss、成本和最优资源决策不变，而局部质量弹性满足
\begin{equation}\label{eq:q2-coordinate-elasticity-v37}
\widetilde\gamma(z)=\gamma(Q)\frac{z}{Q\varphi'(Q)}.
\end{equation}
对应质量端点之间的 $g$ 差与有效数据倍率保持不变. 
\end{proposition}
\begin{proof}
单调坐标变换为可行质量区间之间的双射，代入模型和成本即得同一目标与可行资源集. 又由链式法则，
\[
\frac{\dd\widetilde g}{\dd\ln z}
=g'(\ln Q)\frac{\dd\ln Q}{\dd\ln z}
=g'(\ln Q)\frac{z}{Q\varphi'(Q)},
\]
除以 $-\beta$ 得弹性公式. 端点函数值未变，有限响应差及其指数倍率也未变. 
\end{proof}
因此，质量弹性的数值依赖所用坐标，有限替代关系则在一致变换下保持不变. %第一问文本质量代理与 B6 质量等级共用标尺仍需重叠锚点或额外校准；式~\eqref{eq:q1-general-identification-v37} 的识别范围不能未经该接口直接作为 B6 的质量区间. 

\subsection{秩一必要条件与参数不确定性}

\subsubsection{变量投影与嵌套选型}
记非线性参数 $\boldsymbol\eta=(\alpha,\beta,\delta,\delta_{\min})$，并把常指数模型视作 $\delta_{\min}=\delta$ 的受限情形. 固定 $\boldsymbol\eta$ 后，设计矩阵的第 $i$ 行为
\[
\mathbf x_i(\boldsymbol\eta)=
\bigl(1,n_i^{-\alpha},d_i^{-\beta}e^{g(\ln Q_i)}\bigr),
\quad \boldsymbol b=(E,A,B)^\top.
\]
拟合问题写成
\begin{equation}\label{eq:q2-variable-projection}
F(\boldsymbol\eta):=\min_{\boldsymbol b\in\mathcal B}
\|\mathbf W^{1/2}\{\boldsymbol y-\mathbf X(\boldsymbol\eta)\boldsymbol b\}\|_2^2,
\quad \min_{\boldsymbol\eta\in\mathcal H}F(\boldsymbol\eta),
\end{equation}
其中 $\mathcal B=[0,10]\times[10^{-8},100]^2$，非线性参数满足 $\alpha,\beta\in[0.005,1.5]$ 及式~\eqref{eq:q2-shape}，故端点斜率不超过 $7.5$. 普通拟合 $\mathbf W=\mathbf I$；整单元重抽样时其对角元为重复次数，零次数行删除，等价于重复观测. 数值估计采用变量投影：固定非线性参数后求解带界线性最小二乘，再在约束集合 $\mathcal H$ 上最小化剖面目标 $F(\boldsymbol\eta)$. 

\textbf{拟合性质：变量投影的局部导数：}
若 $\mathbf X$ 光滑，某邻域内线性子问题具有唯一解 $\boldsymbol b^*(\boldsymbol\eta)$，自由列满秩且活动边界固定，则剖面目标可微，并满足
\begin{equation}\label{eq:q2-projection-gradient}
\partial_{\eta_j}F=-2\boldsymbol r^\top\mathbf W
(\partial_{\eta_j}\mathbf X)\boldsymbol b^*,
\quad\boldsymbol r:=\boldsymbol y-\mathbf X\boldsymbol b^*.
\end{equation}
该公式仅用于优化器的局部导数计算；其成立依赖活动边界固定，活动集改变时需重新处理带界线性子问题. 

B6 的 360 行由 45 个完整 $(N,D)$ 单元与 8 个质量等级组成. 主检验分别整组留出一个质量等级和一个完整 $(N,D)$ 单元；另作留一 $N$ 与留一 $D$ 检验. 每个外层训练集内，再作质量等级留一和确定性五折完整单元检验. 仅当曲率模型在两项内层 RMSE 均降低至少 $1\%$、且赢过半数分组，才选曲率形式. B7 新增的两个质量等级共 90 行只作选型后的来源内检验；B8 不参与参数拟合或选择. 

\begin{table}[htbp]\centering\small
\caption{第二问两种质量模型的同口径分组验证}\label{tab:q2-v26-validation}
\begin{tabular}{lrrr}\toprule
验证口径 & 常指数 P 的 RMSE & 曲率 C 的 RMSE & 相对下降\\\midrule
B6 留一质量等级 & 0.0898402 & 0.0673465 & 25.04\%\\
B6 留一完整 $(N,D)$ 单元 & 0.0741682 & 0.0650484 & 12.30\%\\
B6 留一模型规模 & 0.0747492 & 0.0659274 & 11.80\%\\
B6 留一数据规模 & 0.0879606 & 0.0800434 & 9.00\%\\
B7 新质量等级 & 0.0580216 & 0.0508392 & 12.38\%\\\bottomrule
\end{tabular}
\end{table}
%表~\ref{tab:q2-v26-validation} 汇总同口径比较；图示见附录~\ref{app:v39-algorithms} 的图~\ref{fig:q2-v26-validation}. 质量等级的 8 个外层折和完整单元的 45 个外层折均选择曲率模型. 对外层分组误差进行描述性整组重抽样，RMSE 相对下降的 $2.5\%/50\%/97.5\%$ 分位分别为：质量留出 $8.39\%/24.60\%/38.67\%$，完整单元留出 $9.47\%/12.29\%/14.94\%$. 交叉验证训练折彼此重叠，这些分位不是独立重复实验的推断区间. 


\subsubsection{单调积分斜率的备选模型与增量验证}\label{sec:q2-shape-v38}
曲率函数把质量斜率限制为 $\ln Q$ 的仿射函数. 为检查这一参数化是否限制 B6 来源内的预测，在不改变式~\eqref{eq:unified-chain-v9} 及原模型选择结果的条件下，令 $u=\ln Q$，在 $[\ln0.1,0]$ 上取等距结点 $u_0<\cdots<u_4=0$，并设
\begin{equation}\label{eq:q2-integrated-slope-v38}
g_{\rm I}(u)=\sum_{j=0}^{3}s_j
\bigl[u_{j+1}-\max\{u,u_j\}\bigr]_+,
\quad 0\le s_j\le5\beta.
\end{equation}
这里 $[v]_+=\max(v,0)$. 在第 $j$ 个结点区间内部，$-g_{\rm I}'(u)=s_j$，从而有效数据弹性为 $\gamma_{\rm I}(Q)=s_j/\beta\in[0,5]$；结点处允许斜率跳变，而 $g_{\rm I}(0)=0$ 固定了质量 $1$ 的参照值. 对应 Loss 仍为 $E+An^{-\alpha}+Bd^{-\beta}\exp\{g_{\rm I}(\ln Q)\}$. 

\begin{proposition}[积分斜率质量响应的形状保证]\label{prop:q2-integrated-shape-v38}
若 $A,B,\alpha,\beta>0$ 且式~\eqref{eq:q2-integrated-slope-v38} 的斜率约束成立，则 $Q\mapsto g_{\rm I}(\ln Q)$ 在 $[0.1,1]$ 上连续、非增，Loss 关于质量非增；在每个开结点区间，其有效数据弹性为常数 $s_j/\beta\in[0,5]$. 若该区间的 $s_j>0$，质量提高使数据项严格下降. 
\end{proposition}
\begin{proof}
每一项 $[u_{j+1}-\max\{u,u_j\}]_+$ 连续且随 $u$ 非增，在 $(u_j,u_{j+1})$ 上的导数为 $-1$，其余开结点区间导数为零. 求和后得 $g_{\rm I}'=-s_j$；对 $Bd^{-\beta}e^{g_{\rm I}}$ 应用链式法则即得 Loss 的方向与严格性. 
\end{proof}

固定非线性参数 $\alpha,\beta,s_0,\ldots,s_3$ 后，对 $E,A,B$ 求带界最小二乘, 外层在 $\alpha,\beta\in[0.005,1.5]$, $0\le s_j\le5\beta$ 下作三起点约束优化. 五折验证按完整 $(N,D)$ 资源单元分组，每折只用 36 个训练单元重新拟合曲率模型与积分斜率模型，再评价 9 个留出单元的 72 行. 这个五折比较与表~\ref{tab:q2-v26-validation} 的逐单元留一设计不同，两组 RMSE 不横向拼接. 

同一 B6 完整资源单元五折比较中，曲率与积分斜率模型的合并 RMSE 分别为 $0.06669707$ 与 $0.06589885$；后者虽有微小优势，成对重抽样差值区间仍跨零. 因此积分斜率只是待检验的形状候选，不替换式~\eqref{eq:q2-v26-parameters} 的全 B6 主拟合或第三问参数. 完整表与逐折图见附录~\ref{app:v39-algorithms} 的表~\ref{tab:q2-integrated-cv-v38}、图~\ref{fig:q2-integrated-cv-v38}. B6 属半合成资料，B7 与其共享基础记录. 


\textbf{理论条件的程序检查：}
B6/B7 的中心化质量响应矩阵的第一奇异值能量占比分别为 $88.18\%$ 和 $87.07\%$，故共同质量响应的秩一必要条件是近似而非精确满足. 对 $z=Q^2$ 的一致坐标变换，101 个网格点的质量响应差在当前浮点计算中为零，且弹性按命题~\ref{prop:q2-quality-coordinate-v37} 更新；这些结构检查不证明两个来源共享质量机制. 

\subsubsection{参数、质量替代率与来源边界}
第一阶段用 B1 的 1176 条真实 Pythia 轨迹拟合经典 $N$--$D$ 分支，再将其冻结用于不同附件检验. 表~\ref{tab:q2-v39-external} 同时报绝对误差与排序相关：B2 的高 RMSE 与尚可的排序相关说明跨族 Loss 截距和量尺不可直接搬用；B3 是同源插值，不是 4000 次新增训练；B4、B5 的仅 $N$ 对照仍能得到高排序相关，故这两组结果支持粗排序，却不足以单独识别 Token 指数或质量斜率. 
\begin{table}[H]\centering\small
\caption{B1 固定参数的来源内外检验}\label{tab:q2-v39-external}
\begin{tabular}{lrrr}\toprule 数据与角色 & 行数 & RMSE & Spearman $\rho$\\\midrule
B2 半合成跨族 & 1029 & 1.2431 & 0.7881 \\
B3 Pythia 插值 & 4000 & 0.0038 & 1.0000 \\
B4 跨族 & 57 & 0.2927 & 0.9830 \\
B4 仅按 $N$ & 57 & 2.9768 & 0.9555 \\
B5 文献 & 44 & 0.1976 & 0.9588 \\
B5 仅按 $N$ & 44 & 3.1542 & 0.8542 \\
B10 估算对照 & 128 & 0.0011 & 1.0000 \\
\bottomrule\end{tabular}\end{table}

B9 的 132 条参数元数据覆盖 $N=100$--$10000$B，整体远离 B1 的 $0.07$--$11.97$B 观测规模；B10 的 128 条 Loss 为估算值. 表中 B10 的小误差只反映同一估算口径的数值对照，不是独立的大模型训练验证. 质量响应仍主要由 B6 的半合成 $N,D,Q$ 交叉表估计；B8 迁移检验出现反向排序，提醒跨生成机制外推的局限. 
全 B6 拟合为
\begin{align}\label{eq:q2-v26-parameters}
(E,A,B)&=(0.860047,\ 0.540173,\ 1.593498),\\
(\alpha,\beta,\delta)&=(0.279054,\ 0.095607,\ 0.229965),\\
(\delta_{\min},\kappa_Q)&=(0,\ -0.049936).
\end{align}
式~\eqref{eq:q2-v26-parameters} 是本文在 B6 上估计的曲率模型参数，第二、三问的全部质量相关结论均以其为输入；它与式~\eqref{eq:q2-power-curved} 的常指数模型 $\mathrm P$ 是同一半合成表上的两个并列拟合，后者给出 $\beta=0.069902$, $\gamma_Q=1.21646$. 两套参数的可比差异正是命题~\ref{prop:q2-shape} 所刻画的曲率项：常指数模型把质量弹性固定为常数，曲率模型允许其随 $Q$ 变化. 表~\ref{tab:q2-v26-validation} 给出两者的分组验证对照，第三问的模型形式敏感性则报告常指数设定下的全部结果.

计算使用完整精度. 此时 $\gamma(0.1)=0$, $\gamma(1)=2.40532$，且在质量锚点 $Q_0=0.6074742$ 处有 $\gamma(Q_0)=1.88464$. 最低质量端的零斜率来自活动约束；200 次完整 $(N,D)$ 单元联合抽样中，$\delta_{\min}$ 恰为 $0$ 的有 $128$ 次（$64\%$），绝对值不超过 $10^{-15}$（即浮点精度）的有 $196$ 次（$98\%$），最大值为 $4.8\times10^{-3}$. 因此它只支持当前受限模型在低质量端的形状，不能证明真实训练系统在 $Q=0.1$ 精确饱和. 

图~\ref{fig:q2-v26-quality} 的横轴为质量 $Q\in[0.1,1]$（对数刻度），纵轴为质量响应因子 $F(Q)=e^{g(\ln Q)}$. 固定参数规模、数据量和配比时，$F(Q)$ 等于质量为 $Q$ 与质量为 $1$ 时的数据损失项之比；例如 $F(Q)=1.2$ 表示前者的数据损失项高出 $20\%$. 蓝线为常指数模型，橙线为单调曲率模型，全 B6 拟合的两条响应曲线分别为
\[
F_{\mathrm P}(Q)\approx Q^{-0.085033},\quad
F_{\mathrm C}(Q)\approx\exp\!\left[-0.229965\ln Q-0.049936(\ln Q)^2\right].
\]
本文以完整 $(N,D)$ 资源单元为单位重抽样 $200$ 次，每次保留所抽单元内全部质量等级并重新拟合两种模型. 在每个固定 $Q$ 处，分别计算 $200$ 次拟合所得 $F(Q)$ 的 $2.5\%$ 与 $97.5\%$ 分位数，连接这些上下边界得到浅蓝、浅橙色带. 因此，实线表示完整数据的拟合结果，色带表示各质量水平下由参数估计波动形成的逐点 $95\%$ 分位范围. 所有拟合均满足 $F(1)=1$，故两条曲线及其色带在 $Q=1$ 处汇聚于 $1$. 

同一联合重抽样还给出曲率模型 $\gamma(Q_0)$ 的 $95\%$ 分位范围 $[1.71497,2.06961]$. 词数加权锚点下点值为 $1.97391$，范围为 $[1.79621,2.16765]$. 这些替代率范围由每次拟合的完整参数向量计算，保留式~\eqref{eq:q2-local-quality} 中 $\beta,\delta,\kappa_Q$ 的联合变化. 
\begin{figure}[H]\centering
\includegraphics[width=0.93\linewidth]{q2_v26_quality_response_with_joint_band.pdf}
\caption{常指数与单调曲率模型的质量响应及重抽样区间}\label{fig:q2-v26-quality}
\end{figure}
图~\ref{fig:q2-v26-quality} 比较常指数模型与单调曲率模型的质量响应因子. 蓝色与橙色实线分别表示基于完整 B6 数据的拟合结果. 对应浅色带由 200 次完整 $(N,D)$ 单元重抽样并重新拟合得到, 表示各 $Q$ 位置上的 $2.5\%$--$97.5\%$ 分位范围. 纵轴以 $Q=1$ 时的数据损失项为基准进行归一化, 因而两条响应曲线及其重抽样区间均在 $Q=1$ 处汇聚于 $1$.

在 $n=1,d=150,Q_0=0.6074742,\psi=0$ 下，将质量提高 $0.1$，常指数模型预测 Loss 降低 $0.0197105$、等效规模为 $1.14249$B；曲率模型分别为 $0.0308417$ 与 $1.23452$B. 两者均按命题~\ref{prop:q2-finite-substitution-v15} 在各自模型内计算；差异反映质量响应形式变化，不能把两组预测差解释为已经观测到的训练收益. 

图~\ref{fig:q2-v26-profiles} 给出参数剖面. 列归一化 Jacobian 条件数由常指数模型的约 $345$ 降至曲率模型的约 $196$，但参数仍有明显相关性. 单个 $(N,D)$ 单元的最大训练 SSE 份额为 $6.39\%$，前 $10\%$ 单元合计为 $26.87\%$，未触发预设稳健损失候选条件. 两处全 B6 拟合所得剖面极小值在不同初值下的目标差均小于 $3.0\times10^{-14}$；该结果表明数值解稳定，但并不把非凸外层的局部最优条件提升为全局最优定理. 
\begin{figure}[htbp]\centering
\includegraphics[width=0.93\linewidth]{q2_v26_parameter_profiles.pdf}
\caption{常指数模型与曲率模型的 $\beta,\delta$ 参数剖面. 参数组合的稳定性需要联合评估. }\label{fig:q2-v26-profiles}
\end{figure}

B1 的真实轨迹独立给出
\begin{align*}
(E,A,B)&=(1.68979756,0.35398032,1.24030558),\\
(\alpha,\beta)&=(0.33997658,0.27987813),
\end{align*}
留一模型规模 RMSE 为 $0.000151180$. 由于 $n,d$ 以十亿为单位，把系数换算回原始单位（$A_{\mathrm{raw}}=A\cdot10^{9\alpha}$, $B_{\mathrm{raw}}=B\cdot10^{9\beta}$）后为
$(E,A_{\mathrm{raw}},B_{\mathrm{raw}})=(1.6898,\ 406.23,\ 409.67)$；
与 Hoffmann 等报告的 $(E,A,B,\alpha,\beta)=(1.69,\ 406.4,\ 410.7,\ 0.34,\ 0.28)$\cite{hoffmann} 相比，
$E,\alpha,\beta$ 在四位有效数字上一致，$A_{\mathrm{raw}}$ 与 $B_{\mathrm{raw}}$ 分别相差 $0.04\%$ 与 $0.25\%$. 
这一对照说明本文的拟合流程能够在真实轨迹上复现文献级的经典标度律，故 B1 分支可作为跨来源检验的基准；
它同时说明第三问使用的 B6 指数（$\beta=0.0956$）与真实轨迹（$\beta=0.2799$）的差异是\textbf{数据来源}差异而非估计误差，
第~\ref{sec:q3-exponent-source} 节据此报告指数来源的敏感性情景. 
该结果只检查同族轨迹的拟合与泛化；B3 插值点不能增加独立样本量，B2、B4、B5 的跨族排序亦不证明绝对 Loss 标尺可迁移. 尤其 B6 的 45 个单元均表现为质量提高、Loss 下降，而 B8 校准与外推子集的质量斜率为正. 没有共同质量标尺和生成公式时，将它们合并估计会掩盖机制冲突. 

\subsubsection{第二问输出与证据边界}
曲率模型的选择依据来自原分组样本外误差，而定理~\ref{thm:q2-log-marginal-v14} 的有限等价关系与局部微分则由模型形式解析推出；二者承担不同证据角色. 表~\ref{tab:q2-v26-validation} 支持在 B6 来源内采用状态相关的 $\gamma(Q)$；附录表~\ref{tab:q2-integrated-cv-v38} 仅显示另一单调候选的五折误差略低，其配对差区间跨零，故不改变第三问使用的全 B6 曲率参数. 式~\eqref{eq:q2-v26-parameters} 给出的参数满足命题~\ref{prop:q2-shape} 的形状约束. 特别地，在后续优化区间 $Q\in[Q_0,1]$ 上，点估计满足
\[
\gamma(Q)\ge \gamma(Q_0)=1.88464>0,
\]
从而质量提高在该区间内对应严格正的有效数据弹性. 该经验事实使 $w(Q)=\mathcal G(Q)^{-1}$ 在决策区间严格递减，为第三问的全局单调性与饱和判据提供条件，但并不把 B6 的半合成质量机制提升为跨来源定律. B8 的反向响应、第一问有限映射覆盖和未识别的 $\lambda_{\mathrm{mix}}$ 仍构成独立的迁移限制. 

\section{问题三：算力约束下的资源配置与稳健优化}\label{sec:q3}

第二问给出的 $\gamma(Q)$ 随质量变化. 第三问由预算活动性得到固定质量唯一资源对，随后以质量成本下包络证明全局最优质量的预算单调性；完备候选用于数值求解，有限饱和充要条件和充分预算连接高预算结构. 在名义主线之后，再讨论参数不确定性、最大后悔与平均--尾部风险. 数学最优性均相对于给定 Loss 和成本代理，经验支持域与稳健性另行标记. 

\subsection{固定质量问题的唯一最优资源对}

\subsubsection{总预算与条件优化问题}
记 $c:=C/10^{18}$，$k:=6+\eta_{\mathrm{att}}L_{\mathrm{ctx}}$，$\eta_{\mathrm{att}}=2\times10^{-4}$. 单位 Token 的质量处理成本记为 $c_Q(Q)$，并令 $h(Q):=[c_Q(Q)-c_Q(Q_0)]/10^9$. 则总成本为
\begin{equation}\label{eq:q3-budget-v26}
C=6ND+\eta_{\mathrm{att}}L_{\mathrm{ctx}}ND+D\{c_Q(Q)-c_Q(Q_0)\},
\quad d[kn+h(Q)]\le c.
\end{equation}
主模型的 $Q_0=0.6074742$，三个质量成本情景为
\begin{equation}\label{eq:q3-cost-families}
h(Q)=\begin{cases}
0.01s_h(e^{6Q}-e^{6Q_0}),&\text{指数成本},\\
5s_h(Q^4-Q_0^4),&\text{幂成本},\\
2s_h\ln\!\dfrac{1+10Q}{1+10Q_0},&\text{对数成本},
\end{cases}\quad s_h>0.
\end{equation}
默认 $s_h=1$. 这些成本在 $[Q_0,1]$ 非负、光滑且递增；成本倍数及代理质量起点另作敏感性分析. C7 支持的上下文为 $2048,4096,8192,32768,131072$ Token. 注意力与基础训练成本的比值为 $\eta_{\mathrm{att}}L_{\mathrm{ctx}}/6$，在 $L_{\mathrm{ctx}}=30000$ 时相等，这是成本定义的直接结果. 
题面单位 Token 的质量成本系数经 $h=[c_Q(Q)-c_Q(Q_0)]/10^9$ 换算为式~\eqref{eq:q3-cost-families} 的 $0.01,5,2$，与 $c=C/10^{18}$ 的 B 单位资源表达一致. 把注意力项与基础训练项相等，可得
\begin{equation}\label{eq:q3-v39-critical-context}
6ND=\eta_{\mathrm{att}}L_{\mathrm{ctx}}ND
\quad\Longleftrightarrow\quad
L_{\mathrm{ctx}}^{\mathrm{crit}}
=\frac6{\eta_{\mathrm{att}}}=\frac6{2\times10^{-4}}=30000.
\end{equation}
C7 的五个可行值中，2048、4096、8192 三个低于临界值；32768、131072 两个高于临界值. 对应注意力与基础训练之比分别约为 $0.0683,0.1365,0.2731,1.0923,4.3691$；32768 是最邻近的高侧取值. 

固定配比令 $B_p:=B\exp(\beta\lambda_{\mathrm{mix}}\psi)>0$，主情景 $\psi=0$、$B_p=B$. 优化问题及其值函数为
\begin{equation}\label{eq:q3-value-v12}
V(c):=\min_{\substack{n,d>0,\ Q_0\le Q\le1\\d[kn+h(Q)]\le c}}
\{E+An^{-\alpha}+B_p d^{-\beta}e^{g(\ln Q)}\}.
\end{equation}
这里 $g$ 是第二问的质量响应，$h$ 是本问的处理成本，二者不能互换. 

\begin{lemma}[最优预算活动]\label{lem:q3-budget-active}
若 $A,B_p,k,c,\alpha,\beta>0$、$h\ge0$，则式~\eqref{eq:q3-value-v12} 的任何最优点均满足 $d[kn+h(Q)]=c$. 
\end{lemma}
\begin{proof}
若预算严格松弛，固定 $n,Q$ 后可适当增加 $d$. 因 $\partial_d\mathcal L=-\beta B_p d^{-\beta-1}e^{g(\ln Q)}<0$，Loss 严格下降，与最优性矛盾. 
\end{proof}
引理的前提没有额外数据量上限；把经验范围人为作为硬约束后，预算活动性需要重新检查. 

\begin{definition}[对数尺度剖面]\label{def:q3-profile}
对固定 $Q$，令 $x:=\ln n\in\mathbb R$，$s(x,Q):=ke^x+h(Q)>0$. 消去 $d=c/s$ 后定义
\begin{subequations}\label{eq:q3-profile-v26}
   \begin{equation}
       \ell_c(x,Q):=E+\mathcal L_N+\mathcal L_D,
   \end{equation}
   \begin{equation}
       \mathcal L_N:=Ae^{-\alpha x},\quad \mathcal L_D:=B_pc^{-\beta}s(x,Q)^\beta e^{g(\ln Q)},\quad r:=\frac{ke^x}{s(x,Q)}\in(0,1].
   \end{equation}
\end{subequations}
\end{definition}
\begin{theorem}[固定质量的唯一预算最优资源对]\label{thm:q3-unique-scale}
对任意固定 $Q\in[Q_0,1]$，若 $A,B_p,k,c,\alpha,\beta>0$ 且 $h(Q)\ge0$，则 $x\mapsto\ell_c(x,Q)$ 严格凸且强制，存在唯一极小点 $x^*(Q)$. 对应的$n^*(Q)=e^{x^*(Q)},\quad d^*(Q)=\frac c{kn^*(Q)+h(Q)}$
是固定质量预算子问题的唯一最优资源对，并由下式唯一确定：
\begin{equation}\label{eq:q3-scale-root}
c^\beta=\frac{\beta B_pk}{\alpha A}
n^{\alpha+1}[kn+h(Q)]^{\beta-1}e^{g(\ln Q)}.
\end{equation}
\end{theorem}
\begin{proof}
直接求导得
\begin{subequations}\label{eq:q3-convex-derivatives}
    \begin{align}
\ell_x&=-\alpha\mathcal L_N+\beta\mathcal L_Dr,\\
\ell_{xx}&=\alpha^2\mathcal L_N+
\beta\mathcal L_Dr\{(1-r)+\beta r\}>0.
\end{align}
\end{subequations}

当 $x\to-\infty$ 时 $\mathcal L_N\to\infty$；当 $x\to\infty$ 时 $s\sim ke^x$，故 $\mathcal L_D\to\infty$. 利用命题~\ref{prop:bg-convex-minimum-v26} 可得唯一全局极小点. 将 $\ell_x=0$ 整理即得式~\eqref{eq:q3-scale-root}；引理~\ref{lem:q3-budget-active} 使 $d$ 也被唯一决定. 另从右侧对 $\ln n$ 求对数导数得到 $\alpha+1+(\beta-1)r>0$，故也可由单调根唯一性直接求解. 
\end{proof}
%这一结果保留了常指数模型的尺度唯一性，但不要求 $\gamma(Q)$ 为常数. 严格凸性成立于 $x=\ln n$；不能据此断言完整 $(x,Q)$ 问题联合凸. 

\subsubsection{质量边际与预算收益均衡}
令 $a(Q)=-g'(\ln Q)=\beta\gamma(Q)$. 由链式法则
\begin{equation}\label{eq:q3-qderivative-v9}
\ell_Q=\mathcal L_D\left\{\frac{\beta h'(Q)}{kn+h(Q)}-\frac{a(Q)}Q\right\}.
\end{equation}
因此内点的质量条件是收益弹性与成本弹性的相等，而不再含单一常数质量指数：
\begin{equation}\label{eq:q3-quality-stationary}
Qh'(Q)=\gamma(Q)[kn+h(Q)].
\end{equation}
若 $a(Q)>0$，先由该式得到 $n_{stat}(Q)=[Qh'(Q)/\gamma(Q)-h(Q)]/k$；仅当此值为正时，才可代入式~\eqref{eq:q3-scale-root} 生成相应预算候选. 若 $a(Q)=0$，应使用未相除的式~\eqref{eq:q3-qderivative-v9}. 

\textbf{KKT 的单位算力边际均衡:}
设 $Q_0<Q<1$、$h'(Q)>0$，且为满足 KKT 条件的局部最优点. 其预算乘子 $\mu>0$ 满足
\begin{equation}\label{eq:q3-cost-marginal-v15}
\frac{\alpha An^{-\alpha-1}}{kd}
=\frac{\beta\mathcal L_D}{d(kn+h)}
=\frac{a(Q)\mathcal L_D}{Qdh'(Q)}=\mu.
\end{equation}
其中左右两式来自参数规模和质量方向, 中间式来自数据量方向，
该关系由对 $n,d,Q$ 的 KKT 驻点式分别除以对应正边际成本得到，用于解释内点候选的资源均衡；全局性仍由后续质量剖面比较给出. 
预算梯度的 $d$ 分量为 $kn+h>0$，内点满足相关约束资格；质量端点还需加入边界乘子. KKT 是候选筛选条件，全局性由后续质量比较给出. 

\subsection{仿射成本下包络与预算单调定理}\label{sec:q3-global-envelope-v37}

固定质量的边际随预算单调，只描述一条局部分支. 要比较不同预算的全局最优质量，需把质量与规模的共同作用重新表示为一族仿射成本. 

\begin{definition}[有效质量成本的下包络]\label{def:q3-quality-envelope-v37}
在固定配比、上下文及模型参数下，令
\begin{equation}\label{eq:q3-quality-envelope-v37}
w(Q):=e^{g(\ln Q)/\beta}=\mathcal G(Q)^{-1},\quad
\mathscr H(n):=\min_{Q\in[Q_0,1]}w(Q)[kn+h(Q)].
\end{equation}
其中 $w$ 为有效数据折算的倒数，$\mathscr H$ 为质量指标化仿射族的下包络. 
\end{definition}
预算活动性与正幂函数的单调性给出精确约化
\begin{subequations}\label{eq:q3-envelope-value-v37}
\begin{align}
L_c(n,Q)&:=\ell_c(\ln n,Q)
=E+An^{-\alpha}+B_pc^{-\beta}\{w(Q)[kn+h(Q)]\}^{\beta},\\
V(c)&:=\min_{n>0,Q\in[Q_0,1]}L_c(n,Q)
=E+\min_{n>0}\{An^{-\alpha}+B_pc^{-\beta}\mathscr H(n)^\beta\}.
\end{align}
\end{subequations}
命题~\ref{prop:bg-affine-envelope-v37} 表明 $\mathscr H$ 连续、凹且严格递增. 这里凹的是成本下包络，式~\eqref{eq:q3-envelope-value-v37} 的完整目标不由此成为凸函数. 

\begin{theorem}[全局资源配置的预算单调性]\label{thm:q3-global-monotonicity-v37}
设 $0<Q_0<1$，$A,B_p,k,\alpha,\beta>0$，$h$ 在 $[Q_0,1]$ 上连续且非负，$g$ 连续，且式~\eqref{eq:q3-quality-envelope-v37} 的 $w$ 严格递减. 所有参数、配比和成本均不随预算改变，资源变量除正性和总预算外不受额外上限约束. 则每个 $c>0$ 都存在全局最优解；对任意 $0<c_1<c_2$，任取两预算的全局最优解 $(n_1,Q_1)$ 和 $(n_2,Q_2)$，都有
\begin{equation}\label{eq:q3-global-monotonicity-v37}
n_1<n_2,\quad Q_1\le Q_2.
\end{equation}
%该结论允许质量和规模随预算跳跃，不要求同一预算的完整最优解唯一. 
\end{theorem}
\begin{proof}
由于 $w$ 在紧区间有正下界 $w_{\min}$，有 $\mathscr H(n)\ge kw_{\min}n$. 当 $n\to0^+$ 时参数损失项趋于无穷，当 $n\to\infty$ 时数据损失项趋于无穷. 因此目标关于尺度强制，质量区间又紧，全局最小值能够取得. 

先证明一个最优尺度处的活动质量唯一. 若不同质量 $Q_a,Q_b$ 在同一全局最优尺度 $n_*$ 上均达到 $\mathscr H(n_*)$，则两对 $(n_*,Q_a)$、$(n_*,Q_b)$ 都全局最优. 分别对正尺度 $n$ 使用一阶必要条件，得到
\[
\alpha An_*^{-\alpha-1}
=\beta B_pc^{-\beta}\mathscr H(n_*)^{\beta-1}kw(Q_a)
=\beta B_pc^{-\beta}\mathscr H(n_*)^{\beta-1}kw(Q_b).
\]
公共因子均为正，故 $w(Q_a)=w(Q_b)$；严格递减性迫使 $Q_a=Q_b$. 这只排除最优尺度处的不同活动质量，不排除普通尺度处的下包络折点. 

令 $t_i=B_pc_i^{-\beta}$，则 $t_1>t_2$. 两个预算的最优性分别给出
\begin{align*}
An_1^{-\alpha}+t_1\mathscr H(n_1)^\beta
&\le An_2^{-\alpha}+t_1\mathscr H(n_2)^\beta,\\
An_2^{-\alpha}+t_2\mathscr H(n_2)^\beta
&\le An_1^{-\alpha}+t_2\mathscr H(n_1)^\beta.
\end{align*}
相加得到 $(t_1-t_2)\{\mathscr H(n_1)^\beta-\mathscr H(n_2)^\beta\}\le0$. 由 $\mathscr H$ 严格递增，知 $n_1\le n_2$. 若相等，两预算在该尺度具有相同的唯一活动质量，而尺度一阶条件又要求 $t_1=t_2$，矛盾，故 $n_1<n_2$. 

最后令 $a_Q(n)=kw(Q)n+w(Q)h(Q)$. 质量在各自尺度最优意味着
\[
a_{Q_1}(n_1)\le a_{Q_2}(n_1),\quad
a_{Q_2}(n_2)\le a_{Q_1}(n_2).
\]
相加并消去截距，得到 $k\{w(Q_1)-w(Q_2)\}(n_2-n_1)\ge0$. 由 $n_2>n_1$ 与 $w$ 严格递减，得到 $Q_1\le Q_2$. 
\end{proof}

\begin{corollary}[饱和持续性与质量跳跃的规模响应]\label{cor:q3-persistent-saturation-v37}
在定理~\ref{thm:q3-global-monotonicity-v37} 的条件下，若 $c_0$ 有一个全局最优解满足 $Q=1$，则所有 $c>c_0$ 的每个全局最优解均满足 $Q=1$. 若某个正预算处最优质量的左右极限不同，则全局最优参数量也在该处向上跳跃. 
\end{corollary}
% 结论显然, 不予证明
%\begin{proof}
%第一结论直接由式~\eqref{eq:q3-global-monotonicity-v37} 与 $Q\le1$ 得到. 对第二结论，任意紧正预算区间内，前述强制性给出最优尺度的一致正下界和有限上界. 单调性保证左右极限存在，连续性保证极限仍为该预算的全局最优解. 若左右质量不同但尺度相同，就违反定理证明中的活动质量唯一性；结合尺度单调性，跳跃只能向上. 
%\end{proof}

几何上，求 $\mathscr H(n)$ 等于以方向 $(kn,1)$ 在线性目标下最小化曲线
$\{(w(Q),w(Q)h(Q)):Q\in[Q_0,1]\}$. 只有被这些方向支撑的下凸包部分可能成为最优质量，其余质量可排除. %取凸包不改变线性目标的最小值，但不表示实际语料质量可以按凸组合混合. 连续质量的筛选仍需控制曲线间隙，有限网格凸包只对其离散候选精确. 

本定理比固定质量边际引理更强，且不要求 $h$ 单调或 $\beta\le1$. 它暂不涵盖预算相关参数、资源硬上限、质量与 $N,D$ 的新交互项及多参数情景风险目标；也未断言原始 Token 数随预算单调. 若 $w$ 含平坦区，活动质量唯一性的证明失效，不能直接沿用全部严格结论. 

\subsection{质量剖面的完备候选与边界激活}

\begin{definition}[剖面目标的完备候选集]\label{def:q3-complete-candidate-v29}
设连续函数 $f$ 定义在闭有界区间 $I$ 上. 若 $\mathcal C\subseteq I$ 满足 $\argmin_I f=\argmin_{\mathcal C}f$，则称 $\mathcal C$ 为完备候选集. 完备性表示不遗漏任何全局最小点，不意味着该集合有限. 
\end{definition}

\begin{theorem}[质量剖面的完备候选与边界激活]\label{thm:q3-active-transition}
设 $0<Q_0<1$，$A,B_p,c,k,\alpha,\beta>0$，$h$ 在包含 $[Q_0,1]$ 的开区间上为 $C^1$、在该闭区间非负，$g$ 在对应对数区间的邻域为 $C^1$. 则 $n^*(Q)$ 及质量剖面 $\bar\ell_c(Q):=\ell_c(\ln n^*(Q),Q)$ 在闭区间连续、内部为 $C^1$，且
\begin{subequations}\label{eq:q3-candidate-set-v25}
    \begin{align}
\mathcal C(c)&:=\{Q_0,1\}\cup
\{Q\in(Q_0,1):Qh'(Q)=\gamma(Q)[kn^*(Q)+h(Q)]\},\\
\mathcal Q^*(c)&:=\argmin_{Q\in[Q_0,1]}\bar\ell_c(Q)
=\argmin_{Q\in\mathcal C(c)}\bar\ell_c(Q).
\end{align}
\end{subequations}
若 $Q_0$ 为全局最优质量，必要条件为 $Q_0h'(Q_0)\ge\gamma(Q_0)[kn^*(Q_0)+h(Q_0)]$；
若 $1$ 为全局最优质量，必要条件为 $h'(1)\le\gamma(1)[kn^*(1)+h(1)]$. 
\end{theorem}
\begin{proof}
固定质量的一阶条件 $\ell_x=0$ 对 $(x,Q)$ 连续可微，且 $\ell_{xx}>0$. 由背景隐函数引理~\ref{thm:bg-implicit-simple-root}，根在每个内部点有唯一 $C^1$ 延拓；唯一性把局部分支拼成同一分支，端点的一侧限制给出连续性. 背景剖面导数公式~\ref{prop:bg-profile-derivative} 与 $\ell_x=0$ 给出 $\bar\ell_c'(Q)=\ell_Q(x^*(Q),Q)$. 连续剖面在紧区间取得最小值，每个内部最小点必须使该导数为零，端点又已包含，故候选集完备. 两端的单边导数符号分别非负、非正，代入式~\eqref{eq:q3-qderivative-v9} 即得边界必要条件. 
\end{proof}

\begin{lemma}[固定质量边际的预算单调性与局部阈值]\label{lem:q3-fixed-quality-threshold}
在定理~\ref{thm:q3-unique-scale} 的条件下，固定 $Q\in[Q_0,1]$，并进一步假设 $h$ 在该点可微且 $\gamma(Q)>0$. 把唯一尺度解的预算依赖写为 $n^*(Q,c)$，定义质量边际符号函数
\begin{equation}\label{eq:q3-fixed-q-margin}
F_Q(c):=Qh'(Q)-\gamma(Q)\{kn^*(Q,c)+h(Q)\}.
\end{equation}
则 $c\mapsto n^*(Q,c)$ 为 $C^1$ 严格增函数，$c\mapsto F_Q(c)$ 为 $C^1$ 严格减函数. 令
\begin{equation}\label{eq:q3-local-margin-m}
m(Q):=\frac{Qh'(Q)}{\gamma(Q)}-h(Q).
\end{equation}
若 $m(Q)\le0$，则 $F_Q(c)<0$ 对所有 $c>0$ 成立，不存在正的局部边际阈值；若 $m(Q)>0$，则方程 $F_Q(c)=0$ 恰有一个正解 $c_{\rm loc}(Q)$. 记
\begin{equation}\label{eq:q3-local-n}
n_{\rm loc}(Q):=\frac{m(Q)}k,\quad
s_{\rm loc}(Q):=kn_{\rm loc}(Q)+h(Q)=\frac{Qh'(Q)}{\gamma(Q)},
\end{equation}
则该阈值可显式写为
\begin{equation}\label{eq:q3-local-threshold}
c_{\rm loc}(Q)=\left[
\frac{\beta B_pk}{\alpha A}
 n_{\rm loc}(Q)^{\alpha+1}
 s_{\rm loc}(Q)^{\beta-1}
 e^{g(\ln Q)}
\right]^{1/\beta}.
\end{equation}
对内部质量点，$\operatorname{sgn}\bar\ell_c'(Q)=\operatorname{sgn}F_Q(c)$；当 $Q=1$ 时，同一表达式给出左侧质量边际，$c_{\rm loc}(1)$ 即上边界的局部触边候选预算. 
\end{lemma}
\begin{proof}
对固定 $Q$，将式~\eqref{eq:q3-scale-root} 的右侧记为 $\Phi_Q(n)$. 其关于 $\ln n$ 的对数导数为
\[
\alpha+1+(\beta-1)\frac{kn}{kn+h(Q)}
\ge \alpha+\min\{1,\beta\}>0,
\]
故 $\Phi_Q$ 在 $(0,\infty)$ 上严格递增，且从 $0$ 连续趋于 $+\infty$. 由隐函数定理，$c^\beta=\Phi_Q(n^*)$ 唯一确定 $C^1$ 函数 $n^*(Q,c)$，并有 $\partial n^*/\partial c>0$. 因此
\[
F_Q'(c)=-\gamma(Q)k\,\frac{\partial n^*(Q,c)}{\partial c}<0.
\]
又因 $c\to0^+$ 时 $n^*(Q,c)\to0^+$，$c\to\infty$ 时 $n^*(Q,c)\to\infty$，故
\[
\lim_{c\to0^+}F_Q(c)=Qh'(Q)-\gamma(Q)h(Q)=\gamma(Q)m(Q),\quad
\lim_{c\to\infty}F_Q(c)=-\infty.
\]
若 $m(Q)\le0$，严格单调性给出 $F_Q(c)<0$；若 $m(Q)>0$，介值定理与严格单调性给出唯一正根. 根处 $kn^*+h(Q)=Qh'(Q)/\gamma(Q)$，由此得到式~\eqref{eq:q3-local-n}，再代入尺度方程~\eqref{eq:q3-scale-root} 即得式~\eqref{eq:q3-local-threshold}. 最后由式~\eqref{eq:q3-qderivative-v9} 的正因子分解可知剖面导数与 $F_Q$ 同号. 
\end{proof}
该引理只描述固定质量位置的局部边际，局部阈值仍不保证等于首次全局转移预算. 全局最优质量的预算单调性需由前述定理~\ref{thm:q3-global-monotonicity-v37} 的下包络比较另行证明. 

\begin{corollary}[质量改变时的尺度分配方向]\label{cor:q3-quality-scale-direction}
在上述条件下，若 $0<\beta\le1$、$h'\ge0$、$a\ge0$，则固定预算的尺度分支 $n^*(Q)$ 随 $Q$ 非降，相应 $d^*(Q)$ 非增. 
\end{corollary}
\begin{proof}
由式~\eqref{eq:q3-convex-derivatives} 和式~\eqref{eq:q3-qderivative-v9}，
\[
\ell_{xQ}=\beta\mathcal L_Dr
\left\{(\beta-1)\frac{h'}{kn+h}-\frac aQ\right\}\le0.
\]
对尺度驻点隐式求导，$\dd x^*/\dd Q=-\ell_{xQ}/\ell_{xx}\ge0$. 因 $n^*=e^{x^*}$，参数规模非降；而 $d^*=c/(kn^*+h)$ 的分母非降，数据量非增. 
\end{proof}
该方向结论比较同一预算、同一模型的不同固定质量分支；跨模型重估 $\alpha,\beta$ 后的配置变化仍须重新计算. 

定义 $\mathcal S_{\rm sat}:=\{c>0:1\in\mathcal Q^*(c)\}$，并令
\begin{equation}\label{eq:q3-transition-def-v12}
c_{\rm tr}:=\inf\mathcal S_{\rm sat},\quad\inf\varnothing:=+\infty.
\end{equation}
该定义本身不保证下确界取到. 满足严格质量单调条件时，推论~\ref{cor:q3-persistent-saturation-v37} 已给出饱和持续性；下节将以充要条件证明集合非空及首次饱和的上区间结构. 全局切换仍可能发生跳跃，不能预设连续触边. 等式 $h'(1)=\gamma(1)[kn^*(1)+h(1)]$ 只能确定局部边界候选. 一般 $C^1$ 函数的内部驻点可能无穷多，因此全局比较不能仅依赖有限网格或根枚举；本文以可核验的质量区间下界控制未枚举区间的最优值. 

\begin{lemma}[质量区间的解析下界]\label{prop:q3-quality-lower-bound}
设 $h$ 非降，$g$ 非增. 对任意质量区间 $[q_-,q_+]\subseteq[Q_0,1]$，定义
\begin{equation}\label{eq:q3-interval-lower}
\ell^-_{[q_-,q_+]}(x):=E+Ae^{-\alpha x}
+B_pc^{-\beta}[ke^x+h(q_-)]^\beta e^{g(\ln q_+)}.
\end{equation}
则 $\ell_c(x,Q)\ge\ell^-_{[q_-,q_+]}(x)$ 对所有 $x\in\mathbb R,Q\in[q_-,q_+]$ 成立，且右侧关于 $x$ 严格凸. 因此其最小值为整个质量区间最优值的下界. 
\end{lemma}
\begin{proof}
$h(Q)\ge h(q_-)$，而 $g(\ln Q)\ge g(\ln q_+)$；两项正因子分别取下界，其乘积仍为下界. 严格凸性由式~\eqref{eq:q3-convex-derivatives} 同样得到. 先对 $x$、再对 $Q$ 取下确界保持不等式. 
\end{proof}
%可行质量点上的尺度极小值给出上界. 内层求根用异号导数括住极小点，并由背景的凸函数支撑线给出下界；质量外层分裂当前下界最低的区间，直至全域上、下界之差达到预定精度. 这里使用的是含浮点余量的解析数值界，没有采用向外舍入的区间算术. 

\subsection{有限预算质量饱和的充要条件}

\begin{definition}[最终全局饱和]\label{def:q3-eventual-saturation-v29}
若存在 $c_0>0$ 使所有 $c\ge c_0$ 均满足 $\mathcal Q^*(c)=\{1\}$，则称完整问题最终全局饱和. 它要求此后唯一且持续最优，比某个预算上 $1\in\mathcal Q^*(c)$ 更强. 
\end{definition}

\begin{proposition}[固定尺度的质量饱和阈值]\label{prop:q3-envelope-threshold-v37}
在定理~\ref{thm:q3-global-monotonicity-v37} 的条件下，进一步设 $g(0)=0$，因而 $w(1)=1$. 定义扩展实数
\begin{equation}\label{eq:q3-envelope-threshold-v37}
n_{\rm env}:=\max\left\{0,\sup_{Q\in[Q_0,1)}
\frac{h(1)-w(Q)h(Q)}{k[w(Q)-1]}\right\}\in[0,+\infty].
\end{equation}
对任意有限正尺度 $n$，质量 $Q=1$ 达到 $\mathscr H(n)$ 当且仅当 $n\ge n_{\rm env}$；若 $n>n_{\rm env}$，则它是唯一活动质量. 
\end{proposition}
\begin{proof}
因 $w$ 严格递减，对所有 $Q<1$ 有 $w(Q)>1$. 质量 $1$ 不劣于该 $Q$ 等价于
\[
kn+h(1)\le w(Q)[kn+h(Q)]
\quad\Longleftrightarrow\quad
n\ge\frac{h(1)-w(Q)h(Q)}{k[w(Q)-1]}.
\]
对所有 $Q<1$ 同时成立即等价于式~\eqref{eq:q3-envelope-threshold-v37} 的阈值条件；严格超过上确界时，逐质量比较均为严格不等式. 
\end{proof}

\Needspace{12\baselineskip}
\begin{theorem}[有限预算质量饱和的充要条件]\label{thm:q3-finite-saturation-iff-v37}
在命题~\ref{prop:q3-envelope-threshold-v37} 的条件下，下列命题等价：
    \begin{enumerate}
        %\renewcommand{\labelenumi}{(\roman{enumi})}
        \item 存在有限正预算，使 $Q=1$ 属于全局最优解；
        \item 完整问题最终全局饱和；
        \item $n_{\rm env}<+\infty$.
    \end{enumerate}
\end{theorem}
\begin{proof}
若某个有限预算存在饱和全局解，其有限尺度必须满足 $n\ge n_{\rm env}$，故 (1) 推出 (3). 

下证 (3) 推 (2)，取可行比较方案 $Q=1,n=c^{1/2}$ 并用满预算，得到
$0<V(c)-E\le Ac^{-\alpha/2}+B_pc^{-\beta}[kc^{1/2}+h(1)]^\beta\longrightarrow0.$
任意全局最优尺度均满足 $0<A[n^*(c)]^{-\alpha}\le V(c)-E$，且 $V(c)-E\to0$，故 $n^*(c)\to+\infty$；上界与最优解的选取无关，因此该发散对所有全局最优尺度一致成立. 当 $n_{\rm env}<\infty$ 时，充分大预算的每个最优尺度都严格超过 $n_{\rm env}$，由命题~\ref{prop:q3-envelope-threshold-v37}，此时只有质量 $1$ 活动，故最终全局饱和. (2) 推出 (1) 由定义立即成立. 
\end{proof}

\begin{proposition}[单调质量响应下的可计算饱和充分界]\label{prop:q3-envelope-explicit-bound-v38}
在命题~\ref{prop:q3-envelope-threshold-v37} 的条件下，进一步设 $h$ 在 $[Q_0,1]$ 上非负、非降且可微，$H:=\max_{[Q_0,1]}h'(Q)<\infty$；$w(Q)=\exp\{g(\ln Q)/\beta\}$ 满足 $w'(Q)=-w(Q)\gamma(Q)/Q$，其中 $\gamma$ 正且随 $Q$ 非降. 任取 $Q_m\in(Q_0,1)$，有
\begin{equation}\label{eq:q3-envelope-explicit-bound-v38}
n_{\rm env}\le
\max\left\{0,\frac{h(1)}{k[w(Q_m)-1]},
\frac{w(Q_m)H}{k\gamma(Q_m)}\right\}<\infty.
\end{equation}
因此当前条件足以推出定理~\ref{thm:q3-finite-saturation-iff-v37} 的最终全局饱和. 
\end{proposition}
\begin{proof}
若 $Q\le Q_m$，式~\eqref{eq:q3-envelope-threshold-v37} 的分子不超过 $h(1)$，分母至少为 $k[w(Q_m)-1]$. 若 $Q\ge Q_m$，由 $w(1)=1$、$w'h\le0$ 及 $w\le w(Q_m)$，
\[
h(1)-w(Q)h(Q)=\int_Q^1(wh)'(t)\,\dd t
\le w(Q_m)H(1-Q).
\]
又因 $t\le1$、$w(t)\ge1$、$\gamma(t)\ge\gamma(Q_m)$，
\[
w(Q)-1=\int_Q^1\frac{w(t)\gamma(t)}{t}\,\dd t
\ge\gamma(Q_m)(1-Q).
\]
分别代入割线比并取上确界得到式~\eqref{eq:q3-envelope-explicit-bound-v38}；分母均正且有限. 
\end{proof}

\begin{corollary}[首次全局饱和的上区间结构]\label{cor:q3-saturation-upper-interval-v37}
在定理~\ref{thm:q3-finite-saturation-iff-v37} 的条件下，若 $n_{\rm env}<\infty$，则式~\eqref{eq:q3-transition-def-v12} 中的 $\mathcal S_{\rm sat}$ 非空. 若 $c_{\rm tr}>0$，有
\[
\mathcal S_{\rm sat}=[c_{\rm tr},+\infty),
\quad\mathcal Q^*(c)=\{1\}\quad(c>c_{\rm tr});
\]
若 $c_{\rm tr}=0$，则所有正预算都唯一质量饱和. 在正临界预算本身可以存在多个全局最优配置. 
\end{corollary}
%\begin{proof}
%定理给出非空性，推论~\ref{cor:q3-persistent-saturation-v37} 给出上区间性质. 若下确界为正，取饱和预算列趋于它；局部一致强制性使相应最优尺度有收敛子列，连续性使极限仍为该预算的饱和最优解，因此下确界属于集合. 若下确界为零，对任意正预算都能找到更小的饱和预算，再用持续性即可. 
%\end{proof}

该充要条件要求质量成本割线比有有限上界，不必要求端点质量弹性严格为正. 为说明其扩展范围，取 $a,b>0$，令

\begin{subequations}\label{eq:q3-zero-elasticity-example-v37}
\begin{align}
w(Q)&=1+a(1-Q)^2,\quad g(\ln Q)=\beta\ln w(Q),\\
h(Q)&=b\{(1-Q_0)^2-(1-Q)^2\}.
\end{align}
\end{subequations}

这时 $h\ge0$，$w$ 严格递减，而
$\gamma(Q)=\frac{2aQ(1-Q)}{1+a(1-Q)^2},\quad\gamma(1)=0.$
直接代入式~\eqref{eq:q3-envelope-threshold-v37} 得 $n_{\rm env}=b/(ak)<\infty$，故仍会有限饱和. 该例刻画判据的数学范围，不作为附件数据的新拟合结果. 

$n_{\rm env}$ 是固定尺度的质量阈值，不是预算阈值；计算 $c_{\rm tr}$ 仍须进行资源最优与质量全局比较. 上区间结构为二分定位提供理论依据，但每个预算的判定需要解析候选和区间界，不能仅凭浮点损失相等或粗网格接近端点. 下面保留一个可直接计算的充分预算，用于提供饱和区间的保守上端证据. 

算法~\ref{alg:q3-global-configuration} 将上述定理转为名义配置的求解步骤. 初始网格和可定位驻点只提供可行上界；未被定位的内部极小点仍由质量区间下界覆盖. 局部触边候选预算与首次全局饱和预算的比较另按第~\ref{sec:q3-results} 节进行. 
\begin{algorithm}[H]
\caption{第三问的全局质量--资源配置求解}\label{alg:q3-global-configuration}
\small
\begin{algorithmic}[1]
\Require 基线 $Q_0$，质量响应 $g$，成本 $h$，参数 $E,A,B_p,\alpha,\beta$，预算 $c$，上下文 $L_{\mathrm{ctx}}$，容差 $\varepsilon$
\Ensure 配置 $(n^*,d^*,Q^*)$，目标界 $[L,U]$，预算相对残差 $r_c$
\State $k\gets6+\eta_{\mathrm{att}}L_{\mathrm{ctx}}$；将 $[Q_0,1]$ 划为有限初始区间
\State $\mathcal P\gets\{Q_0,1\}\cup\{\text{初始网格点}\}$
\State 用 $\bar\ell_c'(q)=0$ 的变号区间定位内部驻点，并加入 $\mathcal P$ \Comment{定理~\ref{thm:q3-active-transition}}
\State $U\gets+\infty$；$q^*\gets Q_0$；$\mathcal I\gets\varnothing$
\For{$q\in\mathcal P$}
  \State 括住 $\ell_x(x,q)=0$ 的唯一根，求得 $x_q$ 与 $v_q=\ell_c(x_q,q)$ \Comment{定理~\ref{thm:q3-unique-scale}}
  \If{$v_q<U$} \State $U\gets v_q$；$q^*\gets q$；$x^*\gets x_q$ \EndIf
\EndFor
\For{每个初始区间 $I=[q_-,q_+]$}
  \State $L_I\gets\min_x\ell^-_I(x)$；把 $(I,L_I)$ 加入 $\mathcal I$ \Comment{引理~\ref{prop:q3-quality-lower-bound}}
\EndFor
\While{$\mathcal I\ne\varnothing$ 且 $U-\min_{I\in\mathcal I}L_I>\varepsilon$}
  \State 取出下界最小的区间 $I=[a,b]$；$m\gets(a+b)/2$
  \State 求 $x_m$ 与 $v_m=\ell_c(x_m,m)$；若 $v_m<U$，更新 $(U,q^*,x^*)$
  \State 对 $[a,m]$、$[m,b]$ 重算解析下界；保留下界小于 $U$ 的区间
\EndWhile
\State $L\gets\min\{U,\min_{I\in\mathcal I}L_I\}$，约定 $\min\varnothing=+\infty$
\State $Q^*\gets q^*$；$n^*\gets e^{x^*}$；$d^*\gets c/[kn^*+h(Q^*)]$
\State $r_c\gets\bigl|d^*[kn^*+h(Q^*)]/c-1\bigr|$ \Comment{核对 $U-L$ 与预算活动性}
\State \Return $(n^*,d^*,Q^*),[L,U],r_c$
\end{algorithmic}
\end{algorithm}

\begin{proposition}[保证最终全局饱和的充分预算]\label{thm:q3-sufficient-saturation}
在定理~\ref{thm:q3-active-transition} 的条件下，进一步设 $h'\ge0$，且 $\gamma_{\min}:=\min_{Q\in[Q_0,1]}\gamma(Q)>0$. 任选
\[
M>\max\left\{0,\max_{Q\in[Q_0,1]}
\frac{Qh'(Q)/\gamma(Q)-h(Q)}k\right\},
\]
并定义
\begin{equation}\label{eq:q3-sufficient-budget}
\bar c:=\max_{Q\in[Q_0,1]}
\left\{\frac{\beta B_pk}{\alpha A}
M^{\alpha+1}[kM+h(Q)]^{\beta-1}e^{g(\ln Q)}\right\}^{1/\beta}.
\end{equation}
则 $c>\bar c$ 时 $\mathcal Q^*(c)=\{1\}$. %从而完整问题最终全局饱和. 
\end{proposition}
\begin{proof}
连续函数在紧区间有界，且 $\gamma_{\min}>0$，故上述 $M$ 与有限 $\bar c$ 存在. 式~\eqref{eq:q3-scale-root} 的右侧关于 $n$ 严格递增；当 $c>\bar c$ 时，在每个 $Q$ 上必有 $n^*(Q)>M$. 因此 $kn^*(Q)+h(Q)>Qh'(Q)/\gamma(Q)$，由式~\eqref{eq:q3-qderivative-v9} 得剖面导数处处严格负，唯一全局最小质量为右端点 $1$. 取任意 $c_0>\bar c$ 即满足最终饱和定义. 
\end{proof}
本题的优化区间从 $Q_0>0.1$ 开始，点估计满足 $\gamma_{\min}=\gamma(Q_0)=1.88464>0$；三类质量成本在 $[Q_0,1]$ 上均为非减函数，因此命题~\ref{thm:q3-sufficient-saturation} 的关键假设在名义模型中成立. 充分预算 $\bar c$ 保证其后持续饱和，但并不等同于第一次出现 $Q^*=1$ 的预算. 数值部分采用式~\eqref{eq:q3-sufficient-budget} 的保守上界，以保持这一充分性解释. 

\begin{theorem}[质量饱和约束分支的高预算极限分配]\label{thm:q3-asymptotic}
设 $A,B_p,k,\alpha,\beta>0$，$h_1:=h(1)\ge0$，且 $g(0)=0$. 对每个 $c>0$，将质量固定为 $1$ 的预算问题有唯一最优资源对 $(n_1(c),d_1(c))$，两者均为 $C^1$. 当 $c\to\infty$ 时，$n_1\to\infty$、$h_1/(kn_1)\to0$，且
\begin{subequations}\label{eq:q3-asymptotic-v12}
    \begin{align}
n_1(c)&\sim K_n c^{\beta/(\alpha+\beta)},\quad K_n:=\left(\frac{\alpha A}{\beta B_pk^\beta}\right)^{1/(\alpha+\beta)},\\
d_1(c)&\sim(kK_n)^{-1}c^{\alpha/(\alpha+\beta)}.
\end{align}
\end{subequations}
设 $r_1:=kn_1/(kn_1+h_1)$，则精确参数预算弹性及三类成本极限份额为
\begin{subequations}\label{eq:q3-exact-elasticity-v26}
    \begin{equation}
\frac{\dd\ln n_1}{\dd\ln c}=\frac\beta{\alpha+1+(\beta-1)r_1},\quad  \frac{\dd\ln d_1}{\dd\ln c}=1-r_1\frac{\dd\ln n_1}{\dd\ln c},
  \end{equation}
  \begin{equation}
\frac{6 n_1}{kn_1+h_1}\longrightarrow\frac6k,\quad
\frac{\eta_{\mathrm{att}}L_{\mathrm{ctx}}n_1}{kn_1+h_1}\longrightarrow\frac{\eta_{\mathrm{att}}L_{\mathrm{ctx}}}k,
\quad \frac{h_1}{kn_1+h_1}\longrightarrow0.
\end{equation}
\end{subequations}

本定理首先陈述固定 $Q=1$ 的约束分支，未预先假定它在每个预算下都是完整问题的最优分支. 
\end{theorem}
\begin{proof}
唯一性来自定理~\ref{thm:q3-unique-scale}. 式~\eqref{eq:q3-scale-root} 在 $Q=1$ 时可写为
$n_1^{\alpha+1}(kn_1+h_1)^{\beta-1}=\frac{\alpha A}{\beta B_pk}c^\beta.$
左侧在正半轴严格递增，从零趋于无穷；其对数导数为 $\alpha+1+(\beta-1)r_1>0$. 隐函数定理给出 $C^1$ 性，且右侧随 $c\to\infty$ 发散迫使 $n_1\to\infty$. 以 $kn_1+h_1\sim kn_1$ 代入并取正幂根，得到 $n_1$ 的渐近式；预算等式给出 $d_1$. 对上述根方程作对数微分得到精确弹性，再对 $d_1=c/(kn_1+h_1)$ 微分得到数据弹性. 成本份额由预算中三项分别除以 $c$，最后取 $n_1\to\infty$ 的极限. 
\end{proof}

\begin{corollary}[质量饱和分支的二阶高预算展开]\label{cor:q3-second-order-asymptotic}
在定理~\ref{thm:q3-asymptotic} 的条件下，记
\[
m:=\alpha+\beta,\quad a:=\frac{\beta}{m},\quad
\varepsilon(c):=\frac{h_1}{kK_n}c^{-a}.
\]
则 $\varepsilon(c)\to0$，且当 $c\to\infty$ 时
\begin{subequations}
    \begin{align}
 n_1(c) &=K_n c^a\left[ 1+\frac{1-\beta}{m}\varepsilon -\frac{(1-\beta)(\alpha+1)}{2m^2}\varepsilon^2 +O(\varepsilon^3) \right],\label{eq:q3-n-second-order}\\
 d_1(c)&=\frac{c^{1-a}}{kK_n}\left[1-\frac{\alpha+1}{m}\varepsilon+\frac{(\alpha+1)(2\alpha+3-\beta)}{2m^2}\varepsilon^2 +O(\varepsilon^3) \right].\label{eq:q3-d-second-order}
\end{align}
\end{subequations}
特别地，参数量的一阶有限预算修正可写成
\begin{equation}\label{eq:q3-n-first-correction}
 n_1(c)=K_n c^a
 +\frac{1-\beta}{\alpha+\beta}\frac{h_1}{k}
 +O(c^{-a}).
\end{equation}
%若 $h_1=0$，则 $\varepsilon\equiv0$，式~\eqref{eq:q3-asymptotic-v12} 的主导幂律即为该约束分支的精确尺度关系. 
\end{corollary}
\begin{proof}
令 $t:=K_n c^a$、$y:=n_1/t$. 由定理~\ref{thm:q3-asymptotic} 的精确根方程和 $K_n$ 的定义，得到
\[
y^m\left(1+\frac{\varepsilon}{y}\right)^{\beta-1}=1.
\]
定义
\begin{equation*}
    F(y,\varepsilon):=m\ln y+(\beta-1)\ln\left(1+\frac{\varepsilon}{y}\right).
\end{equation*}
有 $F(1,0)=0$ 且 $\partial_yF(1,0)=m>0$. 故由隐函数定理，$y$ 在 $\varepsilon=0$ 的邻域内可光滑表示为 $\varepsilon$ 的函数. 令
\begin{equation}
    y=1+a_1\varepsilon+a_2\varepsilon^2+O(\varepsilon^3)
\end{equation}
代入 $F(y,\varepsilon)=0$ 并逐阶比较系数，得到
\begin{equation*}
    a_1=\frac{1-\beta}{m},\quad a_2=-\frac{(1-\beta)(\alpha+1)}{2m^2}.
\end{equation*}
这给出式~\eqref{eq:q3-n-second-order}. 又由
\begin{equation*}
    d_1=\frac{c}{kn_1+h_1}=\frac{c^{1-a}}{kK_n}\frac{1}{y+\varepsilon},
\end{equation*}
对 $(y+\varepsilon)^{-1}$ 作 Taylor 展开即得式~\eqref{eq:q3-d-second-order}. 将 $\varepsilon=(h_1/(kK_n))c^{-a}$ 代回参数量展开的一阶项，得到式~\eqref{eq:q3-n-first-correction}. 
\end{proof}

推论~\ref{cor:q3-second-order-asymptotic} 给出了主导幂律的收敛速度. 以 $L_{\mathrm{ctx}}=8192$、指数质量成本和曲率质量模型为例，$10^{22}$ 与 $10^{24}$ FLOPs 下分别有 $\varepsilon\approx0.0788$ 与 $0.0243$. 与表~\ref{tab:q3-v26-config} 的精确优化值比较，单独使用主导幂律时参数量相对误差分别约为 $-14.55\%$ 与 $-5.35\%$，加入二阶修正后降至约 $-0.49\%$ 与 $-0.020\%$；数据量误差由约 $24.91\%$ 与 $8.08\%$ 降至约 $3.54\%$ 与 $0.115\%$. 因此二阶项主要解释已经进入高预算区间后尚未完全达到极限幂律的有限预算偏差；在较低预算下 $\varepsilon$ 不够小，不把该展开作为全预算近似. 

\begin{corollary}[完整问题的条件渐近]\label{cor:q3-global-asymptotic-v29}
若完整问题最终全局饱和，则充分高预算下全局解等于上述 $Q=1$ 分支，故式~\eqref{eq:q3-asymptotic-v12}--\eqref{eq:q3-exact-elasticity-v26} 适用于完整问题. 定理~\ref{thm:q3-finite-saturation-iff-v37} 给出充要判据，命题~\ref{thm:q3-sufficient-saturation} 提供可计算的充分预算. 
\end{corollary}
\begin{proof}
由最终饱和有 $Q^*=1$；固定该质量的最优资源对唯一，故两问题的解从某一预算起相同. 
\end{proof}
曲率质量模型的点估计给出参数幂指数 $\beta/(\alpha+\beta)\approx0.2552$；常指数质量模型对应的指数约为 $0.2003$. 差异来自两种质量响应设定下重新估计的 $\alpha,\beta$，属于模型形式敏感性，而不是同一参数集上的纯解析差异. 

\textbf{无量纲缩放关系:}
若 $h(Q)=s_hH(Q)$，置 $m=kn/s_h$，则除去常数 $E$ 并除以正因子 $A(k/s_h)^\alpha$ 后，剖面目标为
\begin{equation}\label{eq:q3-dimensionless-scaling}
m^{-\alpha}+\Upsilon\,[m+H(Q)]^\beta e^{g(\ln Q)},
\quad \Upsilon:=\frac{B_p}{A}c^{-\beta}s_h^{\alpha+\beta}k^{-\alpha}.
\end{equation}
故固定模型中对应同一无量纲转移值 $\Upsilon_*$ 的预算分支满足
\[
c_*\propto (B_p/A)^{1/\beta}s_h^{(\alpha+\beta)/\beta}k^{-\alpha/\beta}.
\]
该缩放关系适用于所指定的转移分支；满足定理~\ref{thm:q3-finite-saturation-iff-v37} 时，首次饱和起点另由上区间结构确定. 
这是对剖面目标的精确变量代换，而不是新的最优性定理. 
当前点估计的质量成本倍数弹性为 $3.9188$，$k$ 弹性为 $-2.9188$. 一般联合参数扰动的资源稳定性由命题~\ref{prop:q3-parameter-stability} 控制；该界只适用于活动集不变的同一最优分支，跨质量分支切换必须重新比较. 

\begin{corollary}[最优值的预算灵敏度与影子价格]\label{prop:q3-shadow-price}
若某预算邻域中最优分支唯一且可微、活动集不变，并满足背景包络命题~\ref{thm:bg-envelope} 的正则条件，则 $V'(c)=-\mu^*(c)$，物理预算下为 $\dd V/\dd C=-\mu^*/10^{18}$. 
\end{corollary}
\begin{proof}
预算约束写作 $d(kn+h)-c\le0$，关于 $c$ 的偏导为 $-1$；将其代入包络公式即可. 
\end{proof}

\subsection{最优分支稳定性与稳健后悔}

前述结果完成了名义参数下的解析主线. 下面再回答第二问参数估计误差如何传递到第三问决策：先考察活动分支不变时的局部稳定性，再用有限参数情景的最大后悔值处理可能较大的联合扰动. 跨质量分支切换时，局部灵敏度界失效，仍须回到定理~\ref{thm:q3-active-transition} 的完备候选比较. 

\begin{proposition}[固定质量最优分支的参数稳定性]\label{prop:q3-parameter-stability}
固定预算 $c>0$、上下文系数 $k>0$ 与质量 $Q\in[Q_0,1]$. 为避免正参数的量纲直接影响扰动大小，可将平滑参数记为
\begin{equation}\label{eq:q3-stability-parameterization}
\boldsymbol\vartheta
:=\bigl(\ln A,\ln B_p,\ln\alpha,\ln\beta,\delta,\delta_{\min}\bigr)\in\Omega,
\end{equation}
其中 $A,B_p,\alpha,\beta>0$，$0\le\delta,\delta_{\min}\le5\beta$，且
$\kappa_Q=(\delta-\delta_{\min})/(2u_{\min})$、$u_{\min}=\ln0.1$；$\Omega$ 表示满足这些条件的参数域. 取 $\boldsymbol\vartheta_0,\boldsymbol\vartheta_1\in\Omega$，并假设线段
$\boldsymbol\vartheta(t)=(1-t)\boldsymbol\vartheta_0+t\boldsymbol\vartheta_1$，$0\le t\le1$，始终位于 $\Omega$. 记定理~\ref{thm:q3-unique-scale} 给出的唯一尺度解为 $x^*(t)$，并定义
\begin{equation}\label{eq:q3-stability-constants}
m_x:=\min_{0\le t\le1}\ell_{xx}(x^*(t),Q;\boldsymbol\vartheta(t)),\quad
L_{\vartheta}:=\max_{0\le t\le1}
\left\|\nabla_{\boldsymbol\vartheta}\ell_x(x^*(t),Q;\boldsymbol\vartheta(t))\right\|_2.
\end{equation}
则 $m_x>0$、$L_{\vartheta}<\infty$，且
\begin{subequations}
\begin{align}
\left|x^*(1)-x^*(0)\right|
&\le \frac{L_{\vartheta}}{m_x}
\left\|\boldsymbol\vartheta_1-\boldsymbol\vartheta_0\right\|_2,\label{eq:q3-stability-x}\\
\left|\ln\frac{n^*(\boldsymbol\vartheta_1)}{n^*(\boldsymbol\vartheta_0)}\right|
&\le \frac{L_{\vartheta}}{m_x}
\left\|\boldsymbol\vartheta_1-\boldsymbol\vartheta_0\right\|_2,\label{eq:q3-stability-n}\\
\left|\ln\frac{d^*(\boldsymbol\vartheta_1)}{d^*(\boldsymbol\vartheta_0)}\right|
&\le \frac{L_{\vartheta}}{m_x}
\left\|\boldsymbol\vartheta_1-\boldsymbol\vartheta_0\right\|_2.\label{eq:q3-stability-d}
\end{align}
\end{subequations}
在任一点上令 $M:=\alpha\mathcal L_N=\beta\mathcal L_Dr>0$，则两个最直接的局部灵敏度为
\begin{equation}\label{eq:q3-stability-AB}
\frac{\partial x^*}{\partial\ln A}=\frac{M}{\ell_{xx}}>0,
\quad
\frac{\partial x^*}{\partial\ln B_p}=-\frac{M}{\ell_{xx}}<0.
\end{equation}
若同一参数线段上 $Q$ 始终是完整问题的唯一全局最优质量，则式~\eqref{eq:q3-stability-n}--\eqref{eq:q3-stability-d} 同时给出完整最优资源对的有限扰动界. %若质量活动集或全局最优分支发生切换，则本命题不跨越该切换提供统一界. 
\end{proposition}
证明见\ref{app:q3-stability-proof}. 

\begin{corollary}[质量饱和分支的形状参数隔离]\label{cor:q3-saturated-shape-isolation}
在 $Q=1$ 的约束分支上，若第二问的质量响应满足 $g(0)=0$，则固定质量剖面化为
\[
\ell_c(x,1)=E+Ae^{-\alpha x}+B_pc^{-\beta}[ke^x+h(1)]^{\beta},
\]
% 这一段疑似没用
%其中不含 $\delta,\delta_{\min}$ 与 $\kappa_Q$. 因此在 $A,B_p,\alpha,\beta$ 及成本函数固定时，质量形状参数的扰动不会直接改变该饱和分支上的 $(n_1(c),d_1(c))$；它们只通过质量剖面的全局比较影响 $Q=1$ 是否成为或保持完整问题的最优质量. 
\end{corollary}
\begin{proof}
$Q=1$ 时 $\ln Q=0$，故曲率响应 $g(\ln Q)=-\delta\ln Q+\kappa_Q(\ln Q)^2$ 恒为零. 代入剖面目标即可；固定 $Q=1$ 后，尺度方程亦不含上述质量形状参数. 
\end{proof}
这两个结果把第二问的联合参数不确定性分成两类：$A,B_p,\alpha,\beta$ 直接扰动同一活动分支内的资源分配，而 $\delta,\delta_{\min},\kappa_Q$ 首先影响质量分支选择和转移位置. 当前 45 个名义配置均位于 $Q=1$，但联合参数重抽样若跨越质量切换边界，仍须使用完备候选集重新作全局比较，不能把局部稳定性界外推到另一分支. 

\begin{definition}[有限参数情景的最大后悔值]\label{def:q3-minimax-regret}
设 $\Theta=\{\theta_1,\ldots,\theta_M\}$ 为第二问完整联合参数行，所有情景共用同一成本和预算. 记各情景自身最优值为 $V_\theta(c)$，定义同一实际配置的后悔值和稳健问题为
\begin{equation}\label{eq:q3-regret-v26}
R_\theta(n,d,Q):=\mathcal L_\theta(n,d,Q)-V_\theta(c),
\quad \min_{d(kn+h)\le c}\max_{\theta\in\Theta}R_\theta(n,d,Q),
\end{equation}
同时要求 $n,d>0,Q\in[Q_0,1]$. 每个情景的比较基准不同，但候选实际资源必须相同. 
\end{definition}

\begin{lemma}[稳健尺度内层的唯一性及数值误差传递]\label{prop:q3-robust-convex}
若每个情景满足定理~\ref{thm:q3-unique-scale} 的正参数条件，则固定 $Q$ 的 $\max_\theta\{\ell_{c,\theta}(x,Q)-V_\theta(c)\}$ 严格凸且强制，具有唯一尺度极小点. 若各基准已满足 $V^-_\theta\le V_\theta\le V^+_\theta$，则质量区间下界应使用 $\max_\theta(\ell^-_\theta-V^+_\theta)$，可行点的后悔值上界使用 $\max_\theta(\ell_\theta-V^-_\theta)$. 
\end{lemma}
证明见 \ref{app:q3-robust-proof}. 

该定理说明若稳健内层不光滑, 则最优条件为零属于当前活动情景导数的凸包，数值下界通过次梯度支撑线构造；不能对最大值函数直接强行使用单一情景导数. 

\subsubsection{平均后悔与尾部风险的凸扩展}

最大后悔控制最不利单个情景，未直接约束平均代价. 为给这一权衡以明确目标，设情景权重 $\pi_s\ge0$、$\sum_{s=1}^M\pi_s=1$，并记 $R_s(x,Q)=\ell_{c,s}(x,Q)-V_s(c)$. 
\begin{definition}[有限情景的条件风险价值（CVaR）]\label{def:q3-cvar-v37}
设 $0<\rho<1$ 为尾部概率，定义
\begin{equation}\label{eq:q3-cvar-v37}
\operatorname{CVaR}_{\rho}(R):=
\min_{t\in\mathbb R}\left\{t+\frac1{\rho}
\sum_{s=1}^M\pi_s(R_s-t)_+\right\},\quad u_+:=\max\{u,0\}.
\end{equation}
其中 $t$ 是尾部阈值；例如 $\rho=0.1$ 表示最差 $10\%$ 情景的平均后悔. 该式是包含离散情景的尾部风险变分表示\cite{rockafellar_cvar}. 取 $0\le\omega\le1$，平均--尾部风险目标定义为
\[
(1-\omega)\sum_s\pi_sR_s+\omega\operatorname{CVaR}_{\rho}(R).
\]
\end{definition}
\begin{proposition}[固定质量的平均--尾部风险凸规划]\label{prop:q3-cvar-convex-v37}
    若各情景满足定理~\ref{thm:q3-unique-scale} 的正参数条件且共用同一成本，固定 $Q$ 后，上述目标的最小化等价于凸规划
    \begin{equation}\label{eq:q3-cvar-epigraph-v37}
        \begin{aligned}
            \min_{x,t,\boldsymbol u}\quad &
            (1-\omega)\sum_s\pi_sR_s(x,Q)
            +\omega\left[t+\frac1{\rho}\sum_s\pi_su_s\right],                 \\
            \text{s.t.}\quad              & u_s\ge R_s(x,Q)-t,\quad u_s\ge0.
        \end{aligned}
    \end{equation}
    还可把平均后悔作为给定上限的凸约束，单独最小化尾部风险. 
\end{proposition}
\begin{proof}
    每个 $R_s$ 是固定质量严格凸目标减去常数，故关于 $x$ 凸. 约束 $R_s(x,Q)-t-u_s\le0$ 为凸上图约束，其余约束仿射；目标为凸函数的非负加权和与线性项之和. 对每个 $s$ 取 $u_s=(R_s-t)_+$ 即得与定义的等价性. 
\end{proof}
基准误差满足 $V_s^-\le V_s\le V_s^+$ 时，逐情景有 $\ell_s-V_s^+\le R_s\le\ell_s-V_s^-$；平均与 CVaR 的单调性把这些界传递为风险界. 质量外层仍需全局比较，名义预算单调定理不自动适用于不同参数情景的风险目标. 

\begin{lemma}[平均--CVaR 的质量区间下界]\label{lem:q3-cvar-interval-v38}
    设各情景共用非降成本 $h$，且各 $g_s(\ln Q)$ 随 $Q$ 非增. 对质量区间 $[a,b]\subseteq[Q_0,1]$，用成本 $h(a)$、质量响应 $g_s(\ln b)$ 代入预算消元后的情景 Loss，记所得函数为 $\ell^-_{c,s}(x;a,b)$. 若 $V_s^-\le V_s\le V_s^+$，则整个区间内平均--CVaR 问题的目标下界为
    \begin{subequations}\label{eq:q3-cvar-interval-v38}
      \begin{align}
        L_{[a,b]}:&=\min_{x\in\mathbb R}\left[(1-\omega)\sum_s\pi_sR_s^-(x)
        +\omega\operatorname{CVaR}_{1-\rho}({\boldsymbol R^-}(x))\right],\\
         R_s^-(x):&=\ell^-_{c,s}(x;a,b)-V_s^+.
        \end{align}
    \end{subequations}
    可行点的实际风险上界则用情景 Loss 减去 $V_s^-$ 计算. 
\end{lemma}
\begin{proof}
    对任意 $Q\in[a,b]$ 和尺度 $x$，$h(Q)\ge h(a)$、$g_s(\ln Q)\ge g_s(\ln b)$. 预算消元后，情景数据项对这两量均非降，故 $\ell_{c,s}(x,Q)\ge\ell^-_{c,s}(x;a,b)$；再由 $V_s\le V_s^+$ 得逐情景后悔下界. 非负加权均值与 CVaR 均对逐坐标增大保持非降，最后对 $x$ 取极小即得式~\eqref{eq:q3-cvar-interval-v38}. 上界由 $V_s\ge V_s^-$ 同理得到. 
\end{proof}

\textbf{数值求解:}
固定 $Q$ 时，对式~\eqref{eq:q3-cvar-epigraph-v37} 在 $x=\ln n$ 上的凸次梯度作括区二分；质量外层按引理~\ref{lem:q3-cvar-interval-v38} 的最小下界优先分裂区间，并在可行质量点更新上界，直到两界相差不超过容差. 设计情景采用第二问的完整联合参数行，不拆开参数边际取极值. 保存的单情景基准最优值上下界继续传递到风险间隙；以下 8 个求解均在同一实际预算下计算 $d=c/[kn+h(Q)]$. 这些下界含浮点余量，并非采用向外舍入的形式化区间算术. 

%\paragraph{稳健性层的定位. \quad}
%参数稳定性命题描述同一活动分支内的连续扰动，最大后悔值则比较一组离散联合参数情景；它们连同平均--尾部风险刻画第二问参数不确定性的决策影响. 名义全局单调定理与风险目标具有不同假设，后者的最优质量必须独立比较. 

\subsection{第三问结果与数值验证}\label{sec:q3-results}

\subsubsection{名义配置与支持域}
在三类成本、五个上下文及 $10^{19},10^{22},10^{24}$ FLOPs 的 45 个主情景中，曲率质量模型的最优解均为 $Q=1$；常指数质量模型则有 42 个上界解和 3 个内点解. 这些是给定模型下的有限预算数值结果，模型形式比较仍需重新估计参数；全局持续性由推论~\ref{cor:q3-persistent-saturation-v37} 与相应假设保证，不能仅由有限情景数量推出. 表~\ref{tab:q3-v26-config} 给出指数成本、8192 Token 的代表配置. 
\begin{table}[htbp]\centering\small
    \caption{常指数与曲率质量响应下的最优资源配置}\label{tab:q3-v26-config}
    \begin{tabular}{lrrrrrr}\toprule
        预算        & 常指数 $N$ & 曲率 $N$  & 常指数 $D$  & 曲率 $D$   & 常指数 $Q$  & 曲率 $Q$ \\\midrule
        $10^{19}$ & 2.00064 & 1.82546 & 0.533925 & 0.568340 & 0.991354 & 1      \\
        $10^{22}$ & 5.75369 & 7.10178 & 210.082  & 172.718  & 1        & 1      \\
        $10^{24}$ & 12.8925 & 20.7635 & 9791.47  & 6163.27  & 1        & 1      \\\bottomrule
    \end{tabular}\end{table}

\begin{table}[H]\centering\footnotesize
\caption{$10^{22}$ FLOPs、指数质量成本下 C7 五个上下文的名义最优配置}\label{tab:q3-v39-context}
\begin{tabular}{rrrrrrr}\toprule $L_{\rm ctx}$ & $N^*(\text{B})$ & $D^*(\text{B})$ & $Q^*$ & 训练 (\%)& 注意力 (\%) & 质量 (\%)\\\midrule
2048 & 7.562 & 191.855 & 1.0 & 87.1 & 5.9 & 7.0 \\
4096 & 7.395 & 184.914 & 1.0 & 82.0 & 11.2 & 6.8 \\
8192 & 7.102 & 172.718 & 1.0 & 73.6 & 20.1 & 6.3 \\
32768 & 5.994 & 126.753 & 1.0 & 45.6 & 49.8 & 4.6 \\
131072 & 4.465 & 67.797 & 1.0 & 18.2 & 79.4 & 2.5 \\
\bottomrule\end{tabular}\end{table}


\begin{table}[H]\centering\small
\caption{指数质量成本下的质量饱和阈值}\label{tab:q3-v39-envelope}
\begin{tabular}{rrr}\toprule $L_{\rm ctx}$ & $n_{\rm env}$ 数值估计（B） & 保守上界\\\midrule
2048 & 1.0004 & 2.8616 \\
4096 & 0.9403 & 2.6897 \\
8192 & 0.8394 & 2.4013 \\
32768 & 0.5108 & 1.4611 \\
131072 & 0.1990 & 0.5694 \\
\bottomrule\end{tabular}\end{table}

\begin{table}[H]\centering\small
\caption{$10^{19}$ FLOPs、8192 Token 下三类质量成本的名义最优解}\label{tab:q3-v39-cost}
\begin{tabular}{lrrrrrr}\toprule 成本 & $N^*(\text{B})$ & $D^*(\text{B})$ & $Q^*$ & 训练 (\%) & 注意力 (\%) & 质量 (\%)\\\midrule
指数 & 1.8255 & 0.5683 & 1 & 62.2 & 17.0 & 20.8 \\
幂函数 & 1.9338 & 0.5238 & 1 & 60.8 & 16.6 & 22.6 \\
对数 & 1.2823 & 0.9366 & 1 & 72.1 & 19.7 & 8.3 \\
\bottomrule\end{tabular}\end{table}

表~\ref{tab:q3-v26-config} 比较常指数质量响应与曲率质量响应下的最优资源配置, 其中 $N$、$D$ 均以十亿为单位. “常指数”列使用该响应形式重新估计得到的 $(\alpha,\beta)=(0.27905,0.06990)$；“曲率”列对应本文主选模型, 其参数为 $(\alpha,\beta)=(0.27905,0.09561)$. 两类模型分别在各自目标函数下求解最优配置, 因此表中重点比较其资源分配结果及结构差异, Loss 数值仅在各自模型内部解释.
表~\ref{tab:q3-v39-context} 与表~\ref{tab:q3-v39-envelope} 显示 C7 五个上下文的同预算名义解. 在 $10^{22}$ FLOPs 档，上下文从 2048 增至 131072 Token 时，注意力成本份额由 $5.94\%$ 升至 $79.36\%$，$N$ 与 $D$ 同时下降；合并三档预算与五个上下文看（指数质量成本），该份额的完整范围为 $4.97\%$--$80.72\%$. 各档 $Q^*$ 均为 1. 质量饱和判据中的 $n_{\rm env}$ 数值估计由 $1.0004$B 降至 $0.1990$B；它是临界参数量的数值估计，并非首次饱和预算. 固定 $10^{19}$ FLOPs 与 8192 Token 时，三类成本分别给出表~\ref{tab:q3-v39-cost} 的 $(N,D)$；对数质量成本方案的 $N$ 低于指数方案而 $D$ 高于指数方案，但三者都选 $Q=1$. 不同成本函数的名义 Loss 目标口径不同，不能仅比较各行 Loss 数值而判定处理技术优劣. 

本节全部配置取自第二问式~\eqref{eq:q2-v26-parameters} 的曲率响应与同表常指数响应两套并列拟合，二者共享同一资源与成本结构，仅质量响应形式不同. 命题~\ref{prop:q2-shape} 保证曲率响应在 $[Q_0,1]$ 上单调，定理~\ref{thm:q3-unique-scale} 与定理~\ref{thm:q3-global-monotonicity-v37} 的全局结论对两套响应同时成立. 图~\ref{fig:q3-v26-config} 展示各成本下的配置变化. 不同质量响应模型具有不同的目标函数，因此其最优 Loss 不能直接作收益差；该图注给出同一对策略在两种模型口径下的交叉评价数值，说明资源选择对质量响应形式敏感. 曲率模型内部，允许优化质量相对固定 $Q_0$ 的 Loss 改善为 $3.49\%$--$5.28\%$，该数值只解释为给定模型内扩大决策空间的收益. 
\begin{figure}[htbp]\centering
    \includegraphics[width=0.94\linewidth]{q3_v26_configuration_comparison.pdf}
    \caption{常指数与曲率质量响应下的最优参数量与数据量. 所有方案均在各自模型和同一成本预算下求解.}\label{fig:q3-v26-config}
\end{figure}
B6 支持域为 $N\in[0.07,11.97]$B、$D\in[10,600]$B、$Q\in[0.1,1]$. 按支持域逐情景核对，$10^{19}$ 档数据量全部低于下界（$D$ 落在 $[10,600]$ 内的情景数为 $0/15$，实际区间 $0.278$--$1.034$B），但参数规模全部在范围内（$15/15$）；$10^{24}$ 档数据量全部高于上界（$0/15$，实际区间 $2207$--$7422$B），参数规模也\textbf{全部}高于上界（$0/15$，实际区间 $13.68$--$22.10$B）；只有 $10^{22}$ 档的 $N$ 与 $D$ 全部落在支持域内（各 $15/15$）. 因此三档之中只有中间档是数据支持域内的内插结果，两端均为模型外推. $Q=1$ 只表示模型变量达到上限，不表示现实语料达到完美质量. 

\textbf{数据--参数配比的健全性核对}
支持域之外还需核对解本身是否符合已知的训练实践量级. 把名义最优解写成 token/参数比 $D^*/N^*$，表~\ref{tab:q3-v41-dn-ratio} 给出与外部量级的对照. 

\begin{table}[htbp]\centering\small
    \caption{名义最优解的数据--参数比 $D^*/N^*$（token/参数）与外部量级的对照}\label{tab:q3-v41-dn-ratio}
    \begin{tabular}{lrrrr}\toprule
        预算        & $D^*/N^*$ & 经典算力最优 $\approx20$ & B1 真实轨迹 $97.3$ & 量级判断        \\\midrule
        $10^{19}$ & $0.311$   & 低约 $64$ 倍          & 低约 $313$ 倍     & 数据严重不足，不可执行 \\
        $10^{22}$ & $24.3$    & 同量级                & 低约 $4$ 倍       & 与经典算力最优一致   \\
        $10^{24}$ & $296.8$   & 高约 $15$ 倍          & 高约 $3.1$ 倍     & 数据相对过剩      \\
        \bottomrule\end{tabular}\end{table}
        
表~\ref{tab:q3-v41-dn-ratio} 比较名义最优解的数据–参数比 $D^*/N^*$ 与外部参考量级, 其中单位为 token/参数. 计算采用 $L_{\mathrm{ctx}}=8192$、指数质量成本和曲率质量响应. 本表的 $D^*/N^*$ 与表~\ref{tab:q3-v41-transfer} 中的 $N^*/D^*$ 互为倒数. 经典参考取 $D/N\approx20$, “B1 真实轨迹”则采用本文 B1 Pythia 日志中 $D/N$ 的中位数 $97.3$.
$10^{19}$ 档建议以 $1.83$B 参数只配 $5.68$ 亿 token，而同预算下的经典算力最优解约为 $D/N\approx20$（Hoffmann 等给出 $N\approx1.7$B、$D\approx36$B），本文 B1 真实 Pythia 轨迹的 $D/N$ 中位数为 $97.3$. 该差异的来源在第~\ref{sec:q3-exponent-source} 节报告：B6 半合成表估计的 $\beta=0.0956$ 远小于真实轨迹的 $0.2799$，使数据项的边际收益被显著低估，最优解因而把资源集中到参数规模与质量上界. 因此\textbf{$10^{19}$ 档的 $(N^*,D^*)$ 不应被解释为可执行的训练配方}，只应解释为“在该半合成响应面下数据相对便宜”这一模型内部推断；三档之中只有 $10^{22}$ 档（$D/N=24.3$）与经典算力最优量级一致，$10^{24}$ 档则相反地落在数据相对过剩一侧. 

\textbf{结构变化的具体构成} 既然三档预算之间不存在质量维度转移，需要回答的是这段时间里\textbf{发生了哪一种结构变化}. 把 $d[kn+h(Q)]=c$ 的三项分别除以预算，得到三类成本的份额，连同最优配置的预算弹性一并列入表~\ref{tab:q3-v41-transfer}. 三档预算下质量份额由 $20.75\%$ 单调降至 $6.31\%$ 与 $2.25\%$，训练份额相应上升. 在 $L_{\mathrm{ctx}}=8192$ 下注意力份额基本稳定；若把 C7 的五个可行取值一并计入，注意力份额则大幅变动（表~\ref{tab:q3-v39-context}），说明它是上下文维度而非预算维度的主导项. 因此该区间的结构变化是\textbf{质量开支的系统性退出}，而不是质量水平的提升或下降：$Q^*$ 始终为 $1$，改变的是为维持该上界所付出的相对代价. 最优配置的预算弹性与推论~\ref{cor:q3-global-asymptotic-v29} 给出的极限 $(\beta/(\alpha+\beta),\alpha/(\alpha+\beta))=(0.2552,0.7448)$ 方向一致，数据量弹性在全部情景下都明显高于参数量弹性；逐档弹性见表~\ref{tab:q3-v41-transfer}. 

\begin{table}[htbp]\centering\small
    \caption{三档预算下最优配置的成本份额与预算弹性}\label{tab:q3-v41-transfer}
    \begin{tabular}{llrrrrrr}\toprule
        质量成本 & 预算        & 训练份额      & 注意力份额     & 质量份额               & $N/D$     & $\mathrm{d}\ln N/\mathrm{d}\ln c$ & $\mathrm{d}\ln D/\mathrm{d}\ln c$ \\\midrule
        指数   & $10^{19}$ & $62.25\%$ & $17.00\%$ & $\mathbf{20.75\%}$ & $3.212$   & $0.197$                           & $0.828$                           \\
        指数   & $10^{22}$ & $73.60\%$ & $20.10\%$ & $\mathbf{6.31\%}$  & $0.0411$  & $0.233$                           & $0.776$                           \\
        指数   & $10^{24}$ & $76.78\%$ & $20.97\%$ & $\mathbf{2.25\%}$  & $0.00337$ & ---                               & ---                               \\
        幂函数  & $10^{19}$ & $60.78\%$ & $16.60\%$ & $22.62\%$          & $3.692$   & $0.192$                           & $0.835$                           \\
        幂函数  & $10^{22}$ & $72.88\%$ & $19.90\%$ & $7.22\%$           & $0.0435$  & $0.230$                           & $0.781$                           \\
        对数   & $10^{19}$ & $72.06\%$ & $19.68\%$ & $8.27\%$           & $1.369$   & $0.231$                           & $0.779$                           \\
        对数   & $10^{22}$ & $77.14\%$ & $21.07\%$ & $1.79\%$           & $0.0313$  & $0.249$                           & $0.754$                           \\
        \bottomrule\end{tabular}\end{table}
表~\ref{tab:q3-v41-transfer} 给出三档预算下的最优配置、成本份额与预算弹性. 计算采用 $L_{\mathrm{ctx}}=8192$、曲率质量响应及 $Q_0=0.607474$. 其中, 成本份额表示各成本项占总预算的比例, 预算弹性采用相邻预算档之间的对数差分计算.

\textbf{预算区间内的结构变化:} 需要区分最优质量水平的转移与成本结构的转移. 对名义曲率模型，8192 Token 下指数质量成本的首次全局饱和预算约为 $6.05\times10^{16}$ FLOPs，低于题目给定的最低预算 $10^{19}$ FLOPs. 因此在 $10^{19}\to10^{22}\to10^{24}$ FLOPs 的三档预算中，最优质量始终为 $Q^*=1$，没有发生质量维度的再次切换.

但质量处理成本的预算占比由 $10^{19}$ FLOPs 时的 $20.75\%$ 降至 $10^{22}$ 和 $10^{24}$ FLOPs 时的 $6.31\%$ 和 $2.25\%$，并在两档之间跌破 $10\%$. 因此，三档预算之间虽然没有最优质量水平的转移，却发生了明显的成本结构变化：质量处理从重要开支项逐渐退化为边际开支项.

近优剖面进一步刻画最优决策的不敏感范围. 在 $10^{19}$ FLOPs、8192 Token、指数成本下，允许绝对 Loss 恶化 $0.001$，逐质量重优化规模得到数值质量区间 $[0.994042,1]$；固定 $Q=1$ 时，参数规模区间为 $[1.546449,2.160025]$B. 前者来自 161 点括区加细化，后者利用固定质量的凸剖面. 这是两个不同的一维近优集合，不是统计置信区间，不能拼成联合可行矩形.


\textbf{配比迁移的条件敏感性:}
第二问的 $\lambda_{\rm mix}$ 在锚点数据上不可识别. 固定第 239 行配方的 $\psi=-0.00644$、指数成本和 $10^{19}$ FLOPs，分别给定 $\lambda_{\rm mix}=0,1,2$ 并重新求名义解，得到表~\ref{tab:q3-v39-mix}. 配比差值进入 $B_p=B\exp(\beta\lambda_{\rm mix}\psi)$，故它改变 Loss 和最优资源比，但表列变化是条件情景，不是 B6 对配比收益的验证. 

现在我们直接回答“配比能否联立优化”的问题：$\lambda_{\rm mix}$ 由 $0$ 增至 $2$ 时 Loss 由 $2.998659$ 单调降至 $2.996589$，相对变化 $0.069\%$；最优参数量由 $1.8255$B 增至 $1.8295$B（$+0.22\%$），数据量由 $0.5683$B 降至 $0.5674$B（$-0.17\%$），最优质量始终为 $1$. 与该量级相比，第二问质量项的模型内 Loss 改善为 $0.0197$--$0.0308$，第一问配比响应面在 1M 检验集上的相对预测误差（宏平均 RMSE 相对目标均值）约为 $4.2\%$. 也就是说，在附件可支持的 $\psi$ 量级（$|\psi|=0.0064$）上，配比项所能引起的 Loss 变化比响应面自身的预测误差小约 $60$ 倍，也比质量项的收益小一个数量级以上. 因此\textbf{在本附件的数据条件下，把配比作为 $N,D,Q$ 之外的联立决策变量没有可识别的收益}；第一问的命题~\ref{prop:q1-quality-collinearity-v15} 已说明配比与质量系数不能分别识别，此处进一步说明即便给定 $\lambda_{\rm mix}$，其影响也远小于质量项. $\lambda_{\rm mix}$ 更大（例如 $|\lambda_{\rm mix}\psi|\geq 0.1$）时该结论不成立，需重新求解. 
\begin{table}[H]\centering\small
\caption{第 239 行配方相对锚点的条件迁移敏感性}\label{tab:q3-v39-mix}
\begin{tabular}{rrrrrr}\toprule $\lambda_{\rm mix}$ & $\psi(\pvec)$ & $N^*$(\text{B}) & $D^*$(\text{B}) & $Q^*$ & Loss\\\midrule
0 & -0.00644 & 1.8255 & 0.5683 & 1 & 2.998659 \\
1 & -0.00644 & 1.8275 & 0.5678 & 1 & 2.997624 \\
2 & -0.00644 & 1.8295 & 0.5674 & 1 & 2.996589 \\
\bottomrule\end{tabular}\end{table}

表~\ref{tab:q3-v39-mix} 比较第 239 行配方相对锚点配方的条件迁移敏感性. 计算采用 $C=10^{19}$ FLOPs、指数质量成本和 $L_{\mathrm{ctx}}=8192$, 其中 $N,D$ 均以 $\text{B}$ 为单位.
\subsubsection{最大后悔值改善及其平均代价}
稳健设计使用 200 行完整曲率参数，并另构造 200 次完整 $(N,D)$ 单元的独立联合重抽样以评价样本扰动敏感性；常指数参数抽样单独用于模型形式敏感性. 表~\ref{tab:q3-v26-regret} 先在每个情景内计算相对下降，再取 18 个情景的中位数；负数表示变差，最坏 $10\%$ 均值为每组 200 个后悔值中最大的 20 个的均值. 
\begin{table}[htbp]\centering\small
    \caption{稳健策略相对名义策略的后悔值下降(18 个情景的中位数)}\label{tab:q3-v26-regret}
    \begin{tabular}{lrrrr}\toprule
        参数评价集      & 最大值     & 平均值        & 最坏 $10\%$ 均值 & 最大值改善 \\\midrule
        设计抽样（曲率）   & 29.48\% & $-9.32\%$  & $-4.84\%$    & 18/18 \\
        独立重抽样（曲率）  & 24.45\% & $-16.17\%$ & $-14.46\%$   & 18/18 \\
        形式敏感性（常指数） & 25.87\% & 20.95\%    & 18.78\%      & 18/18 \\\bottomrule
    \end{tabular}\end{table}
稳健配置控制最大后悔值，但新曲率重抽样的平均及尾部均值反而增大，故不构成所有风险指标上的一致优势. 图~\ref{fig:q3-v26-regret} 报出最坏端收益；名义代价见附录~\ref{app:v39-algorithms} 的图~\ref{fig:q3-v26-cost}；可行插值策略的折中曲线不是已证明完整的 Pareto 前沿. 
\begin{figure}[htbp]\centering
    \includegraphics[width=0.93\linewidth]{q3_v26_regret_tradeoff_across_scenarios.pdf}
    \caption{稳健策略在设计参数、独立重抽样及模型形式敏感性中的后悔值折中. }\label{fig:q3-v26-regret}
\end{figure}

若质量处理成本加倍并保持原实际 $(N,D,Q)$ 不变，预算利用率为 $1.0064$--$1.2438$，全部超预算. 因此不允许按每个成本情景暗中重算 $D$ 来声称共同可行. 成本减半、加倍后另求的 36 个变成本配置均可行（$2$ 档预算、$3$ 个上下文、$3$ 类成本，减半与加倍各 $18$ 个）；它们是成本确定后分别可选的方案. 

\textbf{质量基线口径的敏感性:}\label{par:q3-quality-baseline} 假设 H6 的第 (iii) 点指出一个尚未闭合的口径链：质量处理代价按 Token 计量，而 $Q_0$ 取自文档等权均值. 把基线换成锚点的词数加权对应值 $Q_0^{\mathrm{tok}}=0.6617$（与文档等权值相差 $8.9\%$）后重解，结果见表~\ref{tab:q3-v41-q0}. 三档预算下最优质量均为 $1$，$N$ 的变化不超过 $1.35\%$、$D$ 不超过 $1.95\%$；由于 $g_s(0)=0$ 且成本基准 $h(Q_0)$ 在作差中消去，基线的绝对水平对名义配置的影响很小. 但数值饱和预算由 $6.05\times10^{16}$ 变为 $8.19\times10^{16}$ FLOPs，变化约 $35\%$. 因此口径选择不影响“三档预算均饱和”这一结论（两值都远低于 $10^{19}$），却会按同一比例平移转移点的位置；5.6.4 节报告的全部转移预算都应按此比例理解为口径相关的量. 
\begin{table}[htbp]\centering\small
    \caption{质量基线口径对第三问结果的影响($L_{\mathrm{ctx}}=8192$)}\label{tab:q3-v41-q0}
    \begin{tabular}{lrrrrr}\toprule
        基线                & 预算                                           & $N$（B）                                       & $D$（B）      & $Q^*$ & 质量份额       \\\midrule
        文档等权 $Q_0=0.6075$ & $10^{19}$                                    & $1.8255$                                     & $0.5683$    & $1$   & $20.75\%$  \\
        词数加权 $0.6617$     & $10^{19}$                                    & $1.8008$                                     & $0.5794$    & $1$   & $20.30\%$  \\
        相对变化              &                                              & $-1.35\%$                                    & $+1.95\%$   & ---   & $-0.45$ pp \\\midrule
        文档等权 $Q_0=0.6075$ & $10^{22}$                                    & $7.1018$                                     & $172.7182$  & $1$   & $6.31\%$   \\
        词数加权 $0.6617$     & $10^{22}$                                    & $7.0639$                                     & $174.0310$  & $1$   & $6.10\%$   \\
        相对变化              &                                              & $-0.53\%$                                    & $+0.76\%$   & ---   & $-0.21$ pp \\\midrule
        文档等权 $Q_0=0.6075$ & $10^{24}$                                    & $20.7635$                                    & $6163.2677$ & $1$   & $2.25\%$   \\
        词数加权 $0.6617$     & $10^{24}$                                    & $20.7204$                                    & $6181.4311$ & $1$   & $2.17\%$   \\
        相对变化              &                                              & $-0.21\%$                                    & $+0.29\%$   & ---   & $-0.08$ pp \\\midrule
        数值饱和预算（FLOPs）     & \multicolumn{2}{r}{文档等权 $6.05\times10^{16}$} & \multicolumn{2}{r}{词数加权 $8.19\times10^{16}$} & $+35\%$                          \\\bottomrule
    \end{tabular}\end{table}

\subsubsection{平均--CVaR 配置的求解与第二组重抽样}\label{sec:q3-cvar-results-v38}
取尾部概率 $\rho=0.1$、200 行 B6 完整单元联合参数设计抽样，分别在三种质量成本、$C=10^{19},10^{22}$ FLOPs 和 $L_{\rm ctx}=8192$ 下求 $\omega=0.5$ 的平均--CVaR 配置；对低预算幂成本再求 $\omega=0$ 与 $1$，总计 8 次全局质量搜索. 风险目标的最大数值上下界差为 $1.97913\times10^{-6}$，200 个情景基准最优值的最大数值差为 $1.43783\times10^{-11}$，合并后不超过 $1.97915\times10^{-6}$；最大预算相对残差为数值零. 另把每组保存的情景后悔输入独立的 $t,u_s$ 线性规划，所得 CVaR 与排序尾部公式最大差为 $5.99\times10^{-17}$. %上述均是浮点核查，不等同形式化精度证书；各量均在 $10^{-6}$ 量级或以下. 

在设计参数上，六个 $\omega=0.5$ 配置的加权目标均低于对应的冻结名义与最大后悔配置. 评价时改用第二组 200 个完整资源单元的联合重抽样；仍以各情景自身最优值为基准，表~\ref{tab:q3-cvar-evaluation-v38} 报告相对冻结名义配置的差值，正数为变差. 
\begin{table}[htbp]\centering\small
    \caption{新平均--CVaR 配置相对名义配置的第二组重抽样后悔差(单位为 $10^{-5}$ Loss)}\label{tab:q3-cvar-evaluation-v38}
    \begin{tabular}{llrr}\toprule
        预算/FLOPs  & 质量成本 & 平均后悔差   & $\operatorname{CVaR}_{0.9}$ 差 \\\midrule
        $10^{19}$ & 指数   & $+1.96$ & $+6.65$                       \\
        $10^{19}$ & 幂函数  & $+1.95$ & $+5.76$                       \\
        $10^{19}$ & 对数   & $+2.90$ & $+11.11$                      \\
        $10^{22}$ & 指数   & $+0.97$ & $-4.10$                       \\
        $10^{22}$ & 幂函数  & $+0.82$ & $-4.49$                       \\
        $10^{22}$ & 对数   & $+1.93$ & $-2.10$                       \\\bottomrule
    \end{tabular}
\end{table}
六组评价平均后悔差均为正；低预算尾部差也为正，高预算尾部差为负. 将同一组 200 个评价情景成对重抽样 5000 次，上述六组平均差与六组尾部差的描述性 $95\%$ 分位区间均跨过零. 附录~\ref{app:v39-algorithms} 的图~\ref{fig:q3-cvar-evaluation-v38} 展示低预算幂成本的五种配置：名义方案在这组评价参数上的平均与尾部后悔均低于新风险配置. 风险目标的设计端改善因此不能写成已验证的评价收益；两组参数抽样虽然分开，仍取自同一 B6 半合成表，不是新的真实模型训练实验. 


\subsubsection{局部触边、全局数值转移与充分界}
\begin{table}[htbp]\centering\small
    \caption{8192 Token 下的质量饱和预算（FLOPs）}\label{tab:q3-transition-v8}
    \begin{tabular}{lrrr}\toprule
        成本  & 局部边界候选                  & 全局数值转移                  & 高预算充分界                  \\\midrule
        指数  & $6.048404\times10^{16}$ & $6.048404\times10^{16}$ & $1.905752\times10^{20}$ \\
        幂函数 & $6.577382\times10^{14}$ & $5.104194\times10^{15}$ & $9.020844\times10^{19}$ \\
        对数  & 无正候选                    & $1.013163\times10^{13}$ & $4.221165\times10^{16}$ \\\bottomrule
    \end{tabular}\end{table}
表~\ref{tab:q3-transition-v8} 区分局部候选、数值切换与充分界，图~\ref{fig:q3-v26-transition} 对照两种质量模型的全局数值转移. 局部边界候选列由引理~\ref{lem:q3-fixed-quality-threshold} 在 $Q=1$ 处的显式公式~\eqref{eq:q3-local-threshold} 给出：指数与幂成本分别存在唯一正的局部阈值，而对数成本满足 $m(1)\le0$，因而不存在正的局部候选. 全局数值转移仍需在完备质量候选之间比较；幂成本的局部触边不等于全局切换，对数成本则可在没有局部触边候选的情况下发生全局分支跳跃. 推论~\ref{cor:q3-saturation-upper-interval-v37} 已在当前名义模型满足的条件下证明饱和集合具有唯一上区间起点；表中数值仍是有限扩展扫描与细化所得的近似定位，尚非该起点的形式化区间证书. 保守充分界通常远大于数值切换，质量跳跃也不违反全局单调性. 

\textbf{对题目三档预算的直接回答:}
在当前曲率质量响应下，8192 Token 的指数质量成本首次全局数值饱和约为 $6.048404\times10^{16}$ FLOPs，低于题目三档预算的最低值 $10^{19}$. 因而 $10^{19}\to10^{22}\to10^{24}$ FLOPs 的名义配置均保持 $Q^*=1$，主要改变 $N,D$ 与成本份额，质量维度在这三档之间没有再次转移. 表~\ref{tab:q3-v26-config} 给出指数成本下 $N:1.825\to7.102\to20.764$B、$D:0.568\to172.718\to6163.27$B. 

这一条件答案还可用归档的联合参数扰动直接核对. 对 200 行 B6 曲率模型的完整联合抽样，在五个 C7 上下文、三档预算和指数质量成本下逐一重解全局质量问题，共 $200\times5\times3=3000$ 个配置均有 $Q^*=1$，最大报告目标上下界差为 $1.53\times10^{-11}$. 在 8192 Token 下，200 个\textbf{局部边界候选}预算的中位数、$95\%$ 分位及最大值依次为 $4.04\times10^{16}$、$6.13\times10^{17}$、$3.10\times10^{18}$ FLOPs；五个上下文的局部候选最大值依次为 $4.20,3.77,3.10,1.32,0.261\times10^{18}$ FLOPs. 局部候选不被当作逐次全局切换点；上述 3000 次逐情景全局重解才是三档饱和的数值证据. 命题~\ref{prop:q2-shape} 的上界在名义拟合和这 200 次联合抽样均未激活：$\delta/(5\beta)$ 的抽样范围为 $[0.411,0.556]$、触边 $0/200$. 这些是同一半合成来源内的参数稳健性，不消除下文质量响应形式与规模指数来源造成的结构差异. %逐次配置和输入摘要见 \path{data/reports/review_v44/}. 

\textbf{该结论有两重条件，必须与结论本身一起引用: }
第一重是质量响应形式：改用常指数质量响应后首次转移约为 $1.3559\times10^{19}$ FLOPs，落在题目区间内，``是否存在观察区间内的质量转移"依赖质量响应形式. 
第二重是指数来源：第~\ref{sec:q3-exponent-source} 节的 S1 与 S2 情景把规模指数换成真实轨迹的估计值后，$10^{19}$ 档最优质量退回下界 $Q_0$，转移同样重新落在区间内. 

因此本文对题目第五段（结构性转移）给出的答案是\textbf{条件答案}：在 B6 半合成表估计的指数与曲率响应这一组合下不存在区间内转移，在另外三类设定下均存在；判定的分歧只能靠 B6 支持域之外的、同时含 $N,D,Q$ 的真实观测来消除. 
此外低、高预算的数据量超出 B6 支持域（表~\ref{tab:q3-v41-dn-ratio} 与第~\ref{sec:q3-results} 节的支持域核对），最低档解的数据--参数比仅 $0.311$ token/参数，这些配置应作为模型外推结果. 
\textbf{下包络阈值的直接核查:}
对归档的五种上下文和三类成本，共 15 组预算轨迹、615 个预算点，程序逐组检查全局最优 $N$ 与 $Q$ 随预算非降，并在 8192 Token 的三类成本转移附近另作 15 次质量全局搜索；新搜索的最大目标间隙为 $2.28932\times10^{-7}$、最大预算相对残差为 $2.22\times10^{-16}$. 式~\eqref{eq:q3-envelope-explicit-bound-v38} 在 15 组均给出有限上界. 8192 Token 时，指数、幂函数、对数成本的 $n_{\rm env}$ 数值估计分别为 $0.8394,0.5231,0.0604$B，解析保守上界分别为 $2.4013,1.9841,0.2804$B. 前一组三个值是连续质量割线比的浮点最大化估计，不是精确上确界；后一组只保证在本模型条件下阈值有限，也不等于首次全局饱和的预算. 
\begin{figure}[htbp]\centering
    \includegraphics[width=0.92\linewidth]{q3_v26_transition_model_sensitivity.pdf}
    \caption{不同上下文下两种质量模型的全局数值转移预算.}\label{fig:q3-v26-transition}
\end{figure}
在该上下文与指数成本下，常指数模型的全局数值转移约为 $1.3559\times10^{19}$ FLOPs，而曲率模型约为 $6.0484\times10^{16}$，相差超过两个数量级. %曲率模型对应的转移位于经验数据量支持域之外，因此主要用于刻画模型结构敏感性；局部候选或充分界的参数抽样分位也不能解释为全局首次转移的统计置信区间. 

定理~\ref{thm:q3-unique-scale} 将固定质量问题化为唯一尺度分支，定理~\ref{thm:q3-active-transition} 再把完整问题缩减为端点与内部驻点构成的完备候选集；稳健情景则由命题~\ref{prop:q3-robust-convex} 控制最大后悔值. 数值求解仅用于比较这些解析候选及给出有限区间上的全局上下界；全部报告配置（45 个名义配置、18 个基准成本稳健配置及 36 个变成本配置）的目标间隙与预算相对残差均不超过 $10^{-6}$. 这些误差量说明数值结果与解析目标一致，但不改变 B6 半合成质量机制的证据等级. 

\subsubsection{规模指数来源的敏感性}\label{sec:q3-exponent-source}
% 这一段或许需要重新考虑表达
前述名义结果的目标函数系数 $(E,A,B,\alpha,\beta,\delta,\kappa_Q)$ 全部来自 B6 半合成表. 
同一论文用真实 Pythia 轨迹（B1）独立估出经典分支的 $(E,\alpha,\beta)=(1.68980,0.33998,0.27988)$，
而 B6 给出 $(0.86005,0.27905,0.09561)$. 两组参数的截距相差 $0.830$ nat；在名义 $10^{19}$ 档配置
$(n,d)=(1.8255,0.5683)$ 上，B6 口径给出 $\mathcal L=2.9987$，B1 口径给出 $\mathcal L=3.4311$，
即\textbf{同一配置下相差 $0.432$ nat}. 为避免把该损失误读为真实交叉熵水平，此处明确：第三问报告的损失只有同一模型内的相对意义，
不能与真实交叉熵水平直接比较. 规模指数的差异更具实质影响：
$\beta/(\alpha+\beta)$ 由 $0.2552$ 变为 $0.4515$，而该比值正是高预算下最优参数量的主导幂指数. 

为量化这一影响，本文构造两个条件情景并重解三档预算下的名义问题（$L_{\mathrm{ctx}}=8192$，指数质量成本）：
\textbf{S1} 只把 $(\alpha,\beta)$ 换成 B1 的估计值、其余参数保持 B6；
\textbf{S2} 把整个经典 $N$--$D$ 分支换成 B1 的 $(E,A,B,\alpha,\beta)$，质量响应仍取 B6（B1 不含质量维度）. 
两情景都只是参数来源的情景替换，不是对 B1 的重新拟合，也不是新的数据证据. 结果见表~\ref{tab:q3-exponent-source}. 

\begin{table}[htbp]\centering\footnotesize
    \setlength{\tabcolsep}{5pt}
    \caption{规模指数来源对名义最优配置的影响}\label{tab:q3-exponent-source}
    \setlength{\tabcolsep}{3pt}
    \begin{tabular}{llrrrr}\toprule
        情景          & 预算        & $N^*$    & $D^*$     & $Q^*$    & 质量份额      \\\midrule
        S0 名义（全 B6） & $10^{19}$ & $1.825$  & $0.568$   & $1$      & $20.75\%$ \\
                    & $10^{22}$ & $7.102$  & $172.718$ & $1$      & $6.31\%$  \\
                    & $10^{24}$ & $20.764$ & $6163.27$ & $1$      & $2.25\%$  \\\midrule
        S1 仅换规模指数   & $10^{19}$ & $0.229$  & $5.721$   & $0.6075$ & $0$       \\
                    & $10^{22}$ & $6.621$  & $184.421$ & $1$      & $6.73\%$  \\
                    & $10^{24}$ & $49.390$ & $2625.28$ & $1$      & $0.96\%$  \\\midrule
        S2 真实轨迹分支   & $10^{19}$ & $0.173$  & $7.551$   & $0.6075$ & $0$       \\
                    & $10^{22}$ & $5.130$  & $233.442$ & $1$      & $8.52\%$  \\
                    & $10^{24}$ & $37.549$ & $3442.74$ & $1$      & $1.26\%$  \\\bottomrule
    \end{tabular}
\end{table}
表~\ref{tab:q3-exponent-source} 比较不同规模指数来源下的名义最优配置. 计算采用 $L_{\mathrm{ctx}}=8192$ 和指数质量成本, 预算单位为 FLOPs, $N,D$ 均以 $\text{B}$ 为单位. 三种情景分别依据各自的损失口径求解, 因而重点比较最优资源配置及其结构变化.
与名义结果最关键的差异出现在最低档预算：S1 与 S2 的最优质量都退到下界 $Q_0$（质量份额为零），
而 $10^{22}$ 与 $10^{24}$ 两档仍取上界 $Q^*=1$. 也就是说，\textbf{在真实轨迹估计的规模指数下，
    质量维度的转移重新出现在题目要求的三档预算之内}，其位置在 $10^{19}$ 与 $10^{22}$ 之间；
名义模型的"三档之间不存在质量维度转移"因此是\textbf{以 B6 半合成表估计的指数为条件}的结论，
而非模型无关的结论. 参数量方向也随指数来源改变：最低档最优参数量下降约一个数量级，
最高档则近似翻倍（逐档数值见表~\ref{tab:q3-exponent-source}）. 
质量份额的退出序列（表~\ref{tab:q3-v41-transfer}）同样只在 S0 下成立. 

需要强调的是，S1 与 S2 都保留 B6 的质量响应 $g$ 与成本 $h$，而 B6 的 $(\alpha,\beta)$ 与 $g$ 是联合估计的，故这两个情景不能解释为"真实系统的配置"，只能说明名义结论对指数来源的依赖强度. 在同时含 $N$、$D$、$Q$ 的真实训练数据出现之前，本文不据此替换名义主结果， 但把"转移是否存在"的判定改为条件表述，并在摘要与结论中并列两种答案. 上述情景是参数来源的情景替换，不是新的数据证据. 
%上述情景的逐配置数值与预算残差见 \path{data/reports/rechecks/q3_exponent_source_sensitivity.csv}. 

\textbf{第三问输出与跨问接口:}
第三问以曲率质量模型的名义配置作为主结果，最大后悔与平均--CVaR 配置用于展示不同风险目标的代价. 新风险配置已经求解，但第二组同源重抽样未确认稳定收益，故不替换名义配置. 第四问只传递经过完整质量剖面比较后的训练分量 $C_{\mathrm{train}}=6ND$ 及其联合参数不确定性；Benchmark 能力仍在第四问中独立建模. 

\section{问题四：技术演进与能力前沿}\label{sec:q4}


前三问主要以验证 Loss为目标，第四问则转向 \textbf{Benchmark} 能力分析. 本文先统一综合能力的计算口径再利用 C3 和 C8 检查历史记录，并核对各项任务的评分尺度. 在此基础上，本文研究模型规模与时间的变化关系，并分析 Loss 与 Benchmark 之间的关系. 使用条件高分位数表示能力前沿，后通过\textbf{有序凸规划}估计该前沿，并通过历史截点回测和有限样本校准检查模型的适用范围. 随后，把第三问得到的可比训练 FLOPs 路径接入第四问，并计算未来 12 个月和 24 个月的条件预测结果. Loss 与 Benchmark 存在联系，但两者并不等价. 因此，本文只把二者之间的关系用于相同口径下的比较和检查.

\subsection{历史截断与一致性计}
能力分数的时间轴采用 C1 的提交日期 \texttt{Submission Date}, C4 中模型身份、参数量、许可证和开源权重等信息的可得性则按发布日期 \texttt{Publication date} 截断. 由于 C1 无统一的版本日期字段, 故不采用版本日期作为时间依据. 历史截点的身份匹配仅使用当时已有的信息, 后续仓库中新增的别名不用于反向追溯历史记录的模型身份，以保持历史时间点上的身份判定口径一致. C4 的 3523 条元数据构成当前可用的统一整理的数据版本, 且记录修改时间均晚于预测起点. C8 共包含 1958 个版本文件, 其中 4 个为无效文件. 对每个历史截点, 取该时点之前最近的可解析版本. 本文进一步逐项重建 C8 的随机基线归一化与任务汇总, 用于检查不同记录在统一能力标尺下的可比性. C3 的历史层与重复排行榜层分别进行审计, 未经证实的零值不纳入统一六任务回归.

\subsection{综合能力标尺与评测口径重建}

第四问首先确定预测目标和评分口径，并沿用 C1 中六项任务的完整平均分作为综合能力指标。利用 C3 检查历史层与排行榜层的一致性，并利用 C8 校正随机猜测基线后统一各项任务的评分尺度。不同的度量方法可能改变能力曲线，也可能影响表观跃迁的判断\cite{schaeffer_mirage}。因此，必须统一六项任务的评分口径，再分析能力的时间变化和能力前沿。

本文以 C1 中 IFEval、BBH、MATH Lvl 5、GPQA、MUSR 和 MMLU-PRO 六项完整成绩的算术均值定义综合能力 $S_i\in[0,100]$. 重新计算所得均值与 C1 官方平均分之间的最大绝对误差为 $1.42\times10^{-14}$, 说明两者在数值精度范围内一致. C1 共包含 4576 条原始记录, 经过许可证筛选、有效性检查和去重后, 得到 1708 条用于观察性分解的样本. 在此基础上，进一步按照开放权重与模型身份一致性进行筛选，最终得到 33 条严格前沿样本. 三个样本集合依次嵌套,严格前沿样本来自 18 个发布组织, 其中预训练模型 22 条, 后训练模型 11 条. 严格身份匹配要求完整 owner/path，或发布者及模型版本词元均一致；仅 basename 相同的条目不能合并使用，将C1/C4 参数量相对差超过 $10\%$ 的匹配排除. 同仓库同日若类型、许可证、参数量或六项成绩相冲突，则不按较高分挑选版本. 

\subsubsection{开源筛选口径:}按题目要求明确入选规则，严格前沿样本依次施加四项条件：
\begin{enumerate}[label=\textbf{(\roman*)},leftmargin=*,itemsep=0pt,topsep=0pt]
\item \textbf{开源权重:} 须与 C4 的 \texttt{Open model weights?} 严格前沿样本要求模型同时满足开放权重与唯一身份匹配条件，因此许可证合规本身并不足以构成入选依据.
\item \textbf{许可证:} \texttt{Hub License} 字段须属于 apache-2.0、mit、bsd-2-clause、bsd-3-clause、bsd-3-clause-clear、cc-by-4.0、cc-by-sa-4.0、cc0-1.0、isc、mpl-2.0 十项允许研究与复现的白名单.
 \item \textbf{去重与冲突:} 按仓库取截至截点的最新记录, 同一仓库同日的类型、许可证、参数量或六项成绩若互相冲突则整组剔除.
\item \textbf{身份与规模一致性:} C1/C4 参数量相对差不超过 $10\%$, 且发布日期不晚于截点.
\end{enumerate}



\subsubsection{C3 一致性与 C8 逐任务口径重建}
C3 共 4599 行, 其中历史层 26 行, 排行榜层 4573 行. 历史层 全部的 Average 均无法由六项任务得分的平均值在 $0.02$ 分容差内重建, 其中 23 行还包含未经证实的零值. 排行榜层4561 行与 C1 在舍入精度范围内一致, 其余 12 行存在版本或数值不一致. 因此, C3 历史层不能作为与六任务主口径直接可比的独立趋势数据, 排行榜层中的重复记录也不作为新增样本使用.

\begin{definition}[随机猜测基线归一化与任务宏平均]\label{def:q4-task-normalization}
对选项数 $K>1$ 的子任务原始得分比例 $r\in[0,1]$，定义
\begin{equation}\label{eq:q4-task-normalization}
T_K(r):=100\max\left\{0,\frac{r-1/K}{1-1/K}\right\}.
\end{equation}
有 $m$ 个子任务的任务分数为 $m^{-1}\sum_{j=1}^mT_{K_j}(r_j)$；先归一化再宏平均. 综合能力仍是 C1 的六任务完整均分 $S_i=\frac16\sum_{j=1}^6S_{ij}$. 
\end{definition}

公式~\eqref{eq:q4-task-normalization} 与排行榜官方口径一致\cite{hf_normalization}. 对于选择题任务, BBH 和 MuSR 按各子任务实际选项数分别确定随机基线, GPQA 的随机基线为 $1/4$, MMLU-PRO 为 $1/10$. IFEval 取 prompt 和 instance 两项严格准确率的均值, MATH 则保留对应版本评分器给出的结果. 当不同子任务的选项数 $K_j$ 不同时, 先按各自随机基线完成归一化, 再计算任务宏平均.


\begin{table}[H]\centering\small
\caption{C8 与 C1 同身份任务评分口径归一化前后的 RMSE 对比}\label{tab:q4-task-alignment}
\begin{tabular}{lrr}\toprule
任务 & 原始比例直接换分 RMSE & 按任务口径归一化后 RMSE\\\midrule
IFEval & 2.499956 & 2.499956\\
BBH & 21.385005 & 0.460472\\
MATH & 10.783909 & 10.783909\\
GPQA & 23.465973 & 0.805803\\
MuSR & 30.793565 & 0.795020\\
MMLU-PRO & 7.694422 & 0.649042\\\bottomrule
\end{tabular}\end{table}
\begin{figure}[H]\centering
\includegraphics[width=0.92\linewidth]{q4_v26_task_alignment.pdf}
\caption{C8 与 C1 同身份任务评分口径归一化前后的 RMSE 对比图 }\label{fig:q4-v26-tasks}
\end{figure}

表~\ref{tab:q4-task-alignment} 和图~\ref{fig:q4-v26-tasks} 给出 C8 与 C1 的任务对齐结果. C8 共包含 1958 个版本文件, 其中 1954 个可解析, 4 个无效. 据此重建 76174 条子任务记录和 11724 条任务汇总记录, 并保留指标、任务版本、shot 与评测配置等信息. 历史起点前共有 1859 个模型版本, 其中 715 个可与许可证和模型类型均符合要求的 C1 记录完成身份匹配, 严格身份子集包含 10 个模型.在严格身份子集中, IFEval、BBH、GPQA、MuSR 和 MMLU-PRO 经任务口径归一化后与 C1 一致到浮点精度, 而 MATH 的 RMSE 仍为 $8.517983$ 分, 仅 3 个模型的差异控制在 $0.02$ 分以内. 官方于 2025-02-14 更新过 MATH 评分器\cite{hf_mathverify}, 这一事实说明评测版本需要单独控制, 但不足以据此判断每条 MATH 残差的具体成因.

\subsection{规模、时间与 Loss--Benchmark 接口}

完成能力口径重建后，先回答历史能力如何随规模和时间变化，再检验训练 Loss 是否提供超出简单规模基线的额外信息. 二者分别承担观察性演化解释和跨指标诊断. 

\subsubsection{规模与时间残差的观察性分解}

在较宽的 1708 条 C1 有效样本上，令 $x_i:=\log_{10}N_i$，并以 C1 提交日期 $T_i$ 定义标准化时间 $t_i:=(T_i-T_0)/365.2425$，拟合包含模型类型效应及规模与时间交互效应的观察性关系
\begin{equation}\label{eq:q4-account-v8}
S_i=\theta_0+\theta_Nx_i+\theta_tt_i+\theta_{Nt}x_it_i
 +\delta_P\mathbf 1_{\{G_i=P\}}+\delta_M\mathbf 1_{\{G_i=M\}}+\varepsilon_i.
\end{equation}
其中 $G_i=P$ 表示第 $i$ 个模型属于预训练模型，$G_i=M$ 表示属于 base merge. C1 的 \texttt{Pretrained}、\texttt{Base merges} 与 \texttt{Chat models/Fine-tuned} 在程序中依次映为 pre、merge、post. 该模型类型划分仅用于 1708 条样本的观察性分解. 在后续严格前沿的联合分位模型中, 仅区分预训练模型与后训练模型.

\setcounter{obsitem}{0}
\obsitem{观察性动力学表示}
对固定类型 $G$，若历史规模路径写成 $x=x(t)$，则拟合能力 $\widehat S(t)=f(x(t),t,G)$ 满足链式法则
\begin{equation}\label{eq:q4-dynamics-v28}
\frac{\dd\widehat S}{\dd t}
=
\underbrace{\left(\theta_N+\theta_{Nt}t\right)\frac{\dd x}{\dd t}}_{\text{规模扩张通道}}
+
\underbrace{\left(\theta_t+\theta_{Nt}x(t)\right)}_{\text{非规模时间残差通道}}.
\end{equation}
式~\eqref{eq:q4-dynamics-v28} 给出了本问采用的观察性动力学模型：能力变化率被分成沿规模方向移动产生的变化和固定规模后的时间残差. 后者作为\textbf{非规模技术进步}的统计代理，但仍可能混合\textbf{数据工程、架构、对齐、评测口径与样本构成变化}，因此不作因果解释.

\obsitem{观察性分解恒等式}
固定模型类型 $G$，记式~\eqref{eq:q4-account-v8} 的拟合值为 $f(x,t,G)$. 首尾窗口以类型标准化的规模中位数与时间为参照，定义两因素对称分解
\begin{subequations}
    \begin{align}
\Delta_{\rm scale}&=\tfrac12\{f(x_1,t_0,G)-f(x_0,t_0,G)+f(x_1,t_1,G)-f(x_0,t_1,G)\},\\
\Delta_{\rm time}&=\tfrac12\{f(x_0,t_1,G)-f(x_0,t_0,G)+f(x_1,t_1,G)-f(x_1,t_0,G)\}.
\end{align}
\end{subequations}
对上述双线性 $f$ 展开，有
\begin{subequations}\label{eq:q4-shapley-v9}
    \begin{align}
\Delta_{\rm scale}&=(x_1-x_0)\left[\theta_N+\tfrac12\theta_{Nt}(t_0+t_1)\right],\\
\Delta_{\rm time}&=(t_1-t_0)\left[\theta_t+\tfrac12\theta_{Nt}(x_0+x_1)\right],\\
\Delta_{\rm scale}+\Delta_{\rm time}&=f(x_1,t_1,G)-f(x_0,t_0,G).
\end{align}
\end{subequations}


把公式~\eqref{eq:q4-account-v8} 的双线性表达代入四个端点的 $f$ 值，常数项与固定类型项在差分中抵消，分别收集 $x_1-x_0$ 和 $t_1-t_0$ 得前两式，再相加消去交互项的路径差，得到最后一式. 
\par
公式~\eqref{eq:q4-shapley-v9} 的意义是把双线性模型从起点到终点的总拟合变化，对称地分配给“规模改变”和“时间改变”两条路径，因此二者严格相加为总变化. 它解决的是交互项导致的\textbf{路径依赖记账问题}，而不是因果识别问题. 上述分解在固定类型 $G$、固定拟合系数下是精确恒等式，窗口间的样本构成变化与遗漏变量仍进入拟合项，故 $\Delta_{\rm time}$ 仅刻画控制规模因素后的观察性时间残差，不是可识别的算法因果贡献. 

\obsitem{观察性证据} 90 日窗口的拟合变化可分为规模分量 $-2.250630$ 分与时间残差 $+7.832644$ 分，总和为 $+5.582014$ 分. 把窗口长度改为 30、60、120 日时，规模分量依次为 $-1.3139$、$-1.5992$、$-2.4541$ 分. 加入发布组织固定效应后，90 日规模分量为 $-2.1614$ 分. 

按题目要求把两部分进一步写成“贡献占比”. 若以净拟合增量 $\Delta_{\mathrm{tot}}=5.582014$ 为分母，则有符号记账份额为
\[
\omega_{\mathrm{scale}}=\frac{-2.250630}{5.582014}=-40.3\%,\quad
\omega_{\mathrm{time}}=\frac{7.832644}{5.582014}=140.3\%,
\]
二者之和为 $100\%$. 超过 $100\%$ 与负份额共同表示两个分量方向相反. 主答案采用上述有符号净变化份额：规模扩张为 $-40.3\%$，非规模时间残差为 $140.3\%$，对应规模分量抵消了部分时间增益. 若只描述两项变动的绝对幅度，占比分别为 $22.3\%$ 与 $77.7\%$. 这不是净能力提升的正向贡献拆分. 两种口径均是式~\eqref{eq:q4-shapley-v9} 的观察性记账量，不具有“技术进步占比”的含义. 


\begin{figure}[H]\centering
\includegraphics[width=0.94\linewidth]{q4_observational_decomposition.pdf}
\caption{能力变化的两因素分解 }\label{fig:q4-decomposition-v11}
\end{figure}
图~\ref{fig:q4-decomposition-v11} 的横轴为首尾窗口长度, 柱分别表示规模分量与时间残差, 黑点连线表示二者合计. 右图加入发布组织固定效应, 并列比较四种窗口长度和两种回归口径. 规模分量与时间残差之和满足式~\eqref{eq:q4-shapley-v9}, 其中时间残差仅作观察性解释, 不表示可识别的因果技术贡献.
\subsubsection{同口径 Loss--Benchmark 桥接}

前三问以交叉熵 Loss 为主要评价指标, 本问则以 Benchmark 分数衡量模型能力. 若模型族或评测口径不同, Loss 与 Benchmark 分数之间就不能直接建立换算关系. 因此, 桥接分析仅使用 C6 中与 B1/Pythia 具有可比性的记录, 检验 Loss 是否能够提供超出简单规模基线的额外能力信息. 

第二问基于 B1 真实 Pythia 轨迹得到的参数用于 C6 的高可比桥接, B6 的半合成 Loss 不进入能力回归. 高可比层包含 7 个同模型族样本, 采用逐模型留一验证, 并与简单基线使用相同的数据划分, 结果见表~\ref{tab:q4-bridge-v8} 中可比层包含 68 条跨验证集记录, 按发布组织整组留出进行检验, 其结果仅在该层内部比较, 不与高可比层合并排名.
\begin{table}[H]\centering\small
\caption{C6 高可比七点的同折留一桥接检验}\label{tab:q4-bridge-v8}
\begin{tabular}{lrr}\toprule
模型 & RMSE（分） & MAE（分）\\\midrule
训练集均值 & 0.424448 & 0.375562\\
$\log_{10}N$ 线性 & 0.449642 & 0.397018\\
Loss 线性 & 0.552834 & 0.478108\\
Loss 岭回归 & 0.439516 & 0.414016\\
\bottomrule\end{tabular}\end{table}

高可比层仅包含 7 个同族模型. 留一检验中, 训练集均值基线的 RMSE 为 $0.424448$, 低于 Loss 线性的 $0.552834$ 和 Loss 岭回归的 $0.439516$. 以均值基线为参照, $\log_{10}N$ 线性、Loss 岭回归和 Loss 线性的 RMSE 分别高出 $5.94\%$、$3.55\%$ 和 $30.25\%$, 三种候选均未优于均值基线, 因此该层没有显示 Loss 能够提供额外的能力预测信息.

中可比层按发布组织整组留出, $\log_{10}N$ 线性、Loss 线性、Loss 岭回归和均值基线的 RMSE 分别为 $10.485$、$11.752$、$11.558$ 和 $11.984$ 分. 相对于均值基线, 四者前三者的 RMSE 分别\textbf{降低 $12.51\%$、$1.94\%$ 和 $3.56\%$}. 因此, Loss 线性和 Loss 岭回归在该层均优于均值基线, 但仍弱于 $\log_{10}N$ 线性模型. 由于高、中可比层的 Loss 标尺和验证单位不同, 两层结果不能直接横向比较, 中可比层仅作为独立诊断, 不用于判断同族 Loss 与能力之间是否存在稳定映射.

基于同族同标尺的高可比层结果, 本文在第四问预测中将 Loss 通道权重设为零. 中可比层得到的不同方向结果表明, 这一处理依赖数据的可比层级. B1 在 $N=1$B、$D=150$B 时的 Loss $2.348910$ 仅作为同口径 Pythia 附近的条件插值示例, 不进入能力预测. 因此, 第三问的最优 Loss 与第四问的能力分数\textbf{之间不建立直接映射}, 二者仅通过训练算力 $6ND$ 联系. 
\subsection{全域有序分位前沿与全局可解性}
 
在统一能力口径并完成规模--时间分析与 Loss--Benchmark 桥接后, 本文进一步构造能力前沿. 对给定的模型规模、时间和类型, 以前景条件分布的指定分位数刻画相应能力水平. 本文采用背景定义~\ref{def:background-quantile} 中的条件 $0.85$ 分位, 并以命题~\ref{thm:bg-pinball-quantile} 对应的平均 Pinball 损失评价预测效果. 

\subsubsection{原分数尺度的联合有序模型}
历史起点固定为 $T_0=$2025-03-13，原协变量为 $\boldsymbol x=(\log_{10}n,t)^\top$、$t=(T-T_0)/365.2425$. 辅助训练集 $\mathcal A$ 中每个协变量的组织加权均值与标准差分别为 $\mu_{A,j}, \sigma_{A,j}$，标准化变量为 $z_j=(x_j-\mu_{A,j})/\sigma_{A,j}$. 数值近零尺度取 $1$, 严格集的未来标签不参与标准化. 

\begin{definition}[组织权重与联合分位模型]\label{def:q4-joint-quantile}
令 $o(i)$ 为发布组织，$n_o$ 为同一来源中组织 $o$ 的记录数，$m_o$ 为该组织的抽样频数. 来源内权重为
\[
w_i:=\frac{m_{o(i)}/n_{o(i)}}{\sum_r m_{o(r)}/n_{o(r)}}.
\]
普通拟合取 $m_o=1$. 设类型 $G\in\{\mathrm{pre},\mathrm{post}\}$，$\tau_j\in\{0.05,0.85,0.95\}$，定义共享斜率的预测曲线
\begin{equation}\label{eq:q4-joint-quantile}
\widehat q_{\tau_j}(\boldsymbol x,G)
:=100(\boldsymbol z^\top\boldsymbol\theta+b_{G,j}).
\end{equation}
\end{definition}
辅助集有 1222 个仓库、350 个组织，与严格训练仓库互斥；整组织留出时两个来源同时排除该组织. 以 $y_i=S_i/100$，辅助回归先求
\begin{equation}\label{eq:q4-auxiliary-ridge}
\min_{a_A,\boldsymbol\theta_A,\zeta_A}
\sum_{i\in\mathcal A}w_i
\{y_i-a_A-\boldsymbol z_i^\top\boldsymbol\theta_A
-\zeta_A\mathbf1_{\{G_i=\mathrm{post}\}}\}^2
+0.01(\|\boldsymbol\theta_A\|_2^2+\zeta_A^2).
\end{equation}
截距 $a_A$ 不参与惩罚, 这里只保留斜率作为严格分位拟合的收缩中心. 严格训练集 $\mathcal T$ 的目标为
\begin{equation}\label{eq:q4-joint-objective}
\min_{\boldsymbol\theta,\boldsymbol b}
\frac13\sum_{i\in\mathcal T}w_i\sum_{j=1}^3
\rho_{\tau_j}(y_i-\boldsymbol z_i^\top\boldsymbol\theta-b_{G_i,j})
+\lambda_{\mathrm{QR}}\|\boldsymbol\theta-\boldsymbol\theta_A\|_2^2,
\quad\lambda_{\mathrm{QR}}=1.
\end{equation}
这样训练目标和原分数尺度的分位评价一致. 先以平方损失估计趋势再加残差分位，一般不等价于式~\eqref{eq:q4-joint-objective}. 训练组织等权也不意味着验证时自动使用组织等权，后文同时给出记录平均与组织平均诊断. 

\begin{theorem}[有序联合分位规划的全域形状与全局可解性]\label{thm:q4-ordered-quantile-qp}
声明形状域为矩形
\[
\mathcal D=[-2,4]\times[-2,2],
\]
其顶点集合记为 $\mathcal V(\mathcal D)$. 对每一类型 $G\in\{\mathrm{pre},\mathrm{post}\}$ 与
$0<\tau_1<\tau_2<\tau_3<1$，考虑共享斜率的仿射分位函数
\[
\widehat q_{\tau_j}(\boldsymbol x,G)
=100\{\boldsymbol z(\boldsymbol x)^\top\boldsymbol\theta+b_{G,j}\},
\quad j=1,2,3,
\]
其中 $\boldsymbol z(\boldsymbol x)$ 是由式~\eqref{eq:q4-joint-quantile} 前定义的仿射标准化映射. 设
$w_i\ge0$、$\sum_iw_i=1$、$\lambda_{\mathrm{QR}}>0$，且所有训练输入位于 $\mathcal D$. 若对每个
$v\in\mathcal V(\mathcal D)$ 与每一类型 $G$ 施加
\begin{equation}\label{eq:q4-vertex-order}
0\le\widehat q_{\tau_1}(v,G)
\le\widehat q_{\tau_2}(v,G)
\le\widehat q_{\tau_3}(v,G)\le100,
\end{equation}
并假设 $y_i\in[0,1]$，系数盒有界且包含 $\boldsymbol\theta=0$、全部 $b_{G,j}=1/2$ 的常数预测器，则有：
\begin{enumerate}
\renewcommand{\labelenumi}{(\roman{enumi})}

\item 式~\eqref{eq:q4-vertex-order} 在整个 $\mathcal D$ 上成立；特别地，共享斜率时全域非交叉等价于
$b_{G,1}\le b_{G,2}\le b_{G,3}$，上下界只需在顶点检验. 
\item 以 Pinball 损失的上图变量 $r_{ij}$ 施加
\[
r_{ij}\ge\tau_j(y_i-\widehat q_{ij}/100),\quad
r_{ij}\ge(\tau_j-1)(y_i-\widehat q_{ij}/100),
\]
并取 $0\le r_{ij}\le1$，则式~\eqref{eq:q4-joint-objective} 与一个非空紧可行域上的凸二次规划等价，因而至少存在一个全局最优解. 
\item 所有全局最优解具有相同的共享斜率 $\boldsymbol\theta^\star$；类型--分位截距
$\boldsymbol b$ 则不一定唯一. 
\item 令全部线性约束写为 $\mathbf D\boldsymbol\omega\ge\boldsymbol a$，盒约束写为
$\boldsymbol l\le\boldsymbol\omega\le\boldsymbol u$，二次规划目标记为 $f$. 对任意盒内点
$\boldsymbol\omega$ 和任意 $\boldsymbol\mu\ge0$，设$\boldsymbol r=\nabla f(\boldsymbol\omega)-\mathbf D^\top\boldsymbol\mu,$则
\begin{equation}\label{eq:q4-qp-lower}
\begin{aligned}
\mathrm{LB}:={}&f(\boldsymbol\omega)
-\boldsymbol\mu^\top(\mathbf D\boldsymbol\omega-\boldsymbol a)\\
&+\min_{\boldsymbol l-\boldsymbol\omega\le\boldsymbol v\le
\boldsymbol u-\boldsymbol\omega}
\{\boldsymbol r^\top\boldsymbol v+\lambda_{\mathrm{QR}}\|\boldsymbol v_\theta\|_2^2\}
\le f^*.
\end{aligned}
\end{equation}
因此任一可行点给出的 $\mathrm{UB}$ 满足
$0\le f(\boldsymbol\omega_{\rm feas})-f^*\le\mathrm{UB}-\mathrm{LB}$，从而得到可计算的全局目标误差证书. 
\end{enumerate}

\end{theorem}
\begin{proof}
矩形 $\mathcal D$ 是其四个顶点的凸包，$\boldsymbol z(\boldsymbol x)$ 为仿射映射，故每个
$\widehat q_{\tau_j}$、相邻分位差以及 $100-\widehat q_{\tau_3}$ 都是 $\boldsymbol x$ 的仿射函数. 由背景命题~\ref{prop:bg-affine-vertices-v26}，顶点上的非负约束保持到整个凸包；共享斜率在相邻分位相减时消去，得到截距顺序，证明 (i). 

Pinball 损失是两条仿射函数的最大值；因权重非负，上图变量总可取为对应损失而不增加目标，即使零权重处的最优上图变量不唯一. 由于训练输入属于 $\mathcal D$，(i) 给出 $0\le\widehat q_{ij}/100\le1$；又有 $0\le y_i\le1$，故 Pinball 损失不超过 $1$，约束 $0\le r_{ij}\le1$ 不排除原问题的最优解. 目标由线性上图损失与半正定二次项组成，其余约束均为线性或盒约束，因此得到凸二次规划. 取系数盒所包含的 $\boldsymbol\theta=0$、全部 $b_{G,j}=1/2$ 的常数预测，并取对应 Pinball 损失为上图变量，即可构造可行点；有限盒使可行域紧，连续目标因而取得全局最小值，证明 (ii). 

若存在两个全局最优解的共享斜率分别为 $\boldsymbol\theta_1\ne\boldsymbol\theta_2$，则它们的任意真凸组合仍可行. Pinball 损失的上图部分为凸函数，而
$\lambda_{\mathrm{QR}}\|\boldsymbol\theta-\boldsymbol\theta_A\|_2^2$ 在 $\boldsymbol\theta$ 上严格凸，因此该真凸组合的目标严格小于两个最优值的同一凸组合，矛盾. 故 $\boldsymbol\theta^\star$ 唯一；二次项不惩罚全部截距，故不能据此推出 $\boldsymbol b$ 唯一，证明 (iii). 

对 $\boldsymbol\mu\ge0$，拉格朗日函数
$L(\boldsymbol w,\boldsymbol\mu)=f(\boldsymbol w)-\boldsymbol\mu^\top(\mathbf D\boldsymbol w-\boldsymbol a)$
在任一原问题可行点上不超过 $f$，故盒内 $L$ 的最小值是 $f^*$ 的下界. 令
$\boldsymbol w=\boldsymbol\omega+\boldsymbol v$；由于 $f$ 只有共享斜率块含二次项，精确展开为
\[
f(\boldsymbol\omega+\boldsymbol v)
=f(\boldsymbol\omega)+\nabla f(\boldsymbol\omega)^\top\boldsymbol v
+\lambda_{\mathrm{QR}}\|\boldsymbol v_\theta\|_2^2.
\]
代入拉格朗日函数并在盒内对 $\boldsymbol v$ 最小化即得式~\eqref{eq:q4-qp-lower}. 任意原问题可行目标都是上界，因此 $\mathrm{UB}-\mathrm{LB}$ 给出全局目标误差界，证明 (iv). 
\end{proof}
非交叉分位约束的思想见文献~\cite{bondell_noncrossing}. 这里的矩形只是形状保证域，不等同于经验支持域. 共享斜率还意味着同类型分位带宽与协变量无关，是需要由样本外误差检验的结构限制. 

\subsubsection{声明域与时间外推}
程序在原协变量 $(\log_{10}n,t)$ 上取
$\mathcal D=[-2,4]\times[-2,2]$ 的顶点，再用辅助集的中心和尺度
将顶点变换为 $\boldsymbol z$, 因此 $\mathcal D$ 不是标准化坐标的矩形. 
严格训练样本的原坐标范围为
$\log_{10}n\in[-0.3063,2.1481]$、
$t\in[-0.7611,-0.0055]$，均位于 $\mathcal D$ 内。
以 2025-03-13 为历史起点, 12 月和 24 月情景对应的时间坐标分别为
$t=1,2$，所报名义前沿情景的
$\log_{10}n\in[0.9835,2.8883]$,这些点也都在 $\mathcal D$ 内.
标准化时间坐标 $z_t=5.98, 10.34$ 虽大于 $2$，
却不能据此判定预测点超出原坐标声明域, 
因此定理~\ref{thm:q4-ordered-quantile-qp} 的形状结论覆盖这些名义预测点.

\begin{corollary}[全局目标间隙到共享斜率误差]\label{cor:q4-slope-gap-v37}
在定理~\ref{thm:q4-ordered-quantile-qp} 的条件下，令 $\widehat{\boldsymbol\omega}$ 为可行近似解，$\boldsymbol\theta^\star$ 为唯一最优共享斜率. 若可靠上下界给出
\[
0\le f(\widehat{\boldsymbol\omega})-f^*\le\varepsilon_{\mathrm{opt}}:=\mathrm{UB}-\mathrm{LB},
\]
则
\begin{equation}\label{eq:q4-slope-gap-v37}
\|\widehat{\boldsymbol\theta}-\boldsymbol\theta^\star\|_2
\le\sqrt{\varepsilon_{\mathrm{opt}}/\lambda_{\mathrm{QR}}}.
\end{equation}
对于同一类型、同一分位水平的两个输入，记其标准化差为 $\Delta\boldsymbol z$，则预测差的数值误差满足
\begin{equation}\label{eq:q4-contrast-gap-v37}
|\Delta\widehat q-\Delta q^*|
\le100\|\Delta\boldsymbol z\|_2\sqrt{\varepsilon_{\mathrm{opt}}/\lambda_{\mathrm{QR}}}.
\end{equation}
\end{corollary}
\begin{proof}
取任一最优解 $\boldsymbol\omega^*$. 对 $0<t<1$，可行域凸与目标斜率块的二次结构给出
\begin{align*}
f^*&\le f((1-t)\boldsymbol\omega^*+t\widehat{\boldsymbol\omega})\\
&\le(1-t)f^*+tf(\widehat{\boldsymbol\omega})
-\lambda_{\mathrm{QR}}t(1-t)\|\widehat{\boldsymbol\theta}-\boldsymbol\theta^\star\|_2^2.
\end{align*}
移项、除以 $t$ 并令 $t\to0^+$，得
$f(\widehat{\boldsymbol\omega})-f^*\ge\lambda_{\mathrm{QR}}\|\widehat{\boldsymbol\theta}-\boldsymbol\theta^\star\|_2^2$，从而得到第一式. 同类型同分位相减消去截距，再由 Cauchy--Schwarz 不等式得到第二式. 
\end{proof}

对包含全量拟合在内的 279 份有序联合分位拟合记录逐份检查, 全局目标的最大已记录间隙为 $6.70588\times10^{-8}$. 取已报告的 $\lambda_{\mathrm{QR}}=1$, 由各次目标间隙得到的最大归一化斜率数值误差界为 $\sqrt{6.70588\times10^{-8}}=2.59\times10^{-4}$. 声明域内同类型、同分位两个输入之间的最大标准化距离不超过矩形对角长度 $\sqrt{6^2+4^2}=7.2111$. 因此, 式~\eqref{eq:q4-contrast-gap-v37} 给出的最大预测对比误差界为 $100\times7.2111\times2.59\times10^{-4}=0.1867$ 分. 严格训练样本自身的最大标准化距离为 $5.49$, 对应的误差界为 $0.1422$ 分.

\textbf{上述结果表明, 在本文声明域内, 有序联合分位模型的数值优化误差受到较小的上界控制. }

算法~\ref{alg:q4-ordered-frontier} 将顶点形状约束、凸规划的目标界、历史选型和第三问训练 FLOPs 接口串成可复核的第四问主线。历史截点只使用当时可得的记录, 复用晚期与整组织留出仅用于报告适用性，不反过来选择主候选. 未来曲线的组织重抽样范围与个体分位带也分开记录.
\begin{algorithm}[H]
\caption{第四问的有序分位前沿拟合与条件情景预测}\label{alg:q4-ordered-frontier}
\small
\begin{algorithmic}[1]
\Require 六任务 C1/C4 历史记录、第三问求解器、历史起点 $T_0$ 与拟合容差 $\varepsilon$
\Ensure 选中前沿、拟合目标界、回溯指标及 12/24 月条件分位预测
\State $\mathcal M\gets\{\text{联合时间},\text{联合规模--时间}\}$；原坐标域 $\mathcal D\gets[-2,4]\times[-2,2]$
\For{历史截点 $s$ 及候选 $M\in\mathcal M$}
  \State 按 $s$ 建立严格集 $\mathcal T_s$ 与互斥辅助集 $\mathcal A_s$，计算组织权重
  \State 仅由 $\mathcal A_s$ 求标准化中心、尺度与辅助岭回归斜率 $\boldsymbol\theta_A$
  \State 建立共享斜率和类型--分位截距 $\widehat q_{\tau_j}=100(\boldsymbol z^\top\boldsymbol\theta+b_{G,j})$
  \State 在 $\mathcal V(\mathcal D)$ 施加 $0\le\widehat q_{.05}\le\widehat q_{.85}\le\widehat q_{.95}\le100$
  \State 引入 Pinball 上图变量 $r_{ij}$，按式~\eqref{eq:q4-joint-objective} 求凸二次规划
  \State 由可行解得 $\mathrm{UB}$，用非负乘子与式~\eqref{eq:q4-qp-lower} 求 $\mathrm{LB}$ \Comment{定理~\ref{thm:q4-ordered-quantile-qp}}
  \If{约束违反或 $\mathrm{UB}-\mathrm{LB}>\varepsilon$} \State 细化求解；仍未达标则记录失败 \EndIf
  \State 只用截点前 C4 推算验证规模，记录 $0.85$ 分位的开发期 Pinball 损失
\EndFor
\State 以开发期损失选 $M^*$；复用晚期、整组织留出与时间切分校准另作诊断
\State 在 $T_0$ 前全部严格记录上重拟合 $M^*$，保存形状松弛量与 $\mathrm{UB}-\mathrm{LB}$
\For{每种质量成本、$\zeta\in\{0,0.5,1\}$ 及 $H\in\{1,2\}$ 年}
  \State $F_H\gets F_0 10^{\zeta r_C H}$；重解第三问，使 $6N_H D_H=F_H$ \Comment{只传训练 FLOPs}
  \State 核对第三问的全局质量最优性、预算残差与训练算力匹配残差
  \For{$G\in\{\mathrm{pre},\mathrm{post}\}$}
    \State $\boldsymbol x_H\gets(\log_{10}[N_{\mathrm{ref},G}N_H/N_0],H)$；核对 $\boldsymbol x_H\in\mathcal D$
    \State 输出 $\widehat q_{.05}(\boldsymbol x_H,G),\widehat q_{.85}(\boldsymbol x_H,G),\widehat q_{.95}(\boldsymbol x_H,G)$
  \EndFor
\EndFor
\State 同步重抽样发布组织并配对第二问完整参数行，报告有效抽样的曲线范围
\State \Return $M^*$、$[\mathrm{LB},\mathrm{UB}]$、回溯指标、条件预测与曲线抽样范围
\end{algorithmic}
\end{algorithm}

\subsubsection{斜率惩罚强度的敏感性}
前述零时间斜率下的点预测主要由规模倍率驱动，因此进一步检验该结果对惩罚强度选择的敏感性. 保持辅助集、严格样本及其余设定不变，对 $\lambda_{\mathrm{QR}}\in \left\{ 0.01,0.05,0.1,0.25,0.5,1,2,5,10 \right\} $ 分别重新求解同一凸二次规划（$\lambda_{\mathrm{QR}}=0$ 不满足定理~\ref{thm:q4-ordered-quantile-qp} 对正斜率惩罚的要求），因此不纳入比较. 九档结果列于附录~\ref{app:v39-algorithms} 的表~\ref{tab:q4-lambda-sensitivity}.

九组拟合的共享时间斜率绝对值均不超过 $0.105$ 分/年, 仅在 $\lambda_{\mathrm{QR}}=0.25$ 时取 $0.1048$ 分/年, 说明主模型得到的近零时间斜率并非由单一惩罚强度决定. 随着惩罚强度变化，规模斜率由 $1.85$ 增至 $5.25$ 分/十倍. 相应地，在 $\zeta=1$ 的 24 月情景下，有序联合分位候选对预训练模型的点预测范围为 $22.641$--$23.133$ 分，跨度 $0.49$ 分；后训练模型为 $32.277$--$33.263$ 分，跨度 $0.99$ 分. 因此，附录表~\ref{tab:q4-lambda-sensitivity} 中各结构候选的点预测均对应主选设定 $\lambda_{\mathrm{QR}}=1$，其余惩罚强度结果作为敏感性分析报告.
%预训练模型的点预测范围为 $22.641$--$23.133$ 分，跨度 $0.49$ 分；后训练模型为 $32.277$--$33.263$ 分，跨度 $0.99$ 分. 因此，表~\ref{tab} 中的名义点预测均对应主选设定 $\lambda_{\mathrm{QR}}=1$，其余惩罚强度结果作为敏感性分析报告.

作为补充，归档的滚动检验比较了 $\lambda_{\mathrm{QR}}=0.1$ 与 $10$ 的预测表现. 两者开发期平均 Pinball 损失分别为 $4.726347$ 和 $4.670423$，复用晚期分别为 $1.210677$ 和 $0.906569$. 同期 $90\%$ 区间分数分别为 $30.318614$ 和 $45.90053$. 点预测指标与区间指标对惩罚强度呈现不同响应，因此本文仍按照开发期点预测损失确定主候选，保留 $\lambda_{\mathrm{QR}}=1$，并将上述结果作为不同评价指标之间的敏感性对照.


\subsection{求解误差、历史覆盖与有限样本校准}

\subsubsection{开发、复用晚期与跨组织检验}
时间外检验采用 2024-09-13 至 2025-02-13 的六个月度截点, 预测窗口设为 1 个或 3 个日历月, 共形成 9 个有效验证折. 其中, 开发期截至 2024-12-13, 包含 16 次预测、8 个仓库和 2 个组织. 晚期样本外验证包含 8 次预测、4 个仓库和 4 个组织, 且均为后训练模型. 另将跨越开发期与晚期分界的 7 次预测单独报告. 各验证折的主预测仅使用预测时点可获得的日期、模型类型以及当时 C4 信息和参考规模推算得到的 $N$, 不引入未来实际规模. 此外, 本文设置整组织留出检验：将目标组织同时从严格样本和辅助样本中剔除, 并保留已知的 $N$ 作为条件输入, 从而检验模型在给定规模等条件下向新组织迁移的适用性. 由此形成时间外、跨分界与组织外相互补充的历史验证体系.

候选模型限定为联合时间和联合规模--时间两种，并依据开发期损失选择后者，固定 $\lambda_{\mathrm{QR}}=1$. 近半年经验分位、独立分位回归和辅助趋势模型作为比较基准，模型选择仅依据开发期损失，晚期结果不参与选择. 验证指标为原 $0$--$100$ 分尺度上的平均 $0.85$ Pinball 损失
\begin{equation}\label{eq:q4-validation-pinball-v26}
\mathrm{PB}_{0.85}:=\frac1m\sum_{i=1}^m
\rho_{0.85}(S_i-\widehat q_{0.85,i}).
\end{equation}

独立分位基准在开发期和整组织留出中分别出现 15 次和 5 次相邻分位顺序不一致, 联合有序模型在两类检验中均保持分位有序. 在惩罚强度敏感性检验中, $\lambda_{\mathrm{QR}}=10$ 时的晚期损失降为 $0.906569$, 对应区间分数上升为 $45.90053$. 固定主候选的区间分数为 $30.63021$，不根据晚期结果重新选择参数. 此外, 去除 MATH 后的晚期损失为 $1.660028$, 对应不同的五任务目标.

\begin{table}[H]\centering\small
\caption{条件分位模型验证结果}\label{tab:q4-v26-rolling}
\begin{tabular}{lrrrr}\toprule
模型 & 开发期 & 复用晚期 & 跨分界窗口 & 整组织留出\\\midrule
近半年经验分位 & 4.224773 & 2.579856 & 3.389216 & 3.076414\\
独立分位基准 & 5.794666 & 2.448574 & 2.260022 & 1.789583\\
辅助趋势基准 & 3.819872 & 2.720884 & 2.108135 & 1.815348\\
联合时间 & 4.964083 & 1.693356 & 3.409182 & \textbf{4.124014}\\
联合规模--时间 & 4.726347 & 1.189958 & 3.587237 & 3.724345\\\bottomrule
\end{tabular}\end{table}

\begin{table}[H]\centering\small
\caption{各验证条件的样本构成}\label{tab:q4-v26-rolling-conditions}
\begin{tabular}{lrrrll}\toprule
验证条件 & 预测数 & 仓库数 & 组织数 & 窗口 & 规模输入\\\midrule
开发期 & 16 & 8 & 2 & 1 月窗 8 与 3 月窗 8 & 由当时 C4 推算 \\
复用晚期 & 8 & 4 & 4 & 1 月窗 4 与 3 月窗 4 & 由当时 C4 推算 \\
跨分界窗口 & 7 & 5 & 4 & 3 月窗 7 & 由当时 C4 推算 \\
整组织留出 & 33 & 33 & 18 & 不适用 & 已知 $N$ \\\bottomrule
\end{tabular}\end{table}
\begin{figure}[H]\centering
\includegraphics[width=0.94\linewidth]{q4_v26_validation.pdf}
\caption{各条件分位数前沿候选在不同验证条件下的表现. 晚期仅 4 个已查看仓库；整组织留出使用已知规模，与时间窗预测的输入条件不同. }\label{fig:q4-v26-validation}
\end{figure}

表~\ref{tab:q4-v26-rolling} 给出同一六任务目标下的 $0.85$ Pinball 损失, 各列对应不同验证条件. 具体预测数、仓库数与组织数见表~\ref{tab:q4-v26-rolling-conditions}. 前三列按记录等权计算, 末列对应整组织留出结果. 表~\ref{tab:q4-v26-rolling-conditions} 中的“窗口”表示计算相应损失所采用的预测期, 开发期与晚期样本外验证均同时包含 1 月窗和 3 月窗.表~\ref{tab:q4-v26-rolling} 与图~\ref{fig:q4-v26-validation} 显示: 联合规模--时间模型在晚期样本外验证中的 Pinball 损失相对经验分位和辅助趋势基准分别下降 $53.88\%$ 和 $56.27\%$, 表现出较好的时间外预测能力. 进一步将晚期预测窗口限定为 1 个月时, 该模型的损失为 $1.251657$, 仍低于经验分位的 $3.955578$ 和辅助趋势基准的 $2.340656$.作为权重口径的补充检验, 整组织留出采用组织等权后, 联合规模--时间模型、辅助趋势基准和独立分位基准的损失分别为 $2.247946$、$1.552935$ 和 $1.576110$. 因此, 记录等权、组织等权以及不同预测窗口共同构成相互补充的验证口径, 用于从时间外预测和组织外迁移两个方面评价条件分位模型的表现.


\begin{table}[htbp]\centering\small
\caption{$90\%$ 个体分位带的经验覆盖}\label{tab:q4-v39-coverage}
\begin{tabular}{lrrrr}\toprule 角色 & 行数 & 组织数 & 覆盖率 & 平均宽度/分\\\midrule
开发期 & 16 & 2 & 0.500 & 18.37 \\
整组织留出 & 33 & 18 & 0.606 & 17.44 \\
复用晚期 & 8 & 4 & 0.750 & 17.07 \\
跨分界窗口 & 7 & 4 & 0.714 & 18.31 \\
\bottomrule\end{tabular}\end{table}


\subsubsection{时间切分校准及其适用条件}
在每个历史外层截点内, 按仓库日期将较早记录用于模型拟合, 较晚记录用于校准, 同日记录整体划入同一集合. 校准组织同时从两个拟合来源中排除. 模型拟合完成后不再使用校准标签重新估计参数, 所需规模依据更早拟合截点的信息推算.

\begin{definition}[分位带的校准分数与区间分数]\label{def:q4-calibration}
对已固定的上下分位函数，定义
\begin{subequations}
    \begin{align}\label{eq:q4-cqr}
R_i&:=\max\{\widehat q_{0.05}(X_i)-S_i,\ S_i-\widehat q_{0.95}(X_i)\},\\
k_\varepsilon&:=\lceil(m+1)(1-\varepsilon)\rceil,
\quad c_\varepsilon:=\begin{cases}\max(0,R_{(k_\varepsilon)}),&k_\varepsilon\le m,\\+\infty,&k_\varepsilon>m,\end{cases}\\
I(X)&:=[\widehat q_{0.05}(X)-c_\varepsilon,
\widehat q_{0.95}(X)+c_\varepsilon]\cap[0,100],\quad\varepsilon=0.1.
\end{align}
\end{subequations}
若 $c_\varepsilon=+\infty$，约定 $I(X)=[0,100]$. 以
\begin{equation}\label{eq:q4-interval-score}
\mathrm{IS}_{\varepsilon}(l,u;y):=(u-l)
+\frac2\varepsilon(l-y)_++\frac2\varepsilon(y-u)_+
\end{equation}
同时评价宽度与漏覆盖代价，其中 $(v)_+:=\max(v,0)$. 
\end{definition}

\begin{proposition}[可交换条件下的有限样本边际覆盖]\label{thm:q4-cqr-conditional}
条件于拟合信息，若 $m$ 个校准分数与一个新观测分数可交换，且新得分属于 $[0,100]$，则式~\eqref{eq:q4-cqr} 满足
\[
\Prob\{S_{\rm new}\in I(X_{\rm new})\}\ge1-\varepsilon.
\]
该结论是新观测的边际覆盖，不是对每个规模、类型或时间的条件覆盖保证. 
\end{proposition}
\begin{proof}
由背景的可交换秩引理~\ref{thm:bg-exchangeable-rank-v26}，当 $k_\varepsilon\le m$ 时，新分数不超过校准第 $k_\varepsilon$ 个次序统计量的概率至少为 $k_\varepsilon/(m+1)\ge1-\varepsilon$. 将修正量截为非负只会扩大区间，且得分本已位于 $[0,100]$，与该支持集求交不会删除真实得分. 若秩超过样本量，则返回完整支持区间，覆盖必为一. 这是分位回归校准方法的标准秩论证\cite{romano_cqr}. 
\end{proof}
当 $m<9$ 时, 名义 $90\%$ 覆盖所需的经验秩超过现有样本量, 当前样本无法直接支持该分位点, 因此本文不以样本最大值替代缺失秩.  组织级诊断先取各组织内 $R\_i$ 的最大值, 再在组织层面进行排序，相应理论结果以组织单位与未来单位的可交换性为前提. 结合现有时间切分、组织结构、窗口设置与评测版本口径, 统一报告历史回溯中的实测覆盖结果, 并将其作为后续区间表现的经验参考.


\begin{table}[htbp]\centering\small
\caption{复用晚期的区间校准：8 次预测、4 个仓库}\label{tab:q4-v26-calibration}
\begin{tabular}{lrrrr}\toprule
方案 & Pinball 损失 & 实测覆盖 & 平均宽度/分 & 区间分数\\\midrule
全部历史联合拟合 & 1.189958 & 75.0\% & 17.0728 & 30.6302\\
时间切分、未校准 & 4.632594 & 0.0\% & 0.8233 & 146.4512\\
时间切分、记录校准 & 4.632594 & 100.0\% & 26.6454 & 26.6454\\
时间切分、组织校准 & 4.632594 & 100.0\% & 81.6613 & 81.6613\\\bottomrule
\end{tabular}\end{table}
表~\ref{tab:q4-v26-calibration} 汇总覆盖率、区间宽度与区间分数, 相应图示见附录~\ref{app:v39-algorithms} 的图~\ref{fig:q4-v26-calibration}, 其中 10 个组织校准对应完整支持区间. 随着覆盖率提高, 区间宽度相应增加. 各切分模型使用对应历史截点之前的数据拟合, 因此其点预测与全历史样本拟合结果分别计算. 校准仅调整上下分位带, 不改变切分模型的 $0.85$ 分位点预测, 因此后三行的 Pinball 损失保持一致.
\subsubsection{有序前沿的凸组合与开发期定权}

历史检验表明，不同基准在不同验证条件下各有优劣. 保留多个候选时，可以在有序分位函数的凸集中组合，并用明确的预测损失约束其长期代价. 

\begin{proposition}[有序分位候选的凸组合]\label{prop:q4-convex-aggregate-v37}
设第 $k$ 个候选在共同输入域上满足
$0\le q^{(k)}_{\tau_1}\le\cdots\le q^{(k)}_{\tau_J}\le100$. 
若 $a_k\ge0$、$\sum_{k=1}^Ka_k=1$，并对所有分位共用权重，定义
\begin{equation}\label{eq:q4-convex-aggregate-v37}
q^{\rm agg}_{\tau_j}(x,G):=\sum_{k=1}^Ka_kq^{(k)}_{\tau_j}(x,G),
\end{equation}
则组合保持相同的有序性与上下界，且对任意真实分数 $S$ 和任一分位水平，
\[
\rho_{\tau_j}(S-q^{\rm agg}_{\tau_j})
\le\sum_ka_k\rho_{\tau_j}(S-q^{(k)}_{\tau_j}).
\]
\end{proposition}
\begin{proof}
非负加权保持每个相邻分位差及端点界；Pinball 损失关于预测值凸，应用 Jensen 不等式即得损失界. 
\end{proof}

有序性不等于组合已经是真实条件分位数，损失凸性也不表示每次预测都优于所有候选. 标准顺序反馈设定下，还可得到累计比较保证. 
\subsubsection{开发期定权的实际组合}\label{par:q4-aggregation-v38}
以表~\ref{tab:q4-v26-rolling} 的联合规模--时间与联合时间两候选为输入，只在开发期 16 个标签上选择一个对各输出层共用的权重 $\lambda\in[0,1]$. 对 $0.85$ 分位，求解分段线性的凸目标
\begin{equation}\label{eq:q4-aggregation-lp-v38}
\min_{0\le\lambda\le1}\frac1{16}\sum_{i=1}^{16}
\rho_{0.85}\!\left(S_i-\lambda q^{(\mathrm{size,time})}_{0.85,i}
-(1-\lambda)q^{(\mathrm{time})}_{0.85,i}\right).
\end{equation}

线性规划与全部有效残差折点的直接枚举均得到最优权重 $\lambda=1$, 对应开发期平均 Pinball 损失为 $4.726347$, 两种求解方法得到完全一致的权重结果. 将该权重统一用于区间下界、$0.85$ 分位点预测与区间上界后, 在开发期、跨分界窗口、晚期样本外验证及整组织留出的 64 条对齐记录上均满足 $0\le \mathrm{low}\le q_{0.85}\le \mathrm{high}\le100$. 

其中 8 条记录在 $5\times10^{-14}$ 的数值容差内满足 $q_{0.85}=\mathrm{high}$, 表明组合后的分位结构在数值计算中保持有序约束.由于最优权重为 $\lambda=1$, 组合结果与原联合规模--时间候选逐条一致. 四类验证条件下的平均 Pinball 损失分别为 $4.726347$、$3.587237$、$1.189958$ 和 $3.724345$, 与表~\ref{tab:q4-v26-rolling} 中对应候选的结果完全一致. 因此, 权重优化、分位有序约束与已有预测结果在数值上相互吻合, 构成对凸组合求解过程和实现结果的双重一致性核验.

\subsubsection{把有序基准纳入候选集的定权结果}\label{par:q4-aggregation-baseline}


上述凸组合仅包含两个联合分位候选, 而命题~\ref{prop:q4-convex-aggregate-v37} 适用于任意满足有序性约束的候选. 6.4.1 节构造的有序辅助趋势基准同样满足该条件, 因而可进一步纳入统一候选集. 在开发期 16 个预测标签上求解三候选的单纯形权重后, 得到最优权重 $(0,0,1)$, 即开发期数据全部选择有序辅助趋势基准, 对应平均 Pinball 损失为 $3.819872$；同期联合规模--时间候选的损失为 $4.726347$. 这一结果表明, 凸组合框架能够在保持分位有序性的同时, 根据开发期预测表现对不同结构候选进行统一比较与数据驱动定权.

不同验证条件下各候选的表现有所差异. 在复用晚期 8 条记录上, 辅助趋势基准与联合规模--时间候选的平均 Pinball 损失分别为 $2.720884$ 和 $1.189958$. 在跨分界 7 条记录上, 两者分别为 $2.108135$ 和 $3.587237$. 本文继续采用 6.4.1 节预先设定的选型规则, 即在有序联合分位候选内部依据开发期损失选择主候选. 三候选的验证损失与开发期权重见

\begin{table}[H]\centering\footnotesize
\setlength{\tabcolsep}{4pt}
\caption{训练算力增长下的资源倍率}\label{tab:q4-aggregation-baseline}
\begin{tabular}{lrrrr}\toprule
候选 & 开发期 $(n=16)$ & 复用晚期 $(n=8)$ & 跨分界 $(n=7)$ & 开发期最优权重\\\midrule
联合规模--时间（主候选） & $4.726347$ & $1.189958$ & $3.587237$ & $0$\\
联合时间 & $4.964083$ & $1.693356$ & $3.409182$ & $0$\\
有序辅助趋势基准 & $3.819872$ & $2.720884$ & $2.108135$ & $1$\\\midrule
三候选凸组合 & $3.819872$ & $2.720884$ & $2.108135$ & ---\\\bottomrule
\end{tabular}
\end{table} 
表~\ref{tab:q4-aggregation-baseline} 中的权重仅依据开发期 16 个标签拟合, 复用晚期与跨分界结果作为样本外对照. 即时反馈权重更新的累计后悔界及 $T=256$ 的合成协议核验见附录~\ref{app:q4-online-extension}. 
\subsection{训练算力映射与条件未来前沿}

最后把第三问的资源决策接入能力前沿. 跨问只传递可比训练 FLOPs，不把第三问总预算或 B6 Loss 直接当作 Benchmark 解释变量. 

\subsubsection{从历史训练算力到第三问资源路径}
C4 限定截至 2025-03-13 发表的 Language 域、开放权重及正训练 FLOPs 后有 346 条记录，完整训练 FLOPs、$N,D$ 的记录为 247 条. 七个年度训练算力 $90\%$ 分位的对数增长斜率为 $r_C=0.604052$. 令 $y$ 为日历年份，历史观察关系写为
\begin{equation}\label{eq:q4-compute-v8}
\log_{10}C_{\rm train}=a_C+\kappa_C\log_{10}n+\xi_C(y-2020),
\quad \kappa_C=1.174053,\quad\xi_C=0.249006.
\end{equation}
若增长情景为 $\zeta r_C H$ 个 $\log_{10}$ 单位，且关系系数在未来 $H$ 年固定，则起终点相减得到
\begin{equation}\label{eq:q4-empirical-increment}
\Delta\log_{10}n=\frac{(\zeta r_C-\xi_C)H}{\kappa_C}.
\end{equation}
其中 $\zeta\in{0,0.5,1}$ 分别对应零增长、历史增长斜率减半和历史增长斜率延续. 在该经验关系下, 即使训练算力保持不变, 参数规模仍可随年份项变化. 作为另一种结构情景, 若固定 $D$, 由 $C_{\rm train}=6ND$ 可得
$\Delta\log_{10}n=\zeta r_C H$.
两种表达对应不同的条件设定.

第三问中, 每种质量成本均以总预算 $C_0=10^{22}$ FLOPs、上下文长度 8192 及其自身最优配置为锚点. 先计算该配置对应的训练 FLOPs$F_0:=6N_0D_0$,再定义目标训练 FLOPs$F_H:=F_0,10^{\zeta r_C H}$. 这里只匹配训练分量：
\begin{equation}\label{eq:q4-training-match-v8}
6N^*(C_H)D^*(C_H)=F_H.
\end{equation}

\begin{lemma}[饱和分支上训练 FLOPs 匹配根的唯一性]\label{prop:q4-training-root}
在第三问 $Q=1$ 分支，设 $A,B_p,k,\alpha,\beta>0$、$h_1=h(1)\ge0$，且 $g(0)=0$. 给定 $T_H:=F_H/10^{18}>0$，尺度满足
\begin{equation}\label{eq:q4-training-root}
\frac{n^{\alpha+\beta+1}}{kn+h_1}
=\frac{\alpha A}{\beta B_pk}\left(\frac{T_H}{6}\right)^\beta.
\end{equation}
该式具有唯一正根，回代 $d=T_H/(6n)$、$c=d(kn+h_1)$ 即得该约束分支所需总预算. 要成为完整问题的匹配解，还须验证该预算下 $Q=1$ 全局最优. 
\end{lemma}
\textbf{证明说明:} 证明见\ref{app:q4-training-root-proof}. 
由式~\eqref{eq:q4-training-root} 得到 $Q=1$ 约束分支后，仍需回到第三问的完整质量剖面比较以确认其全局最优性.

\begin{table}[H]\centering\small
\caption{指数成本下的资源倍率}\label{tab:q4-v26-resources}
\begin{tabular}{llrrrl}\toprule
质量成本 & 年数 & $N$ 倍率 & $D$ 倍率 & 总预算倍率 & B6 支持域\\\midrule
指数 & 1 & 1.363034 & 2.948119 & 3.950885 & $N,D$ 范围内\\
幂函数 & 1 & 1.355220 & 2.965117 & 3.942384 & $N,D$ 范围内\\
对数 & 1 & 1.406468 & 2.857076 & 3.997592 & $N,D$ 范围内\\
指数 & 2 & 1.877598 & 8.600043 & 15.671419 & $N,D$ 均外推\\
幂函数 & 2 & 1.857838 & 8.691512 & 15.609424 & $N,D$ 均外推\\
对数 & 2 & 1.985810 & 8.131404 & 16.003896 & $N,D$ 均外推\\\bottomrule
\end{tabular}\end{table}
表~\ref{tab:q4-v26-resources} 中均有 $Q^*=1$. 指数成本一年点为 $N=9.679971$B、$D=509.193693$B，两年点为 $N=13.334295$B、$D=1485.38353$B. 两年所有成本的点估计均超出 B6 的 $N,D$ 范围. 一年范围内也不能证明跨来源质量标尺统一. 图~\ref{fig:q4-v26-resources} 展示的是资源倍率，不是能力分数. 
\begin{table}[H]\centering\small
\caption{三档算力增长下的 12/24 月资源倍率}\label{tab:q4-v39-slow-resources}
\begin{tabular}{rrrrr}\toprule
$\zeta$ & 月数 & $N$ 倍率 & $D$ 倍率 & 总预算倍率\\\midrule
0 & 12/24 & 1.000 & 1.000 & 1.000\\
0.5 & 12 & 1.166 & 1.719 & 1.987\\
0.5 & 24 & 1.363 & 2.948 & 3.951\\
1 & 12 & 1.363 & 2.948 & 3.951\\
1 & 24 & 1.878 & 8.600 & 15.671\\\bottomrule
\end{tabular}\end{table}
表~\ref{tab:q4-v39-slow-resources} 给出指数质量成本下三档训练算力增长对应的名义资源倍率, 基准为第三问 $10^{22}$ FLOPs 下的自身最优解. 由于减半增长 24 月与延续增长 12 月具有相同的 $\zeta H$, 两种情景对应的资源点重合. 表中各资源点均按照第三问包含质量与注意力开销的总预算口径计算.
\begin{figure}[H]\centering
\includegraphics[width=0.93\linewidth]{q4_v26_resource_paths.pdf}
\caption{三类质量成本下的训练 FLOPs 匹配资源路径 }\label{fig:q4-v26-resources}
\end{figure}

\subsubsection{主模型条件前沿与联合不确定性}
共同发布组织抽样同时重估辅助趋势、严格分位、参考规模及 C4 关系. 第二问的完整曲率参数来自另一数据来源, 按原抽样编号配对, 以保持参数之间的行内相关性.  200 次联合抽样全部成功, 后续不确定性分析基于这 200 次有效抽样. 分数回归采用训练期的类型参考规模. 第三问则使用机制锚点 $N_0$, 两者对应不同的建模口径, 因此 B6 的半合成 Loss 与 Pythia Benchmark 分别保留在各自模型中, 跨问连接仍通过统一的训练算力口径实现.
% 共同发布组织抽样同时重估辅助趋势、严格分位、参考规模及 C4 关系. 第二问的完整曲率参数来自另一数据来源, 按原抽样编号配对, 以保持参数之间的行内相关性. 200 次联合抽样中共有 199 次有效, 后续不确定性分析基于这些有效抽样结果.分数回归采用训练期的类型参考规模, 第三问则使用机制锚点 $N_0$. 两者对应不同的建模口径, 因此 B6 的半合成 Loss 与 Pythia Benchmark 分别保留在各自模型中, 跨问连接仍通过统一的训练算力口径实现.
\begin{table}[H]\centering\small
   \caption{主模型的 12/24 月能力分数参考值与两类区间}
   \label{tab:q4-v39-slowdown}
   
   \begin{tabular}{lrrrr}\toprule 类型 & 目标日        & 点估计   & 曲线抽样范围          & 个体分位带           \\\midrule
                预训练                & 2026-03-13 & 33.20 & $[15.53,33.74]$ & $[9.54,36.35]$  \\
                预训练                & 2027-03-13 & 33.20 & $[15.53,33.74]$ & $[9.54,36.35]$  \\
                后训练                & 2026-03-13 & 29.44 & $[24.42,30.77]$ & $[14.53,30.33]$ \\
                后训练                & 2027-03-13 & 29.44 & $[24.42,30.77]$ & $[14.53,30.33]$ \\
                \bottomrule\end{tabular}\end{table}

 %\begin{figure}[H]\centering
 %  \includegraphics[width=0.97\linewidth]{q4_slowdown_frontier.pdf}
%\end{figure}
表~\ref{tab:q4-v39-slowdown} 给出主模型在 12/24 月两个期限上的能力点估计与不确定性区间. 历史起点为 2025-03-13, 对应情景日期分别为 2026-03-13 和 2027-03-13, 单位均为综合能力分. 预训练与后训练模型的点估计分别为 33.20 分和 29.44 分. 由于主模型采用近 180 日的类型经验分位, 两个期限的点估计保持一致；三档算力增长情景主要通过表~\ref{tab:q4-v39-slowdown} 中的资源倍率体现. 两类区间刻画不同层面的不确定性：曲线抽样范围反映经验分位估计随样本变化产生的波动, 个体分位带则对应单个模型分数的 $0.05$--$0.95$ 分位范围. 后者未经历史覆盖校准, 因而本文将其作为个体分数分布范围的条件性描述.


主模型采用严格前沿入选口径. 前述 90 日观察性分解使用许可、有效性和去重后的 1708 条较宽样本, 属于另一数据口径. 为检验经验分位对样本选择的敏感性, 表~\ref{tab:q4-baseline-caliber} 并列报告三种入选口径在同一 180 日窗口、同一六任务目标下的结果.


保留的经验分位基准不含时间趋势，其两个期限的参考值因而相同. 由于该基准的数值取决于样本入选口径，而本节的经验分位对照采用严格前沿口径. 前述 90 日观察性分解采用许可、有效性和去重后的 1708 条较宽样本，两者不可视作同一入选集. 本节并列报告三种口径在同一 180 天窗口、同一六任务目标下的结果，见表~\ref{tab:q4-baseline-caliber}. 

\begin{table}[H]\centering\footnotesize
\setlength{\tabcolsep}{4pt}
  \caption{不同样本入选口径下的近 180 日六任务经验分位}\label{tab:q4-baseline-caliber}
  \begin{tabular}{lrrrr}\toprule
    入选口径             & 预训练 $n$ & 预训练 $q_{0.85}$         & 后训练 $n$ & 后训练 $q_{0.85}$        \\\midrule
    严格: 许可+去重+身份+开源权重 & $8$     & $33.1968$ & $4$     & $29.4380$ \\
    许可去重: 许可白名单+去重    & $98$    & $18.4544$ & $773$   & $38.1002$ \\
    最宽: 六任务+类型            & $182$   & $19.4413$ & $1873$  & $33.4021$ \\\bottomrule
  \end{tabular}
\end{table}
表~\ref{tab:q4-baseline-caliber} 中, $n$ 为入选模型数. 严格口径的预训练和后训练分位分别为 $33.1968$ 分和 $29.4380$ 分, 即表~\ref{tab:q4-v39-slowdown} 中主模型点估计的未舍入值. 另外两种口径用于衡量样本筛选对经验分位的影响.

严格口径仅包含 $8$ 条预训练记录和 $4$ 条后训练记录, 两类分位相差 $3.7588$ 分. 许可去重口径纳入 $98$ 条和 $773$ 条记录, 对应分位为 $18.4544$ 分和 $38.1002$ 分. 最宽口径纳入 $182$ 条和 $1873$ 条记录, 对应分位为 $19.4413$ 分和 $33.4021$ 分. 样本口径扩大后, 两类模型的分位排序反转, 显示主模型结论对入选规则敏感. 本章通过滚动检验与整组织留出报告其历史验证范围. 全历史最宽口径的预训练和后训练分位分别为 $18.9166$ 分和 $32.3742$ 分, 样本量为 $330$ 和 $2501$. 全历史严格口径分别为 $25.9633$ 分和 $29.2899$ 分, 样本量为 $22$ 和 $11$. 最宽口径的 180 日与全历史结果相差最多 $1.03$ 分.


\subsubsection{联合分位的规模--时间增量恒等分解}
保持类型与分位水平不变，令原单位斜率为 $b_N=100\theta_N/\sigma_{A,N}$、$b_t=100\theta_t/\sigma_{A,t}$，预测期限 $H$ 以年计. 本文无时间趋势衰减的联合候选满足
\begin{equation}\label{eq:q4-additive-contribution}
\Delta q_\tau=
\underbrace{b_N\log_{10}(N_H/N_0)}_{\Delta_N}
+\underbrace{b_tH}_{\Delta_t}.
\end{equation}
给定报告阈值 $\varepsilon_{\mathrm{share}}>0$, 当 $|\Delta q_\tau|>\varepsilon_{\mathrm{share}}$ 时, 定义有符号份额 $\omega_N=\Delta_N/\Delta q_\tau$ 和 $\omega_t=\Delta_t/\Delta q_\tau$, 并有 $\omega_N+\omega_t=1$. 公式~\eqref{eq:q4-additive-contribution} 由附录表~\ref{tab:q4-lambda-sensitivity} 主候选一行的起点与终点直接相减得到, 将结构候选前沿的总变化分解为规模项与时间项. 当前原单位规模斜率对应同一行的 $b_N$（分/十倍）, 时间斜率在数值上接近零. 在指数质量成本与历史算力增长情景下, 规模项随预测期限由一年延长至两年而相应增加, 时间项接近零, 点估计对应的规模与时间份额约为 $100\%$ 与 $0\%$. 200 次有效联合抽样进一步给出不同参数抽样下规模项与时间项的份额分布, 一年情景下时间份额的上界约为 $85\%$. 因此, 点估计与联合抽样结果形成互补, 既给出主情景下的明确分解, 又展示不同参数组合下贡献结构的变化范围, 使前沿增量的来源能够得到更完整的量化描述.

\subsubsection{第四问输出与证据边界}
第四问在数据证据、统计建模与资源情景之间建立了层次清晰的分析体系, 以下三点概述其贡献与适用范围.

\textbf{(1) 评测口径统一与前沿模型的结构保证.} C3 与 C8 的口径分析把历史记录统一到同一套六任务能力标尺: 六任务均分与 C1 官方平均分的最大绝对误差为 $1.42\times10^{-14}$, 各任务按随机猜测基线归一化后与 C1 一致到浮点精度. 定理~\ref{thm:q4-ordered-quantile-qp} 把有序分位模型的全域形状约束、凸规划求解与全局目标证书纳入同一框架: 共享斜率下全域非交叉等价于截距有序, 上下界只需在声明域的四个顶点检验, 因而形状保证可以脱离经验支持域单独成立. 279 份有序联合分位拟合记录的最大目标间隙为 $6.70588\times10^{-8}$, 经推论~\ref{cor:q4-slope-gap-v37} 传为共享斜率误差界 $2.59\times10^{-4}$ 与最大预测对比误差界 $0.1867$ 分. 命题~\ref{thm:q4-cqr-conditional} 进一步给出有限样本区间覆盖的秩论证, 使点预测、分位结构与区间分析处于同一统计框架.

\textbf{(2) 第三问资源路径向能力情景的确定性接口.} 跨问只传递训练 FLOPs $6ND$:
以 $C_0=10^{22}$、8192 Token 的 $Q^*=1$ 分支为锚点, 按式~\eqref{eq:q4-training-match-v8}
反求满足训练算力目标的预算. 该分支的唯一性与饱和条件由引理~\ref{prop:q4-training-root} 保证, 表~\ref{tab:q4-v26-resources} 中 18 条资源路径全部经完整质量剖面比较确认 $Q^*=1$, 与该引理的饱和分支条件一致. 这一设计把第三问的资源优化结果与第四问的能力前沿连成 "资源配置—训练算力—能力情景"的连续链条, 同时保持前三问的 Loss 口径与第四问的 Benchmark 口径互不混用.

\textbf{(3) 多层次验证体系及其边界.} 开发期定权、滚动检验、整组织留出与联合抽样互不替代:
开发期 16 个标签含 8 次 1 月窗与 8 次 3 月窗, 复用晚期 8 次预测涉及 4 个仓库与 4 个组织, 整组织留出 33 条记录涉及 18 个组织. 联合规模--时间候选在复用晚期取得显著的样本外优势, Pinball 损失 $1.189958$ 相对经验分位与辅助趋势基准分别下降 $53.88\%$ 与 $56.27\%$, 同时保持分位有序; 但同一候选在开发期与整组织留出上弱于辅助趋势基准 ($4.726347$ 对 $3.819872$、$3.724345$ 对 $1.815348$), 三候选的开发期定权最优权重 $(0,0,1)$ 也指向辅助趋势基准, 故其优势限于复用晚期这一口径. 此外时间外检验只覆盖 1 与 3 个月窗口, 12/24 月结果依赖不含时间或规模斜率的主模型, 个体分位带的历史实测覆盖为 $50.0\%$--$75.0\%$, 低于名义 $90\%$, 故长期数值应作为条件情景而非已验证预测读取.

总体而言, 第四问形成了“数据口径统一—有序分位建模—历史验证—资源路径接入—未来情景分析”的完整分析链条. 各环节之间具有明确的数学接口和数据对应关系, 既增强了能力前沿建模的结构化程度, 也使不同算力增长情景下的模型能力变化能够得到清晰、可量化且具有较强解释性的展示.
% 原稿, 注释备查
%第四问在数据证据、统计建模与资源情景之间建立了较为完整且层次清晰的分析体系. C3/C8 的口径分析首先统一历史记录的六任务能力标尺, 为后续能力前沿建模提供一致的数据基础；定理~\ref{thm:q4-ordered-quantile-qp} 将有序分位模型的全域形状约束、凸规划求解与全局目标证书统一于同一框架, 使模型在数学结构与数值实现之间保持良好一致性；命题~\ref{thm:q4-cqr-conditional} 则进一步补充有限样本区间覆盖的理论依据, 使点预测、分位结构与区间分析形成较为完整的统计框架.第三问的资源路径通过统一的训练 FLOPs $6ND$ 接入式~\eqref{eq:q4-training-match-v8}. 且表~\ref{tab:q4-v26-resources} 中各路径均经完整质量优化确认 $Q^*=1$, 与引理~\ref{prop:q4-training-root} 的饱和分支条件保持一致. 这一设计将第三问的资源优化结果与第四问的能力前沿自然衔接, 在保持不同问题建模口径清晰的同时, 建立了从资源配置、训练算力到能力情景的连续量化链条.在预测层面, 开发期定权、滚动检验、整组织留出和联合抽样共同构成多层次验证体系. 联合规模--时间候选在复用晚期取得较低的 Pinball 损失, 同时保持有序分位结构, 使预测精度与前沿形状约束能够在同一模型中统一表达. 在此基础上, 12/24 月结果进一步将历史起点、训练算力增长路径、模型类型与条件能力前沿整合到统一情景框架中.总体而言, 第四问形成了“数据口径统一—有序分位建模—历史验证—资源路径接入—未来情景分析”的完整分析链条. 各环节之间具有明确的数学接口和数据对应关系, 既增强了能力前沿建模的结构化程度, 也使不同算力增长情景下的模型能力变化能够得到清晰、可量化且具有较强解释性的展示.
\section{模型评价、局限与改进方向}\label{sec:conclusion}

本节不再复述各问结论（逐条答案与出处见第 1.5 节表~\ref{tab:v42-answer-index}，各问的数值结果与适用范围见相应小节），只评价模型本身的结构优缺点、明确解释边界，并给出后续改进的优先级. 

\subsection{模型的主要优势}
本文的主要特点在于解析结构、决策目标与数据证据之间保持一致，并通过可证明的数学性质和可复核的数值证据形成完整的方法链条.
\setcounter{obsitem}{0}
\obsitem{明确了参数的可识别范围与分析边界：}
在固定域质量下，配比系数与质量交互系数存在严格的线性相关关系（命题~\ref{prop:q1-quality-collinearity-v15}）；锚点数据无法单独识别跨来源配比迁移系数（命题~\ref{prop:q2-lambda-unidentified-v15}）；对于未映射域，完整质量由命题（命题~\ref{prop:background-quality-bound}） 给出可识别区间. 这些结果将参数识别条件直接写入模型结构，使后续各问的推断范围具有明确的数学依据.
\obsitem{能够定量刻画预测误差对决策结果的影响：}
候选集合上的一致误差界通过三次比较传递为选择后悔界（命题~\ref{prop:q1-decision-regret-v37}）；有序分位规划的全局目标间隙则通过强凸性进一步转化为共享斜率误差界（推论~\ref{cor:q4-slope-gap-v37}）. 因此，模型不仅给出预测值和最优选择，还能够对预测精度、优化精度与最终决策质量之间的关系进行定量刻画，使“预测是否准确”和“决策是否可靠”处于统一的误差分析框架中.
\obsitem{预算变化下的最优配置具有严谨的数学证明：}
仿射下包络的凹性与严格单调性保证全局最优参数量随预算严格增加，同时最优质量不下降（定理~\ref{thm:q3-global-monotonicity-v37}）. 进一步地，有限预算下的质量饱和与质量成本割线比有界之间具有充要关系（定理~\ref{thm:q3-finite-saturation-iff-v37}）. 由此，参数规模、数据规模与质量投入随预算变化的规律不仅由数值结果展示，还具有相应的结构性判据.
\obsitem{主要数值结果均具有可复核的求解精度保证：}
45 个名义配置的最大目标间隙、预算相对残差、稳健配置的基准误差上界以及 279 份有序分位拟合的目标间隙均被统一记录并归档. 这些数值证书与前述解析结果相互配合，使模型从理论推导、数值求解到最终决策均具有明确的精度依据，并提高了整个建模流程的可复核性与一致性.


\subsection{模型的主要局限}
\setcounter{obsitem}{0}
\obsitem{质量衡量仍是间接的:}
本文的质量变量 $Q$ 由文本信号构造, 并非直接的人类评价或真实下游训练效果. 六个可映射域覆盖锚点配比的 $31.94\%$, 因此质量分析主要反映当前指标体系下的相对差异.
\obsitem{部分结果依赖模型设定:}
第二、三问中的质量响应主要来自经真实数据校准的半合成数据, 成本函数也按照题目给定的训练、质量处理和注意力开销构造. 因此, 得到的最优配置与资源路径对应本文所采用的质量响应和成本设定.
\obsitem{未来能力验证数据还比较少:}
第四问的严格前沿样本为 33 条, 开发期包含 16 个预测标签. 现有数据用于构建条件能力前沿并开展滚动检验和整组织留出分析, 后续随着更多模型和长期数据积累, 可进一步验证 12/24 月能力情景.

\subsection{模型推广}
本文的三层结构可迁移到其他"有限预算下的多要素配置"问题: 先用部分识别界定数据能支持的质量与配比信息, 再把不可直接观测的要素折算为有效资源量, 最后用成本下包络把多维配置问题约化为可全局比较的一维尺度问题. 

其中第二、三问依赖"质量响应单调且存在有效数据折算"这一结构, 因此在语料质量存在非单调效应或多模态训练等场景下, 需要先重新检验该前提.
\subsection{AI 使用声明}
本文章使用了chatGPT工具，版本为GPT-6 Sol / GPT-6 Luna开发公司为 OpenAI版本颁布日期2026年9月22日. \\
主要用于辅助概念理解（\url{https://chatgpt.com/share/6ab68617-b2b0-83e8-b719-ff1822b2f894}）、\\
资料搜索与整理（\url{https://chatgpt.com/share/6ab6857e-54fc-83ee-ac02-d0bb01562831}）、程序实现与调试、数值结果核对、文字整理及版面核查（\url{https://chatgpt.com/share/6ab68617-b2b0-83e8-b719-ff1822b2f894}）\\（\url{https://chatgpt.com/s/cx_6ab68696564c8191a4953b0863f7668e}）\\
全部模型设定、数学推导、结果解释与结论由参赛作者完成并逐项复核.\\
AI 未用于替代核心建模、推导与论证. 参赛作者对提交稿的模型设定、数据来源、计算结果和结论承担最终复核责任. 本文已列明的程序输出与证据边界供逐项复核.
\bibliographystyle{gmcm}
\bibliography{reference}

% 目录展示正文层级与五个附录标题；附录内部条目仍保留 PDF 书签。
\addtocontents{toc}{\protect\setcounter{tocdepth}{1}}
\clearpage
\begin{appendices}
  % 模板的 \appendix 只改写 \thesection，表/图仍沿用"节号.序号"，在附录中显示为
  % "表 1.46"这类无意义编号。此处把附录内的浮动体编号交给附录计数器，改为"表 A.1"。
  \makeatletter
  % 每个附录起始处重置浮动体计数器，使编号形如"表 A.1"、"表 D.1"而非跨附录累加。
  \newcommand{\resetappendixfloats}{%
    \setcounter{table}{0}%
    \setcounter{figure}{0}%
  }
  \setcounter{section}{0}
  \setcounter{table}{0}
  \setcounter{figure}{0}
  \renewcommand{\thetable}{\@Alph\c@section.\@arabic\c@table}
  \renewcommand{\thefigure}{\@Alph\c@section.\@arabic\c@figure}
  \makeatother
  \renewcommand{\thesubsection}{\thesection.\arabic{subsection}}
  \renewcommand{\thesubsubsection}{\thesubsection.\arabic{subsubsection}}
  \section{预备记号与通用数学工具与}\label{sec:math-background}

本部分列出正文四问所用数学工具及相关定义如可识别性、概率单纯形与 Fisher--Rao 几何、对数弹性等, 正文保留本题主定理的假设和关键证明. 线性映射与内积记号参考 Axler\cite{axler}，信息几何参考 Amari 与 Pistone\cite{amari2016,pistone}，约束优化参考 Boyd--Vandenberghe\cite{boyd}，分位数与 Pinball 损失参考 Koenker\cite{koenker}. 

\paragraph{基本记号. \quad}
对非空集合 $\mathcal A$ 与实函数 $f$，$\inf_{x\in\mathcal A}f(x)$ 表示函数值集合的下确界，即所有下界中的最大者；定义
\[
\argmin_{x\in\mathcal A}f(x)
:=\left\{x\in\mathcal A:\ f(x)=\inf_{y\in\mathcal A}f(y)\right\}.
\]
当下确界未被任何点取得时，$\argmin$ 可以为空. 若 $U\subseteq\mathbb R^m$ 为开集，记 $f\in C^1(U)$ 表示 $f$ 在 $U$ 上可微且一阶偏导数连续. 

\subsection{梯度与约束方向}\label{sec:bg-differential}

多元函数在约束集上的局部变化可由方向导数描述. 对可微映射，方向导数由 Jacobian 作用于方向向量得到；对标量函数，则退化为梯度与方向向量的内积. 

\begin{definition}[一阶可行方向]\label{def:bg-feasible-direction}
设约束写成 $a(\mathbf x)=0$，且 $a$ 可微. 若从可行点 $\mathbf x$ 沿方向 $\mathbf u$ 做小扰动仍希望在一阶近似下保持约束，则必须有
\[
Da(\mathbf x)[\mathbf u]=\nabla a(\mathbf x)^\top\mathbf u=0.
\]
满足该等式的 $\mathbf u$ 称为该等式约束的一阶可行方向. 
\end{definition}

对单个等式约束 $a(\mathbf x)=0$，所有满足 $\nabla a(\mathbf x)^\top\mathbf u=0$ 的一阶可行方向构成一个线性子空间；当 $\nabla a(\mathbf x)\neq0$ 时，该子空间的维数为 $m-1$. 


\begin{proposition}[正交投影与约束梯度]\label{prop:bg-orthogonal-projection}
设 $\mathbf q\in\mathbb R^m$ 为单位向量，超平面
\[
H_{\mathbf q}:=\{\mathbf z\in\mathbb R^m:\mathbf q^\top\mathbf z=0\}.
\]
则
\[
\mathbf P_{\mathbf q}:=\mathbf I-\mathbf q\mathbf q^\top
\]
是到 $H_{\mathbf q}$ 的正交投影矩阵. 若 $f:\mathbb R^m\to\mathbb R$ 可微，则 $f$ 在约束方向空间 $H_{\mathbf q}$ 上的一阶梯度为
\begin{equation}\label{eq:bg-projected-gradient}
\nabla_H f=\mathbf P_{\mathbf q}\nabla f.
\end{equation}
若 $\nabla_H f\neq0$，所有满足 $\mathbf z\in H_{\mathbf q}$、$\|\mathbf z\|_2=1$ 的方向中，最陡下降方向为 $-\nabla_H f/\|\nabla_H f\|_2$. 
\end{proposition}

\begin{definition}[参数的结构可识别性]\label{def:bg-identifiability}
设模型输出写成 $m(x;\boldsymbol\theta)$，观测设计点集合为 $\mathcal X$. 若
\[
m(x;\boldsymbol\theta)=m(x;\boldsymbol\theta')\quad\forall x\in\mathcal X
\]
必然推出 $\boldsymbol\theta=\boldsymbol\theta'$，则称参数 $\boldsymbol\theta$ 在设计 $\mathcal X$ 上可识别；否则称存在不可识别参数方向. 
\end{definition}

\begin{proposition}[线性参数模型的秩判据]\label{prop:bg-linear-identifiability}
在线性模型 $\mathbf y=\mathbf X\boldsymbol\theta$ 中，参数在给定设计矩阵 $\mathbf X$ 上唯一可识别，当且仅当
\[
\ker\mathbf X=\{\mathbf0\},
\]
即 $ \exists \mathbf h\in\ker\mathbf X$, $\forall t \in \mathbb{R}$，$\boldsymbol\theta$ 与 $\boldsymbol\theta+t\mathbf h$ 给出相同拟合值. 
\end{proposition}

\begin{lemma}[连续单调方程的唯一根]\label{lem:bg-monotone-root}
设 $R:(0,\infty)\to(0,\infty)$ 连续且严格递增，并满足
\[
\lim_{x\to0^+}R(x)<a<\lim_{x\to+\infty}R(x)
\]
（允许两端极限分别为 $0$ 与 $+\infty$）. 则方程 $R(x)=a$ 在 $(0,\infty)$ 上恰有一个解. 
\end{lemma}

\subsection{概率单纯形与 Fisher--Rao 信息几何}\label{sec:bg-info-geometry}

概率向量满足“分量非负、总和为 $1$”，因此其可行集合是概率单纯形，而不是整个 $\mathbb R^K$. 若一个分量增加，为保持总和不变，其他分量的总变化量必须相应减少. Fisher--Rao 度量为这种受约束的概率变化规定局部长度，并与相邻概率分布的统计可区分性相联系. 

\begin{definition}[概率单纯形及其相对内部]\label{def:background-simplex}
设 $K\ge2$. 定义
\[
\Delta_{K-1}:=\{\pvec\in\mathbb R_{\ge0}^{K}:\one^\top\pvec=1\},
\quad
\Delta_{K-1}^{\circ}:=\{\pvec\in\Delta_{K-1}:p_i>0\ \forall i\}.
\]
这里 $\Delta_{K-1}^{\circ}$ 指单纯形在仿射包络 $\{\pvec:\one^\top\pvec=1\}$ 内的相对内部. 对相对内部点 $\pvec$，约束函数为 $a(\pvec)=\one^\top\pvec-1$. 由定义~\ref{def:bg-feasible-direction}，一阶可行方向满足 $\one^\top\mathbf u=0$，故定义
\[
T_{\pvec}\Delta_{K-1}^{\circ}
:=\{\mathbf u\in\mathbb R^K:\one^\top\mathbf u=0\}.
\]
这里 $T_{\pvec}\Delta_{K-1}^{\circ}$ 是一个 $K-1$ 维线性子空间，可理解为“在不改变总概率的前提下允许的变化方向”. 
\end{definition}

\begin{definition}[KL 散度]\label{def:bg-kl}
对 $\pvec,\mathbf r\in\Delta_{K-1}^{\circ}$，定义
\[
D_{\mathrm{KL}}(\pvec\|\mathbf r)
:=\sum_{i=1}^{K}p_i\ln\frac{p_i}{r_i}.
\]
KL 散度一般不对称，因此不是通常意义下的距离. 当 $\mathbf r$ 接近 $\pvec$ 时，其二阶展开在切空间 $T_{\pvec}\Delta_{K-1}^{\circ}$ 上给出局部二次型. 
\end{definition}

以下使用 Landau 大 $O$ 记号：当 $\varepsilon\to0$ 时，$R(\varepsilon)=O(\varepsilon^m)$ 表示存在常数 $C>0$ 与 $\varepsilon_0>0$，使得 $0<|\varepsilon|<\varepsilon_0$ 时有 $|R(\varepsilon)|\le C|\varepsilon|^m$. 

\begin{lemma}[KL 散度的局部二阶展开]\label{lem:bg-kl-fisher}
设 $\pvec\in\Delta_{K-1}^{\circ}$、$\mathbf u\in T_{\pvec}\Delta_{K-1}^{\circ}$，并取足够小的 $\varepsilon$ 使 $p_i+\varepsilon u_i>0$. 则
\begin{equation}\label{eq:bg-kl-taylor}
D_{\mathrm{KL}}(\pvec\|\pvec+\varepsilon\mathbf u)
=\frac{\varepsilon^2}{2}\sum_{i=1}^{K}\frac{u_i^2}{p_i}+O(\varepsilon^3).
\end{equation}
\end{lemma}

引理~\ref{lem:bg-kl-fisher} 说明：两个很接近的概率分布之间，KL 散度的一阶差异被总和约束消掉，最先出现的非零项是一个二次型. 这个二次型自然给出切空间 $T_{\pvec}\Delta_{K-1}^{\circ}$ 上的局部内积. 

\begin{definition}[Fisher--Rao 度量]\label{def:bg-fisher}
对 $\mathbf u,\mathbf v\in T_{\pvec}\Delta_{K-1}^{\circ}$，定义
\begin{equation}\label{eq:bg-fisher}
g_{\pvec}(\mathbf u,\mathbf v)
:=\sum_{i=1}^{K}\frac{u_iv_i}{p_i}.
\end{equation}
因为 $p_i>0$，有 $g_{\pvec}(\mathbf u,\mathbf u)>0$（除非 $\mathbf u=0$），所以 $g_{\pvec}$ 在每个切空间上都是内积. 引理~\ref{lem:bg-kl-fisher} 可写成
\[
D_{\mathrm{KL}}(\pvec\|\pvec+\varepsilon\mathbf u)
=\frac{\varepsilon^2}{2}g_{\pvec}(\mathbf u,\mathbf u)+O(\varepsilon^3),
\]
这就是 Fisher--Rao 度量与局部统计可区分性的联系\cite{amari2016,pistone}. 
\end{definition}

\begin{theorem}[平方根映射把 Fisher--Rao 几何化为球面几何]\label{thm:bg-sqrt-isometry}
定义平方根映射
\[
\Phi:\Delta_{K-1}^{\circ}\to\mathbb R^K,
\quad
\Phi(\pvec):=(\sqrt{p_1},\ldots,\sqrt{p_K})^\top.
\]
则 $\|\Phi(\pvec)\|_2=1$，即 $\Phi(\pvec)$ 位于单位球面的正象限；并且对任意 $\mathbf u,\mathbf v\in T_{\pvec}\Delta_{K-1}^{\circ}$，
\begin{equation}\label{eq:bg-sqrt-isometry}
g_{\pvec}(\mathbf u,\mathbf v)
=4\left\langle D\Phi(\pvec)[\mathbf u],D\Phi(\pvec)[\mathbf v]\right\rangle .
\end{equation}
任意光滑曲线的 Fisher--Rao 弧长等于其平方根像在单位球面上的球面弧长的$2$倍.
\end{theorem}

\begin{corollary}[Fisher--Rao 测地距离]\label{cor:bg-fr-distance}
对 $\pvec,\mathbf r\in\Delta_{K-1}^{\circ}$，令
\[
\theta:=\arccos\langle\Phi(\pvec),\Phi(\mathbf r)\rangle
=\arccos\!\left(\sum_{i=1}^{K}\sqrt{p_ir_i}\right).
\]
则
\begin{equation}\label{eq:bg-fr-distance}
d_{\mathrm{FR}}(\pvec,\mathbf r)=2\theta
=2\arccos\!\left(\sum_{i=1}^{K}\sqrt{p_ir_i}\right).
\end{equation}
\end{corollary}

\begin{remark}
Fisher--Rao 度量中的 $1/p_i$ 要求 $p_i>0$，因此上述微分公式定义在单纯形内部. 当 $p_i\to0^+$ 时，对应方向的局部权重 $1/p_i$ 发散；边界点需通过内部极限或其他适用于边界的参数化处理. 
\end{remark}

\paragraph{平方根坐标的几何解释. \quad}
平方根变换 $\phi_i=\sqrt{p_i}$ 将概率单纯形内部映射到单位球面的正象限，并把 Fisher--Rao 局部度量化为常数倍的欧氏内积. 该坐标变换规定了自然的几何尺度，但并不决定其上应采用哪一种具体统计回归模型. 

\subsection{对数微分、弹性与幂律}\label{sec:bg-elasticity}

当变量具有不同量纲时，绝对导数 $\partial y/\partial x$ 会随计量单位改变. 对数弹性刻画相对变化之间的局部比例关系，因此是无量纲的边际量. 

\begin{definition}[局部对数弹性]\label{def:background-elasticity}
设 $y=y(x_1,\ldots,x_m)>0$ 可微，且 $x_i>0$. 定义 $y$ 对 $x_i$ 的局部对数弹性为
\begin{equation}\label{eq:bg-elasticity}
\varepsilon_{y,x_i}
:=\frac{\partial\ln y}{\partial\ln x_i}
=\frac{x_i}{y}\frac{\partial y}{\partial x_i},
\end{equation}
其中其余变量保持不变. 若相对变化很小，则一阶 Taylor 近似给出
\[
\frac{\Delta y}{y}\approx
\sum_{i=1}^{m}\varepsilon_{y,x_i}\frac{\Delta x_i}{x_i}.
\]
\end{definition}

\begin{proposition}[乘积幂律的对数微分]\label{prop:bg-log-monomial}
若
\[
y=a\prod_{i=1}^{m}x_i^{b_i},\quad a>0,\ x_i>0,
\]
则
\begin{equation}\label{eq:bg-log-monomial}
\dd\ln y=\sum_{i=1}^{m}b_i\,\dd\ln x_i,
\quad
\varepsilon_{y,x_i}=b_i.
\end{equation}
\end{proposition}


命题~\ref{prop:bg-log-monomial} 表明，幂律在对数坐标中转化为线性关系，幂指数直接成为局部弹性. 若目标函数由若干幂项相加构成，则总目标的弹性取决于各项在当前状态下的相对大小. 


\begin{definition}[渐近等价]\label{def:bg-asymptotic}
设 $f(c),g(c)>0$. 若
\[
\lim_{c\to\infty}\frac{f(c)}{g(c)}=1,
\]
则称 $f$ 与 $g$ 在 $c\to\infty$ 时渐近等价，记作 $f(c)\sim g(c)$. 
\end{definition}

\paragraph{渐近等价的运算规则. \quad}
若 $f\sim g$、$u\sim v$ 且各函数最终为正，则对任意固定实数 $r$，
\begin{equation}\label{eq:bg-asymptotic-rules}
f^r\sim g^r,\quad fu\sim gv,\quad \frac{f}{u}\sim\frac{g}{v}.
\end{equation}
这些关系直接来自极限的乘积、商与连续幂运算，例如
$(f/g)^r\to1$、$(f/g)(u/v)\to1$. 


\subsection{约束优化、凸分析与预算灵敏度}\label{sec:bg-optimization}

KKT 条件是不等式约束下拉格朗日乘子法的一阶推广. 在满足适当正则条件的局部最优点，目标函数的负梯度可表示为活动约束法向量的非负线性组合. 


\begin{lemma}[简单根的隐函数形式]\label{thm:bg-implicit-simple-root}
设 $F(x,q)$ 在 $(x_0,q_0)$ 附近连续可微，满足
\[
F(x_0,q_0)=0,\quad \frac{\partial F}{\partial x}(x_0,q_0)\neq0.
\]
则在 $q_0$ 的某个邻域内，方程 $F(x,q)=0$ 唯一确定一个连续可微函数 $x=x^\star(q)$，且\cite{bonnans}
\begin{equation}\label{eq:bg-implicit-derivative}
\frac{\dd x^\star}{\dd q}
=-\frac{F_q(x^\star(q),q)}{F_x(x^\star(q),q)}.
\end{equation}
\end{lemma}

\begin{proposition}[剖面函数的一阶导数]\label{prop:bg-profile-derivative}
设 $f(x,q)$ 连续可微，$x^\star(q)$ 是可微的内点驻点分支，即
\[
\frac{\partial f}{\partial x}(x^\star(q),q)=0.
\]
定义剖面函数 $\bar f(q):=f(x^\star(q),q)$，则
\begin{equation}\label{eq:bg-profile-derivative}
\frac{\dd\bar f}{\dd q}
=\frac{\partial f}{\partial q}(x^\star(q),q).
\end{equation}
\end{proposition}


\begin{definition}[活动约束、Lagrangian 与 LICQ]\label{def:bg-kkt-objects}
考虑
\begin{equation}\label{eq:bg-general-nlp}
\min_{\mathbf x\in\mathbb R^m} f(\mathbf x),
\quad g_j(\mathbf x)\le0,\quad j=1,\ldots,q.
\end{equation}
在可行点 $\mathbf x$，定义活动集
\[
\mathcal A(\mathbf x):=\{j:g_j(\mathbf x)=0\}.
\]
给定乘子 $\lambda_j\ge0$，Lagrangian 定义为
\[
\mathscr L(\mathbf x,\boldsymbol\lambda)
:=f(\mathbf x)+\sum_{j=1}^{q}\lambda_jg_j(\mathbf x).
\]
若活动约束的梯度集合
$\{\nabla g_j(\mathbf x):j\in\mathcal A(\mathbf x)\}$
线性无关，则称在 $\mathbf x$ 满足线性无关约束资格（LICQ）. 
\end{definition}

\begin{theorem}[KKT 一阶必要条件]\label{thm:bg-kkt}
设 $f,g_1,\ldots,g_q$ 在局部最优点 $\mathbf x^\star$ 附近连续可微，并且 $\mathbf x^\star$ 满足 LICQ. 则存在乘子 $\lambda_j^\star\ge0$，使
\begin{align}
\nabla f(\mathbf x^\star)+\sum_{j=1}^{q}\lambda_j^\star\nabla g_j(\mathbf x^\star)&=0,\label{eq:bg-kkt-stationarity}\\
g_j(\mathbf x^\star)&\le0,\label{eq:bg-kkt-primal}\\
\lambda_j^\star g_j(\mathbf x^\star)&=0.\label{eq:bg-kkt-complementary}
\end{align}
式~\eqref{eq:bg-kkt-stationarity}、\eqref{eq:bg-kkt-primal} 与~\eqref{eq:bg-kkt-complementary} 分别称为驻点条件、原始可行性与互补松弛\cite{boyd,bonnans}. 
\end{theorem}

\begin{remark}
若问题为凸优化，并满足适当的严格可行条件，则 KKT 条件也可成为全局最优的充分条件\cite{boyd}. 对于一般非凸问题，KKT 条件仅提供局部最优点的必要条件，不能保证全局最优. 
\end{remark}

\begin{definition}[参数化最优值函数]\label{def:background-value}
若预算参数 $c$ 进入约束 $G(\mathbf x)\le c$，定义
\[
V(c):=\inf\{f(\mathbf x):G(\mathbf x)\le c,\ \mathbf x\text{ 满足其余约束}\}.
\]
\end{definition}

\begin{proposition}[活动预算约束的局部灵敏度]\label{thm:bg-envelope}
设在某个预算 $c$ 的邻域内存在唯一且可微的最优分支 $\mathbf x^\star(c)$，活动约束集合保持不变，且 KKT 条件成立. 若预算约束 $G(\mathbf x)\le c$ 在该分支上活动，其 KKT 乘子为 $\mu^\star(c)$，则
\begin{equation}\label{eq:bg-shadow-price}
V'(c)=-\mu^\star(c).
\end{equation}
因此 $\mu^\star$ 表示增加一单位预算所带来的最优目标值的一阶下降量，即预算的影子价格. 
\end{proposition}


\paragraph{局部最优与灵敏度. \quad}
KKT 条件描述活动约束下的一阶平衡关系；最优值灵敏度公式描述约束右端参数变化时最优目标值的一阶响应. 两者均为局部性质，在非凸问题中不能替代不同可行分支之间的全局目标值比较. 


\paragraph{凸性、次梯度与多面体约束. \quad}
下面的凸分析工具用于第三问的唯一尺度与稳健内层，也用于第四问把连续形状约束化为有限顶点约束. 

\begin{definition}[严格凸性、强制性与次梯度]\label{def:bg-convex-v26}
在凸集 $D$ 上，若任意不同 $x,y\in D$ 及 $0<t<1$ 均有
$f(tx+(1-t)y)<tf(x)+(1-t)f(y)$，则称 $f$ 严格凸；把小于改为不大于即为凸性. 对 $f:\mathbb R\to\mathbb R$，若 $|x|\to\infty$ 时 $f(x)\to\infty$，称其强制. 凸函数在 $x_0$ 的次梯度 $v$ 满足
\[
f(x)\ge f(x_0)+v(x-x_0)\quad\text{对所有 }x;
\]
所有这类 $v$ 组成次微分 $\partial f(x_0)$. 可微时次梯度唯一，等于 $f'(x_0)$. 
\end{definition}

\begin{proposition}[唯一极小点、有限最大值与支撑下界]\label{prop:bg-convex-minimum-v26}
连续、严格凸且强制的实函数有唯一全局极小点，其充要最优条件为 $0\in\partial f(x^*)$. 有限个严格凸且强制的实函数 $f_1,\ldots,f_m$ 的最大值 $F=\max_jf_j$ 仍严格凸且强制；若各 $f_j$ 可微，则
\[
\partial F(x)=\operatorname{conv}\{f_j'(x):f_j(x)=F(x)\}.
\]
对任意凸函数，已知 $x_-, x_+$ 两处的次梯度 $v_-\le0\le v_+$，则极小点可在两点之间寻找，且两条支撑线最大值在该区间的下界为函数极小值的下界. 
\end{proposition}


\begin{definition}[凸包与顶点表示]\label{def:bg-polytope-v26}
有限点集 $v_1,\ldots,v_m$ 的凸包为
$\operatorname{conv}\{v_j\}:=\{\sum_j\omega_jv_j:\omega_j\ge0,\sum_j\omega_j=1\}$. 
有限点集的凸包称为多胞形；二维矩形恰为其四个顶点的凸包. 
\end{definition}
\begin{proposition}[仿射不等式的顶点检验]\label{prop:bg-affine-vertices-v26}
若 $a(x)=b+\boldsymbol w^\top x$，且 $a(v_j)\ge0$ 对所有顶点成立，则在其凸包内处处 $a(x)\ge0$. 
\end{proposition}


\paragraph{仿射下包络与全局比较. \quad}
局部驻点描述一条分支的变化；比较不同分支随预算的选择，还需要利用目标对参数的整体结构. 下面的下包络性质把这一步转化为可比较的直线族. 
\begin{definition}[仿射族的下包络]\label{def:bg-affine-envelope-v37}
设 $\mathcal I$ 为紧集，$a,b:\mathcal I\to\mathbb R$ 连续. 仿射族 $a(q)x+b(q)$ 的下包络定义为
\[
\mathscr H(x):=\min_{q\in\mathcal I}\{a(q)x+b(q)\}.
\]
在给定 $x$ 达到最小值的 $q$ 称为最小化指标. 
\end{definition}
\begin{proposition}[下包络的凹性与单调性]\label{prop:bg-affine-envelope-v37}
上述 $\mathscr H$ 连续且凹. 若 $a_{\min}:=\min_{q\in\mathcal I}a(q)>0$，则对 $x_2>x_1$，
\[
\mathscr H(x_2)-\mathscr H(x_1)\ge a_{\min}(x_2-x_1)>0.
\]
\end{proposition}

第三问以质量为指标构造上述直线族，再用两个预算下的最优性不等式证明全局单调性. 这是单调比较静态的题目专属形式\cite{milgrom_shannon}，无需假定完整资源问题联合凸. 

\subsection{条件分位数、Pinball 损失与有限样本秩}\label{sec:bg-quantile}

当研究对象是分布的上沿而非均值时，固定分位水平的条件分位数比样本最大值具有更稳定的统计含义. 样本最大值会随样本量变化，而条件分位数由条件分布本身决定，并可通过条件期望 Pinball 损失的最小化来刻画. 

\begin{definition}[条件分布函数与条件分位数]\label{def:background-quantile}
设随机变量 $S$ 表示能力分数，$X$ 表示条件变量. 定义条件分布函数
\[
F_{S\mid X}(s\mid x):=\Prob(S\le s\mid X=x).
\]
对 $\tau\in(0,1)$，定义条件 $\tau$ 分位数为广义逆
\begin{equation}\label{eq:background-quantile}
q_\tau(x):=F^{-1}_{S\mid X}(\tau\mid x)
:=\inf\{s:F_{S\mid X}(s\mid x)\ge\tau\}.
\end{equation}
并记左极限
$F_{S\mid X}(a^-\mid x):=\lim_{t\to a^-}F_{S\mid X}(t\mid x)$. 
\end{definition}


\paragraph{分位数的单调性. \quad}
对固定条件 $x$，若 $0<\tau_1<\tau_2<1$，则
\[
q_{\tau_1}(x)\le q_{\tau_2}(x).
\]
因为满足 $F_{S\mid X}(s\mid x)\ge\tau_2$ 的点必然也满足
$F_{S\mid X}(s\mid x)\ge\tau_1$，较高分位水平对应集合包含于较低分位水平对应集合；取下确界即得上述顺序. 经验分布的经验分位数满足同样的单调性. 


\paragraph{样本最大值与样本量. \quad}
设 $S_1,\ldots,S_m$ 独立同分布，公共分布函数为 $F$，并令
$M_m:=\max_{1\le i\le m}S_i$. 则
\begin{equation}\label{eq:bg-sample-maximum}
\Prob(M_m\le s)
=\Prob(S_1\le s,\ldots,S_m\le s)
=F(s)^m.
\end{equation}
若 $0<F(s)<1$，该概率随 $m$ 增大而减小，因此样本最大值的分布显式依赖样本量. 


\begin{definition}[Pinball 损失]\label{def:bg-pinball}
对预测误差 $e=S-a$，定义
\begin{equation}\label{eq:q4-pinball-v13}
\rho_\tau(e)
:=e\{\tau-\mathbf 1_{\{e<0\}}\}
=\begin{cases}
\tau e,&e\ge0,\\
(\tau-1)e,&e<0.
\end{cases}
\end{equation}
当预测 $a$ 低于观测 $S$ 时，误差按权重 $\tau$ 计；当预测过高时，误差按权重 $1-\tau$ 计. 因此 $\tau>1/2$ 会更重地惩罚低估，适合估计上分位位置\cite{koenker,koenker_bassett,koenker_hallock}. 本文只使用 $0.05$、$0.85$、$0.95$ 三个内部分位水平，不涉及尾部外推；极值分位回归的渐近理论与本文的有序凸规划是不同的问题，可参见文献~\cite{chernozhukov_extreme}. 
\end{definition}

\begin{lemma}[Pinball 损失及其条件期望的凸性]\label{lem:bg-pinball-convex}
固定观测 $s$，函数 $a\mapsto\rho_\tau(s-a)$ 为凸函数；因此条件期望损失（风险函数）
\[
R_x(a):=\E[\rho_\tau(S-a)\mid X=x]
\]
也是凸函数. 
\end{lemma}

\begin{proposition}[条件分位数的 Pinball 损失最小化表征]\label{thm:bg-pinball-quantile}
假设 $\E[|S|\mid X=x]<\infty$，使 $R_x(a)$ 有限. 则 $R_x(a)$ 的极小点集合恰由满足
\begin{equation}\label{eq:q4-quantile-subgradient-v13}
F_{S\mid X}(a^-\mid x)\le\tau\le F_{S\mid X}(a\mid x)
\end{equation}
的 $a$ 构成. 若条件分布在该位置连续且严格递增，则极小点唯一，且
\[
a^\star=q_\tau(x)=F^{-1}_{S\mid X}(\tau\mid x).
\]
\end{proposition}

\begin{corollary}[经验分位数最小化样本 Pinball 损失]\label{cor:bg-empirical-quantile}
给定样本 $s_1,\ldots,s_n$，任一经验 $\tau$ 分位数都是
\[
\argmin_{a\in\mathbb R}\sum_{i=1}^{n}\rho_\tau(s_i-a)
\]
中的一个解. 
\end{corollary}

\begin{remark}
平方损失的条件期望极小点是条件均值，而 Pinball 损失的条件期望极小点是条件分位数. 因此，均值回归回答“平均模型怎样变化”，高分位回归回答“分布上沿在哪里”，两者是不同统计目标，不能相互替代. 
\end{remark}


\paragraph{有限样本秩与可交换性. \quad}
条件分位数描述总体的概率位置，而有限样本区间还需要把校准分数的排序转化为覆盖概率；这一步由可交换秩结论完成. 

\begin{definition}[可交换性与次序统计量]\label{def:bg-exchangeable-v26}
随机向量 $(R_1,\ldots,R_m)$ 若在任意有限指标置换后联合分布不变，则称可交换；它不要求各分量相互独立. 将样本从小到大排列后的第 $k$ 个值记为 $R_{(k)}$，称为第 $k$ 个次序统计量. 
\end{definition}
\begin{lemma}[可交换分数的保守秩界]\label{thm:bg-exchangeable-rank-v26}
若 $R_1,\ldots,R_m,R_{\rm new}$ 可交换，则对 $1\le k\le m$，
\[
\Prob\{R_{\rm new}\le R_{(k)}\}\ge\frac{k}{m+1},
\]
其中右侧次序统计量仅由前 $m$ 个校准分数计算. 无并列值时等号成立；有并列值时该不等式保守. 
\end{lemma}


\paragraph{小结. \quad}
方向导数、正交投影与秩判据处理约束方向和参数可识别性；概率单纯形与 Fisher--Rao 几何处理配比变量；对数弹性与渐近等价处理标度律和高预算尺度；隐函数、剖面导数、KKT、凸性、次梯度、仿射下包络与包络定理共同处理预算优化及全局比较；条件分位数、Pinball 损失与可交换秩则处理能力前沿及有限样本校准. 上述五组工具分别进入正文四问，正文只引用结论而不再重复建立通用数学结论. 

表~\ref{tab:math-dependencies-v26} 汇总上述通用工具与各问直接结论的依赖关系. 
\begin{table}[htbp]\centering\small
\caption{通用数学背景到四问直接结论的依赖关系}\label{tab:math-dependencies-v26}
\begin{tabular}{P{1.0cm}P{5.6cm}P{7.9cm}}\toprule
问题 & 重点数学基础 & 本题直接结论及算法作用\\\midrule
一 & 单纯形、Fisher--Rao 几何、部分识别与一致误差界 & 平方根响应面、质量识别集、配比选择后悔界；提供 $Q_0,\psi$ 及接口条件. \\
二 & 秩与中心化、对数弹性、坐标变换、形状约束 & 共同质量响应的秩一必要条件、有限等价与坐标弹性变换；来源内分组拟合. \\
三 & 严格凸性、仿射下包络、最优性比较、渐近等价 & 唯一尺度、全局预算单调、有限饱和充要条件、完备候选和条件渐近；风险内层凸性. \\
四 & Pinball 损失的凸性、顶点约束、部分强凸性、顺序反馈与可交换秩 & 有序分位规划、斜率误差界、凸组合和累计后悔；带前提的校准与训练算力匹配. \\\bottomrule

\end{tabular}\end{table}
\resetappendixfloats

  \section{技术性证明与计算支撑}\label{app:technical-proofs}
本附录集中给出不改变主理论依赖关系、但用于数值实现或跨问接口的技术性证明. 正文保留第二问有限等价及第三问全局预算单调、有限饱和充要条件的完整证明；本附录给出拟合、风险计算和跨问接口的技术支撑. 

\subsection{变量投影局部导数的推导}\label{app:q2-projection-proof}
\begin{proof}固定正对角权矩阵 $\mathbf W$，对加权残差平方求导，有额外项 $-2\boldsymbol r^\top\mathbf W\mathbf X\partial_{\eta_j}\boldsymbol b^*$. 自由坐标满足正规方程 $\mathbf X_{\rm free}^\top\mathbf W\boldsymbol r=0$，活动坐标在该邻域固定、导数为零，故额外项消失. 自由列满秩保证局部系数解光滑. 活动集改变处须按带界子问题处理，不能直接延用此局部光滑结论. \end{proof}

\subsection{稳健尺度内层唯一性的证明}\label{app:q3-robust-proof}
\begin{proof}
有限个严格凸函数的最大值仍严格凸，证明见背景命题~\ref{prop:bg-convex-minimum-v26}；有限常数平移不改变强制性. 故唯一极小点存在. 由 $V^-_\theta\le V_\theta\le V^+_\theta$，逐情景有 $\ell^-_\theta-V^+_\theta\le R_\theta\le\ell_\theta-V^-_\theta$，再取最大值即可. 这样不会把尚有误差的个体最优值误当作精确基准. 
\end{proof}

\subsection{固定质量参数稳定性命题的证明}\label{app:q3-stability-proof}
\begin{proof}
定理~\ref{thm:q3-unique-scale} 给出 $\ell_{xx}>0$. 沿紧参数线段，唯一根 $x^*(t)$ 由隐函数定理连续可微；连续正函数 $\ell_{xx}$ 在该紧轨迹上取得正最小值 $m_x$，而连续向量函数 $\nabla_{\boldsymbol\vartheta}\ell_x$ 取得有限最大范数 $L_{\vartheta}$. 对恒等式
$\ell_x(x^*(t),Q;\boldsymbol\vartheta(t))=0$ 关于 $t$ 求导，得到
\[
\frac{\dd x^*}{\dd t}
=-\frac{\nabla_{\boldsymbol\vartheta}\ell_x^{\top}
(\boldsymbol\vartheta_1-\boldsymbol\vartheta_0)}{\ell_{xx}}.
\]
由 Cauchy--Schwarz 不等式，
$|\dd x^*/\dd t|\le(L_{\vartheta}/m_x)
\|\boldsymbol\vartheta_1-\boldsymbol\vartheta_0\|_2$；在 $t\in[0,1]$ 上积分即得式~\eqref{eq:q3-stability-x}. 因为 $\ln n^*=x^*$，式~\eqref{eq:q3-stability-n} 随即成立. 固定 $Q$ 时
\[
\ln d^*=\ln c-\ln(ke^{x^*}+h(Q)),\quad
\frac{\partial\ln d^*}{\partial x^*}=-r\in[-1,0),
\]
故中值定理给出
$|\Delta\ln d^*|\le|\Delta x^*|$，得到式~\eqref{eq:q3-stability-d}. 最后，
$\partial_{\ln A}\ell_x=-\alpha\mathcal L_N=-M$，
$\partial_{\ln B_p}\ell_x=\beta\mathcal L_Dr=M$；再次使用隐函数公式即得式~\eqref{eq:q3-stability-AB}. 
\end{proof}

\subsection{训练 FLOPs 匹配根唯一性的证明}\label{app:q4-training-root-proof}
\begin{proof}
固定质量下尺度驻点的一阶必要条件为 $\alpha An^{-\alpha}=\beta B_pd^{-\beta}kn/(kn+h_1)$. 代入训练约束 $d=T_H/(6n)$ 并整理即得式~\eqref{eq:q4-training-root}. 左侧关于 $\ln n$ 的对数导数为
\[
\alpha+\beta+1-\frac{kn}{kn+h_1}>0,
\]
且从零连续增至无穷，故唯一正根存在. 该推导只使用 $Q=1$ 的驻点，因此还需第三问全质量区间比较来确定完整问题的最优性. 
\end{proof}

\subsection{成本倍数与上下文缩放关系的推导}\label{app:q3-scaling-proof}
\begin{proof}代入 $n=s_hm/k$、$kn+h=s_h(m+H)$ 并提取公共正因子即得式~\eqref{eq:q3-dimensionless-scaling}；固定其 $\Upsilon$ 解出 $c$ 即得比例式. \end{proof}

\subsection{顺序反馈下的理论扩展}\label{app:q4-online-extension}
以下结果比较同一顺序预测协议中的固定候选，不参与前述开发期定权或历史回溯选择. 
\Needspace{13\baselineskip}
\begin{proposition}[指数加权分位预测的累计后悔界]\label{thm:q4-online-regret-v37}
设 $y_t=S_t/100\in[0,1]$，各候选在观测 $y_t$ 前给出归一化预测 $q_t^{(k)}\in[0,1]$，并且上一期所有候选损失在下一次预测前已可用. 固定目标分位 $\tau$，令
\[
\ell_t^{(k)}=\rho_\tau(y_t-q_t^{(k)})\in[0,1],\quad
a_{t,k}=\frac{\exp\{-\eta\sum_{s<t}\ell_s^{(k)}\}}
{\sum_{j=1}^K\exp\{-\eta\sum_{s<t}\ell_s^{(j)}\}},\quad\eta>0.
\]
组合 $\widehat q_t=\sum_ka_{t,k}q_t^{(k)}$ 对任意结果序列满足
\begin{equation}\label{eq:q4-online-regret-v37}
\sum_{t=1}^T\rho_\tau(y_t-\widehat q_t)
-\min_k\sum_{t=1}^T\ell_t^{(k)}
\le\frac{\ln K}{\eta}+\frac{\eta T}{8}.
\end{equation}
\end{proposition}
\begin{proof}
证明见\ref{app:q4-online-regret-v37}，采用有界损失的标准指数权重势函数论证\cite{freund_schapire}. 
\end{proof}
该界与可交换性无关，比较的是同一预测序列上的最佳固定候选；取 $\eta=\sqrt{8\ln K/T}$ 得 $\sqrt{T\ln K/2}$ 的累计界. 多分位共用按 $0.85$ 损失更新的权重可以保序，但上述后悔保证只对应所用目标. 实际一月、三月预测含标签延迟和重叠窗口，需按真实反馈时间构造协议或另作延迟分析，不能直接套用即时反馈界；当前晚期样本也不足以检验长期收益. 本文已有历史比较与未来情景仍使用前述冻结候选. 

为检查更新次序，以两条在当期标签出现前固定的合成候选预测运行 $T=256$ 步、即时揭示标签的程序协议，观察到归一化累计后悔 $2.60$，低于式~\eqref{eq:q4-online-regret-v37} 在相同 $K=2$、$\eta=\sqrt{8\ln2/256}$ 下的上界 $9.42$. 这仅是定理实现与信息时序的合成核验，不是 C 系列真实时间窗的表现；标签延迟仍需单独建模. 

顺序反馈还可用于自适应共形校准的长期经验误覆盖控制\cite{gibbs_aci}，但不能据此给任意未来时点或任意组织提供条件覆盖保证. 现有 C4 事后统一快照、评测版本变化和长期标签缺失，仍限制十二月、二十四月情景的统计解释. 

\subsection{指数加权分位预测后悔界的证明}\label{app:q4-online-regret-v37}
\begin{proof}
令 $W_t=\sum_{k=1}^K\exp\{-\eta\sum_{s<t}\ell_s^{(k)}\}$，则 $W_1=K$. 有界变量的指数矩界给出
\[
\ln\frac{W_{t+1}}{W_t}
=\ln\sum_ka_{t,k}e^{-\eta\ell_t^{(k)}}
\le-\eta\sum_ka_{t,k}\ell_t^{(k)}+\frac{\eta^2}{8}.
\]
逐轮相加，并用任一候选 $k$ 的一项作为 $W_{T+1}$ 的下界，得到
\[
-\eta\sum_t\ell_t^{(k)}-\ln K
\le\ln\frac{W_{T+1}}{W_1}
\le-\eta\sum_{t,k}a_{t,k}\ell_t^{(k)}+\frac{T\eta^2}{8}.
\]
移项后对 $k$ 取最小，再由命题~\ref{prop:q4-convex-aggregate-v37} 用组合损失替代加权损失，即得式~\eqref{eq:q4-online-regret-v37}. 该论证按实际可用的历史反馈更新权重，故标签延迟需要额外处理. 
\end{proof}
\resetappendixfloats

  \section{关键数值算法与复核路径}\label{app:v39-algorithms}
第一问按相同 \texttt{index} 对齐每对配比与 Loss 文件，先冻结 A5 分母与锚点；重复交叉验证选型之后，分别在平方根二次面和平方根 RBF 上计算相同的 17 个单域与 136 个双域探针，并在各自固定模型下对训练配方作 80 次重抽样. 第二问以 B1 拟合经典标度律，在 B6 估计质量响应；交叉验证以完整 $(N,D)$ 单元或质量等级分组，每折重估全部待比较参数. 

第三问对每个给定 $(c,Q)$，令 $x=\ln n$，用严格凸问题的一阶最优性条件求 $n^*(Q,c)$，再取 $d=c/(kn+h(Q))$；外层使用端点、驻点及质量区间下界进行全局比较，记录目标上下界和预算残差. 对 $C7$ 五个上下文、三类成本和三档预算逐项保存名义结果，平均--CVaR 另解相应上图问题. 第四问在冻结组织抽样与联合分位拟合后，将 $\zeta r_CH$ 代入训练 FLOPs 目标；零增长复用锚点，减半增长一年新解完整第三问，减半增长两年与延续增长一年共用相同资源目标，再重算每次分位预测与两类区间. 该复用只因训练目标相同，所有新路径仍检查预算、质量最优性和原延续情景回算. 


\subsection{第一问参考方法的补充数值对照}\label{app:q1-method-crosschecks}
表~\ref{tab:q1-utility} 在正文给出多目标结果；这里保留模型族、训练强度、桥接形式与时间项的来源内诊断，均不替换四问的主模型. 
\paragraph{多目标证据. \quad} 表~\ref{tab:q1-utility} 说明“平均性能”和“最坏域性能”不是同一个优化目标. 两者 Spearman 为 $0.838$、前 20 名交集为 $15/20$；从第 239 行改选最坏域最优的第 199 行，可使最坏域改善 $1.66\%$，均值仅恶化 $0.038\%$，且 256 个候选的均值--最坏域 Pareto 前沿只有这两个非支配点. 若改以域间标准差为目标，标准差由 $0.1238$ 降至 $0.0680$（改善 $45.1\%$），代价是均值恶化 $5.77\%$. 超额损失还依赖归一化：原始尺度与主目标都选第 239 行（Spearman $0.986$），按各域可改进空间归一后则改选第 219 行，并使其与最坏域目标的相关降至 $0.642$. 因此后续所有“最优配比”都必须连同目标标尺一起解释. 

\paragraph{代理模型选择本身的风险（RegMix \cite{regmix}）. \quad}
平方根二次面与平方根 RBF 在 A4/A5 重复交叉验证中的宏平均 RMSE 仅差 $0.001304$；两者在历史 1M 候选表上都选第 239 行. 因此，同一候选池不足以证明二次面产生了不同的最优配方，事后对照也不能替代嵌套验证. RegMix 提倡对高排名方案进行稳健筛选，本题附件没有额外训练接口验证这种集成策略；模型族和候选池的不确定性仍应与给定集合内的后悔值分开报告. 

\paragraph{``高质量源未必高份额''（RegMix \cite{regmix}）. \quad}
在六个已映射域上，域质量与该域在选定配方中的份额的 Spearman 相关仅为 $+0.086$：
质量最高的 arxiv（$0.8209$）只占 $9.6\%$，而质量较低的 github（$0.5041$）占 $3.0\%$，
没有质量观测的 pubmed\_abstracts 反而占 $34.9\%$. 这说明配比最优性与单域质量评分是
两个不同的量，与 RegMix 关于``直觉上的高质量数据源未必最优''的实验结论一致，
也再次限制了把 $Q$ 直接乘进配比的写法. 

\paragraph{训练强度分层诊断（次缩放 \cite{acl2025quality}）. \quad}
ACL 2025 报告在约 $20$ tokens/parameter 附近资源边际收益会失衡. B1 的 $D/N$ 中位数为
$97.3$、范围 $0.01$--$4251$，跨越该区间. 按 $D/N<25$、$25$--$100$、$\ge100$ 三档分组后，
经典幂律的残差 RMSE 分别为 $1.68\times10^{-4}$、$1.20\times10^{-4}$ 与 $1.45\times10^{-4}$，
没有随训练强度单调整体恶化的迹象. 因此在 Pythia 的观测范围内，不需要引入额外的
次缩放衰减项；这支持第二问对 $D_{\mathrm{eff}}$ 的单一幂律写法. 

\paragraph{桥接函数形式（FLP \cite{flp}）. \quad}
在 C6 桥接上把线性换成 logistic（$S=100/[1+e^{-k(L-x_0)}]$）后，高可比组的留一 RMSE 由
$0.553$ 升到 $3.011$，中可比组由 $10.071$ 升到 $14.093$. Pythia 的损失区间（$2.09$--$2.60$）
尚未进入分数饱和段，因此线性桥接在该区间是合适的；FLP 建议的 S 型映射只在分数接近
上界时才有优势，本文数据不覆盖该区间. 

\paragraph{时间扩展标度律的可行性（算法进步 \cite{ho2024}）. \quad}
本文尝试在文献汇总点 B5 上拟合带时间效率项的标度律
$\mathcal L=E+AN^{-\alpha}e^{-\eta_N(Y-Y_0)}+BD^{-\beta}e^{-\eta_D(Y-Y_0)}$. 
B5 的 44 个点只覆盖 $2020$、$2022$、$2023$ 三个年份（分别 $9$、$28$、$7$ 点），
B4 的 12 个模型族则完全没有年份字段，因此时间项无法在本附件内被可靠识别：
拟合得到的 $\eta_N$ 为负且标准误大于点估计. 第四问因此改用 C4 的发布时间字段构造
时间轴，而不在第二问的数据上做时间分解. 


% \subsection{质量评分函数与分组验证图}
% 两项表面指标的十一组试算均冻结其余质量规则；图示只检查条件基线与排序的变化. 
% \begin{figure}[htbp]\centering
% \includegraphics[width=0.96\linewidth]{q1_quality_function_sensitivity.pdf}
% \caption{第一问两项表面适宜度函数的敏感性. 点均来自冻结的十一组试算；橙点为联合平台主比较，虚线为正式主口径 $Q_0$. 七域排序相同不代表绝对质量标尺得到外部验证. }\label{fig:q1-function-sensitivity-v11}
% \end{figure}

\subsection{第二问分组验证与备选质量响应}
表~\ref{tab:q2-v26-validation} 是正文主选型证据；下图为相同分组结果的可视化. 
\begin{figure}[htbp]\centering
\includegraphics[width=0.91\linewidth]{q2_v26_validation_rmse_comparison.pdf}
\caption{第二问固定分组口径下的预测误差. B7 的新质量等级不参与模型选择，B8 未并入拟合. }\label{fig:q2-v26-validation}
\end{figure}
积分斜率候选使用另一项完整资源单元五折协议，不能与逐单元留一结果直接拼接. 
\begin{table}[htbp]\centering\small
\caption{B6 完整资源单元五折检验的新增质量响应比较；360 行留出误差合并计算}\label{tab:q2-integrated-cv-v38}
\begin{tabular}{lr}\toprule
质量响应形式 & 原始 Loss 的合并 RMSE\\\midrule
曲率模型，逐折重新拟合 & 0.06669707\\
四段单调积分斜率，逐折重新拟合 & 0.06589885\\\bottomrule
\end{tabular}
\end{table}
\begin{figure}[htbp]\centering
\includegraphics[width=0.86\linewidth]{q2_grouped_resource_validation.pdf}
\caption{同一完整资源单元五折下两种质量响应的留出 RMSE；橙色候选的微小总体优势在成对重抽样下未排除零差. }\label{fig:q2-integrated-cv-v38}
\end{figure}

\subsection{第三问稳健配置的数值折中}
正文表~\ref{tab:q3-v26-regret} 与表~\ref{tab:q3-cvar-evaluation-v38} 给出收益和不利的第二组重抽样结果；以下两图保留名义成本与代表情景的逐项对照. 
\begin{figure}[H]\centering
\includegraphics[width=0.91\linewidth]{q3_v26_nominal_cost_of_robustness.pdf}
\caption{降低有限情景最大后悔值所付出的名义 Loss 代价. }\label{fig:q3-v26-cost}
\end{figure}
\begin{figure}[H]\centering
\includegraphics[width=0.83\linewidth]{q3_cvar_tradeoff_power_1e19.pdf}
\caption{$10^{19}$ FLOPs、幂质量成本下第二组重抽样的平均与尾部后悔值；横纵坐标均以越低越好解释. }\label{fig:q3-cvar-evaluation-v38}
\end{figure}

\subsection{第四问惩罚敏感性与校准图}
惩罚强度的九档重拟合检查时间斜率的零值及规模斜率的变化；校准图与正文表~\ref{tab:q4-v26-calibration} 使用相同回溯记录. 
\begin{table}[htbp]\centering\footnotesize
\setlength{\tabcolsep}{5pt}
\caption{共享斜率对惩罚强度的敏感性；斜率为原分数尺度的换算值，Pinball 为归一化加权目标，所有拟合的系数盒均未激活}\label{tab:q4-lambda-sensitivity}
\begin{tabular}{lrrrrr}\toprule
$\lambda_{\mathrm{QR}}$ & $b_N$（分/十倍） & $b_t$（分/年） & 加权 Pinball & 后训练 24 月 & 预训练 24 月\\\midrule
$0.01$ & $1.852$ & $-1.5\times10^{-14}$ & $0.008916$ & $32.699$ & $22.641$\\
$0.05$ & $1.852$ & $6.7\times10^{-15}$ & $0.008916$ & $32.699$ & $22.641$\\
$0.10$ & $1.852$ & $1.0\times10^{-14}$ & $0.008916$ & $32.699$ & $22.641$\\
$0.25$ & $1.909$ & $0.1048$ & $0.008923$ & $32.972$ & $22.938$\\
$0.50$ & $1.955$ & $-1.8\times10^{-14}$ & $0.008928$ & $32.679$ & $22.656$\\
$1$（主候选） & $1.955$ & $-1.0\times10^{-15}$ & $0.008928$ & $32.679$ & $22.656$\\
$2$ & $2.120$ & $3.7\times10^{-14}$ & $0.009035$ & $32.647$ & $22.680$\\
$5$ & $4.027$ & $-4.6\times10^{-14}$ & $0.012828$ & $32.277$ & $22.956$\\
$10$ & $5.246$ & $3.7\times10^{-14}$ & $0.016112$ & $33.263$ & $23.133$\\\bottomrule
\end{tabular}
\end{table}
\begin{figure}[htbp]\centering
\includegraphics[width=0.93\linewidth]{q4_v26_calibration.pdf}
\caption{区间覆盖与宽度的对照，相应区间分数见表~\ref{tab:q4-v26-calibration}. 高覆盖可能来自几乎完整的分数支持区间，不能单独证明预测精确. }\label{fig:q4-v26-calibration}
\end{figure}



  %\input{content/09_appendix_data}
  \section{程序清单与关键代码}\label{app:v42-programs}

%\subsection{关键代码片段}
下列片段取自上表所列程序本体，行号对应当前版本；完整程序见相应文件。

片段 1：Q1 平方根坐标二次岭与固定候选集（\texttt{q1/mixture\_solver.py} L 110-127） 候选模型集合在任何嵌套或外部评价之前固定；\texttt{sqrt} 坐标给出 $\sqrt{p_i}$ 与 $\sqrt{p_ip_k}$ 项。

\begin{lstlisting}[language=Python,breaklines=true,basicstyle=\small\ttfamily,showstringspaces=false,frame=single,framesep=4pt,rulecolor=\color{gray!50},backgroundcolor=\color{gray!5},title={mixture\_solver.py}]
def candidate_models(jobs: int = 4) -> dict:
    """Finite candidate set fixed before nested or external evaluation."""
    def trees(n=256, leaf=3, features=1.0):
        return ExtraTreesRegressor(n_estimators=n, min_samples_leaf=leaf,
                                   max_features=features, random_state=SEED, n_jobs=jobs)

    def quadratic(kind, alpha):
        """Degree-2 polynomial in raw, Helmert contrast, or sqrt coordinates.

        In raw proportions, the sqrt model contains sqrt(p_i) and
        sqrt(p_i * p_j) terms, as well as p_i terms. Finite-difference effects
        are defined at their probes; they are not polynomial coefficients in p.
        """
        if kind not in {"raw", "contrast", "sqrt"}:
            raise ValueError(f"Unknown quadratic coordinates: {kind}")
        steps = [] if kind == "raw" else [SimplexFeatures(kind)]
        return make_pipeline(*steps, PolynomialFeatures(2, include_bias=False),
                             StandardScaler(), Ridge(alpha=alpha))
\end{lstlisting}

片段 2：Q3 质量区间的解析下界与全局比较骨架（\texttt{q3/core.py} L 155-182）. 区间下界取 $h(a)$ 与 $g(\ln b)$ 的组合，用凸支撑线给出该区间的目标下界；堆结构按最小下界优先分裂。

\begin{lstlisting}[language=Python,breaklines=true,basicstyle=\small\ttfamily,showstringspaces=false,frame=single,framesep=4pt,rulecolor=\color{gray!50},backgroundcolor=\color{gray!5},title={core.py}]
    def node(self,a,b,budget):
        # h is nondecreasing and each g(log Q) is nonincreasing.
        # Splitting them at opposite endpoints gives a global interval lower bound.
        return self.scale(a,budget,h_override=self.cost.h(a),g_override=self.g(b))['lower']

    def solve(self,budget,*,tol=1e-7,max_nodes=100000,refine=True):
        if not math.isfinite(budget) or budget<=0:raise ValueError("Invalid budget")
        count0=self.calls
        points={}
        def at(q):
            q=float(q)
            if q not in points:points[q]=self.scale(q,budget)
            return points[q]['value']
        def lower_bound(a,b):
            if len(self.params)==1 and self.params[0].beta<=1:
                # Each fixed-Q scale optimum increases with Q when beta<=1.
                # Bounds on its profile derivative can certify monotonicity.
                at(a);at(b)
                p=self.params[0]
                na=math.exp(points[a]['z']);nb=math.exp(points[b]['z'])
                hpmin=min(self.cost.hp(a),self.cost.hp(b))
                hpmax=max(self.cost.hp(a),self.cost.hp(b))
                smin=min(p.slope(a),p.slope(b));smax=max(p.slope(a),p.slope(b))
                upper=p.beta*hpmax/(self.k*na+self.cost.h(a))-smin/b
                lower=p.beta*hpmin/(self.k*nb+self.cost.h(b))-smax/a
                if upper < -1e-11:return points[b]['lower']
                if lower > 1e-11:return points[a]['lower']
            return self.node(a,b,budget)
\end{lstlisting}

片段 3：Q4 有序分位规划的顶点约束（\texttt{q4/ordered\_model.py} L 65-82）. 把形状要求化为矩形四个顶点上的线性约束，与 Pinball 损失的上图变量一起构成凸二次规划。

\begin{lstlisting}[language=Python,breaklines=true,basicstyle=\small\ttfamily,showstringspaces=false,frame=single,framesep=4pt,rulecolor=\color{gray!50},backgroundcolor=\color{gray!5},title={ordered\_model.py}]
    vertices=(domain_vertices(names)-center)/scale
    vdesign=[];orders=[]
    for g in range(2):
        for v in vertices:
            for j in range(3):
                row=np.zeros(p);row[:nv]=v;row[nv+g*3+j]=1;vdesign.append(row)
        for j in range(2):
            row=np.zeros(p);row[nv+g*3+j+1]=1;row[nv+g*3+j]=-1;orders.append(row)
    vd=np.array(vdesign);od=np.array(orders)
    # D x >= b: t >= tau*(y-pred); t >= (tau-1)*(y-pred).
    dim=p+nr
    D=np.vstack([np.column_stack([levels[:,None]*design,np.eye(nr)]),
        np.column_stack([(levels-1)[:,None]*design,np.eye(nr)]),
        np.column_stack([vd,np.zeros((len(vd),nr))]),
        np.column_stack([-vd,np.zeros((len(vd),nr))]),
        np.column_stack([od,np.zeros((len(od),nr))])])
    b=np.r_[levels*yy,(levels-1)*yy,np.zeros(len(vd)),-np.ones(len(vd)),np.zeros(len(od))]
    lower=np.r_[np.full(p,-10.),np.zeros(nr)];upper=np.r_[np.full(p,10.),np.ones(nr)]
\end{lstlisting}

片段 4：顶点预测与全局间隙断言（\texttt{pipeline/finalize\_evidence.py} L 304-311）. 对全部矩形顶点逐点检查取值区间与分位间距，并断言归档的目标间隙上界。

\begin{lstlisting}[language=Python,breaklines=true,basicstyle=\small\ttfamily,showstringspaces=false,frame=single,framesep=4pt,rulecolor=\color{gray!50},backgroundcolor=\color{gray!5},title={finalize\_evidence.py}]
    predictions = fit.matrix(vertex_frame)
    assert predictions.min() >= -1e-7 and predictions.max() <= 100 + 1e-7
    assert (np.diff(predictions, axis=1) >= -1e-7).all()
    metrics = pd.read_csv(Q4 / "rolling_metrics.csv")
    role = metrics.loc[
        metrics.role.eq("reused_late") & metrics.input_mode.eq("resource_forecast")
        & metrics.model.eq("joint_size_time") & metrics.weighting.eq("row")
    ].iloc[0]
\end{lstlisting}



% \subsection{复核输出清单}
% 复算脚本的机器可读输出均在 \path{data/reports/rechecks/} 下：
% \path{q1_penalty_ablation.csv} 与 \path{q1_quality_validation.json}（惩罚消融与验收测试）、
% \path{q4_baseline_caliber.csv} 与 \path{q4_baseline_strict_window_records.csv}（基准入选口径与逐条记录）、
% \path{q4_lambda_sensitivity.csv} 与 \path{q4_domain_and_lambda_checks.json}（斜率惩罚敏感性与声明域）、
% \path{q3_exponent_source_sensitivity.csv} 与同名 JSON（规模指数来源情景）、
% \path{q4_aggregation_with_baseline.csv} 与同名 JSON（含基准的凸组合定权）。
% 这些文件与正文各表的对应关系在相应小节注明。
% 本轮条件情景与核验回执另见 \path{data/reports/review_v44/} 和
% \path{data/reports/paper_number_verification_v44.json}。


\end{appendices}

\end{document}
